\documentclass{article}
\usepackage[margin=0.8in]{geometry}
\usepackage{amsmath,amssymb,amsthm,mathtools}
\usepackage{mathrsfs,bm,bbm}
\usepackage{graphicx,booktabs,array,multirow,tabularx,longtable}
\usepackage{enumitem}
\usepackage{xcolor}
\usepackage{algorithm,algpseudocode}
\usepackage{subcaption,tikz}
\usepackage{siunitx,cancel,accents}
\usepackage{microtype}
\usepackage[numbers,sort&compress]{natbib}
\usepackage[colorlinks=true,linkcolor=blue,citecolor=blue,urlcolor=blue]{hyperref}
\usepackage{cleveref}
\setlength{\emergencystretch}{2em}
\newtheorem{assumption}{Assumption}
\newtheorem{definition}{Definition}
\newtheorem{theorem}{Theorem}
\newtheorem{lemma}{Lemma}
\newtheorem{proposition}{Proposition}
\newtheorem{openendedquestion}{Open-ended Question}
\newtheorem{conjecture}{Conjecture}
\newcommand{\coreref}[1]{\ref{#1}}

\begin{document}

\sloppy

\section*{stat\_\allowbreak{}proximal\_\allowbreak{}quotient\_\allowbreak{}efficiency}

\noindent\textit{Specialization:} v1\par

\paragraph{TL;DR.} The observed-law model uses the full local \(L^\infty\) likelihood-ratio ball, the full fixed-defect spectrally separated operator stratum, and every DQM curve that remains in both bridge-solvability loci. Reverse score lifting on this unrestricted joint germ remains unresolved: continuous kernel projections do not themselves parameterize the coupled solvability intersection. Conditional on tangent equality and the theorem's explicit pathwise product condition, the quotient projection is valid. Under the displayed nested construction/calibration/evaluation split, dimension-aware bridge and calibration rates, spectral cutoff, coefficient envelope, and mixed-bias condition, the affine-ensemble estimator is now proved asymptotically linear, variance-ratio consistent, and Wald-valid pointwise. The same conclusions hold uniformly on the declared contiguous path classes when those quotient hypotheses and all primitive bounds hold at every local law and the transported gradients satisfy the stated continuity and uniform-integrability conditions. This inference theorem does not resolve \(\mathrm{oeq:score-lifting-completeness}\). The exact \(36\)-cell efficiency value \(773723312/334026875\) and strict single-pair nonattainment remain valid, and the observed bridge functional equals a causal ATE only under the separate latent identification proposition.

\section{Project justification}

\paragraph{Gap.} Existing nonunique-bridge results do not establish the full tangent closure of the simultaneous outcome- and treatment-solvability germ considered here. The remaining load-bearing gap is reverse score lifting across moving kernels and the arms' shared \((W,X)\) marginal; the repaired estimator theorem is conditional on tangent equality and bridge-score pathwise validity at every law where efficiency is claimed.

\paragraph{Niche.} The contribution is quotient-tangent efficiency for a saturated fixed-defect observed-law model in which both bridge fibers may be non-singleton, together with an explicit nested-cross-fitted affine ensemble that estimates the variance-minimizing quotient score. Its distinctive mathematical burden is the coupled score-lifting theorem, not the finite-dimensional projection algebra.

\paragraph{Fill.} The established contribution consists of the affine normal identity, operator-kernel stability, exact finite tangent and nonattainment witnesses, the corrected outcome-only selected-parameter comparison, dimension-aware learned affine-projection convergence, and pointwise plus locally uniform inference under explicit primitive and local-law conditions. General score lifting remains open and is not inferred from the estimator proof.

\section{Related work}

Bennett et al. provide penalized bridge inference under additional regularity, while Zhang et al. emphasize nonuniqueness; neither result is identified here with the unresolved global affine-quotient tangent bound. Imbens et al. and Ai--Shan concern uniqueness-restricted or finite-dimensional efficiency regimes and therefore do not establish the two-fiber quotient tangent calculation. Severini--Tripathi is prior art for efficiency of quotient-like functionals, so the new burden is the simultaneous proximal incidence geometry rather than Hilbert projection alone. Ai--Chen and the proximal nuisance-rate literature motivate Galerkin and cross-fit calculations; the independent calibration layer, singular-Gram cutoff, four-corner affine conversion, and their dimension-aware perturbation proof are the explicit construction used here. The final Bennett publisher theorem text has not been verified here, and no mathematical claim depends on an asserted difference from inaccessible text.

\section{Setup}

Target (estimand / structural parameter): \(\theta(P)=E_P[h_1(W,X)-h_0(W,X)]\).
Primary functional / estimator / diagnostic: \([\mathcal T_P=\mathcal C_P^\perp\text{ and bridge-score pathwise validity at }P]\Longrightarrow\phi_{\mathrm{eff}}=\phi_0-\Pi_{\mathcal C_P}\phi_0;\quad P=P^\star\Longrightarrow\mathcal T_{P^\star}=\mathcal C_{P^\star}^\perp,\ V_{\mathrm{eff}}(P^\star)=773723312/334026875\)..
\begin{itemize}
  \item \texttt{n} --- \(n\in\mathbb N\); \texttt{asymptotic index}
  \par\smallskip\noindent\emph{Definition.} \texttt{Sample size.}
  \item \texttt{a} --- \(a\in\{0,1\}\); \texttt{index}
  \par\smallskip\noindent\emph{Definition.} \texttt{Treatment-\allowbreak{}arm index.}
  \item \texttt{A} --- \texttt{binary random variable}; \(A\in\{0,1\}\); \texttt{treatment}
  \par\smallskip\noindent\emph{Definition.} \texttt{Observed treatment.}
  \item \texttt{Z} --- \texttt{random element of a standard Borel space}; \texttt{proxy}
  \par\smallskip\noindent\emph{Definition.} \texttt{Observed treatment-\allowbreak{}inducing proxy.}
  \item \texttt{W} --- \texttt{random element of a standard Borel space}; \texttt{proxy}
  \par\smallskip\noindent\emph{Definition.} \texttt{Observed outcome-\allowbreak{}inducing proxy.}
  \item \texttt{X} --- \texttt{random element of a standard Borel space}; \texttt{covariate}
  \par\smallskip\noindent\emph{Definition.} \texttt{Observed baseline covariates.}
  \item \texttt{Y} --- \texttt{real random variable}; \(Y\in\mathbb R\); \texttt{outcome}
  \par\smallskip\noindent\emph{Definition.} \texttt{Observed outcome.}
  \item \texttt{O} --- \texttt{observed random vector}; \texttt{sampling unit}
  \par\smallskip\noindent\emph{Definition.} \(O=(Y,A,Z,W,X)\).
  \item \texttt{nu} --- \texttt{sigma-\allowbreak{}finite measure}; \texttt{dominating measure}
  \par\smallskip\noindent\emph{Definition.} \(\nu\) is a fixed dominating measure on the sample space of \(O\).
  \item \texttt{P} --- \texttt{probability law}; \(P\ll\nu\); \texttt{data-\allowbreak{}generating law}
  \par\smallskip\noindent\emph{Definition.} Observed-data law of \(O\).
  \item \texttt{E} --- \texttt{expectation operator}; \(f\mapsto E_P[f(O)]\); \texttt{integration}
  \par\smallskip\noindent\emph{Definition.} Expectation under \(P\).
  \item \texttt{Ia} --- \texttt{indicator random variable}; \texttt{arm selector}
  \par\smallskip\noindent\emph{Definition.} \(I_a=\mathbf 1\{A=a\}\).
  \item \texttt{ta} --- \texttt{signed arm coefficient}; \texttt{ATE contrast sign}
  \par\smallskip\noindent\emph{Definition.} \(t_a=2a-1\).
  \item \texttt{H} --- \texttt{Hilbert space}; \(\mathcal H=L^2(P_{W,X})\); \texttt{outcome-\allowbreak{}bridge domain}
  \par\smallskip\noindent\emph{Definition.} Square-integrable real functions of \((W,X)\).
  \item \texttt{Ga} --- \texttt{arm-\allowbreak{}specific Hilbert space}; \(\mathcal G_a=L^2(I_a\,dP_{A,Z,X})\); \texttt{conditional-\allowbreak{}moment range}
  \par\smallskip\noindent\emph{Definition.} Square-integrable real functions of \((Z,X)\) with inner product \(\langle u,v\rangle_{\mathcal G_a}=E[I_a u(Z,X)v(Z,X)]\).
  \item \texttt{Ta} --- \texttt{bounded linear operator}; \(T_a:\mathcal H\to\mathcal G_a\); \texttt{outcome-\allowbreak{}bridge operator}
  \par\smallskip\noindent\emph{Definition.} \((T_a h)(Z,X)=E[h(W,X)\mid A=a,Z,X]\).
  \item \texttt{Tastar} --- \texttt{adjoint operator}; \(T_a^*:\mathcal G_a\to\mathcal H\); \texttt{treatment-\allowbreak{}bridge operator}
  \par\smallskip\noindent\emph{Definition.} \((T_a^*q)(W,X)=E[I_a q(Z,X)\mid W,X]\).
  \item \texttt{ga} --- \texttt{arm-\allowbreak{}specific regression}; \(g_a\in\mathcal G_a\); \texttt{outcome-\allowbreak{}bridge right-\allowbreak{}hand side}
  \par\smallskip\noindent\emph{Definition.} \(g_a(Z,X)=E[Y\mid A=a,Z,X]\).
  \item \texttt{Ha} --- \texttt{affine solution set}; \(\mathcal H_a(P)\subseteq\mathcal H\); \texttt{outcome-\allowbreak{}bridge fiber}
  \par\smallskip\noindent\emph{Definition.} \(\mathcal H_a(P)=\{h\in\mathcal H:T_a h=g_a\}\).
  \item \texttt{Qa} --- \texttt{affine solution set}; \(\mathcal Q_a(P)\subseteq\mathcal G_a\); \texttt{treatment-\allowbreak{}bridge fiber}
  \par\smallskip\noindent\emph{Definition.} \(\mathcal Q_a(P)=\{q\in\mathcal G_a:T_a^*q=1\}\).
  \item \texttt{Na} --- \texttt{closed linear subspace}; \(\mathcal N_a(P)\subseteq\mathcal H\); \texttt{outcome-\allowbreak{}bridge null space}
  \par\smallskip\noindent\emph{Definition.} \(\mathcal N_a(P)=\ker(T_a)\).
  \item \texttt{Ma} --- \texttt{closed linear subspace}; \(\mathcal K_a(P)\subseteq\mathcal G_a\); \texttt{treatment-\allowbreak{}bridge null space}
  \par\smallskip\noindent\emph{Definition.} \(\mathcal K_a(P)=\ker(T_a^*)\).
  \item \texttt{kappaa} --- \texttt{rank-\allowbreak{}stratum coordinate}; \texttt{local rank label}
  \par\smallskip\noindent\emph{Definition.} \(\kappa_a(P)=(\dim\mathcal N_a(P),\dim\mathcal K_a(P))\).
  \item \texttt{PiNa} --- \texttt{orthogonal projection}; \texttt{moving-\allowbreak{}kernel coordinate}
  \par\smallskip\noindent\emph{Definition.} \(\Pi_{N,a}(P)\) is the orthogonal projection from \(\mathcal H\) onto \(\mathcal N_a(P)\).
  \item \texttt{PiMa} --- \texttt{orthogonal projection}; \texttt{moving-\allowbreak{}cokernel coordinate}
  \par\smallskip\noindent\emph{Definition.} \(\Pi_{K,a}(P)\) is the orthogonal projection from \(\mathcal G_a\) onto \(\mathcal K_a(P)\).
  \item \texttt{Sa} --- \texttt{bridge signal}; \((h,q)\mapsto S_a(h,q)\); \texttt{arm-\allowbreak{}specific doubly robust signal}
  \par\smallskip\noindent\emph{Definition.} \(S_a(h,q)=h(W,X)+I_aq(Z,X)\{Y-h(W,X)\}\).
  \item \texttt{F} --- \texttt{outcome-\allowbreak{}solvability normal generator}; \((a,h,m)\mapsto F_a(h,m)\); \texttt{normal generator}
  \par\smallskip\noindent\emph{Definition.} \(F_a(h,m)=I_am(Z,X)\{Y-h(W,X)\}\).
  \item \texttt{G} --- \texttt{treatment-\allowbreak{}solvability normal generator}; \((a,q,u)\mapsto G_a(q,u)\); \texttt{normal generator}
  \par\smallskip\noindent\emph{Definition.} \(G_a(q,u)=u(W,X)\{I_aq(Z,X)-1\}\).
  \item \texttt{R} --- \texttt{null-\allowbreak{}space-\allowbreak{}interaction normal generator}; \((a,m,u)\mapsto R_a(m,u)\); \texttt{rank-\allowbreak{}stratum normal generator}
  \par\smallskip\noindent\emph{Definition.} \(R_a(m,u)=I_am(Z,X)u(W,X)\).
  \item \texttt{CP} --- \texttt{closed normal subspace}; \(\mathcal C_P\subseteq L_0^2(P)\); \texttt{candidate tangent normal}
  \par\smallskip\noindent\emph{Definition.} \(\mathcal C_P\) is the closed linear span of all square-integrable \(F_a(h,m)\), \(G_a(q,u)\), and \(R_a(m,u)\) with \(h\in\mathcal H_a(P)\), \(q\in\mathcal Q_a(P)\), \(m\in\mathcal K_a(P)\), and \(u\in\mathcal N_a(P)\), across \(a\in\{0,1\}\).
  \item \texttt{PiC} --- \texttt{orthogonal projection}; \(\Pi_{\mathcal C_P}:L_0^2(P)\to\mathcal C_P\); \texttt{efficiency projection}
  \par\smallskip\noindent\emph{Definition.} Orthogonal projection onto \(\mathcal C_P\).
  \item \texttt{Ustar} --- \texttt{latent Bernoulli variable}; \(U^\star\in\{0,1\}\); \texttt{witness coordinate}
  \par\smallskip\noindent\emph{Definition.} \texttt{Auxiliary latent variable in the finite witness construction.}
  \item \texttt{Pstar} --- \texttt{strictly positive finite observed law}; \texttt{nonattainment witness}
  \par\smallskip\noindent\emph{Definition.} The \(36\)-cell observed law defined in \(\mathrm{def:finite-witness}\).
  \item \texttt{vnull} --- \texttt{three-\allowbreak{}vector}; \(v=(3,-7,3)^\top\in\mathbb R^3\); \texttt{witness null vector}
  \par\smallskip\noindent\emph{Definition.} Common left and right bridge-operator null direction in \(P^\star\).
  \item \texttt{hsel} --- \texttt{arm-\allowbreak{}specific selected outcome bridge}; \(h_a^{\mathrm{sel}}:\{0,1,2\}\to\mathbb R\); \texttt{finite-\allowbreak{}witness comparator bridge}
  \par\smallskip\noindent\emph{Definition.} \(h_a^{\mathrm{sel}}(w)=2w/5+a/5\) for \(w\in\{0,1,2\}\).
  \item \texttt{qsel} --- \texttt{arm-\allowbreak{}specific selected treatment bridge}; \(q_a^{\mathrm{sel}}:\{0,1,2\}\to\mathbb R\); \texttt{finite-\allowbreak{}witness comparator bridge}
  \par\smallskip\noindent\emph{Definition.} In the coordinate order \(z=0,1,2\), \(q_0^{\mathrm{sel}}=(0,8/3,16/3)\) and \(q_1^{\mathrm{sel}}=(16/3,8/3,0)\).
  \item \texttt{alpha} --- \texttt{confidence error probability}; \(\alpha\in(0,1)\); \texttt{confidence level}
  \par\smallskip\noindent\emph{Definition.} \texttt{Nominal two-\allowbreak{}sided Wald error probability.}
  \item \texttt{Kfold} --- \texttt{fixed positive integer}; \(K_{\mathrm{fold}}\ge2\); \texttt{sample-\allowbreak{}splitting dimension}
  \par\smallskip\noindent\emph{Definition.} Number of cross-fitting folds, fixed as \(n\to\infty\).
  \item \texttt{k} --- \texttt{fold index}; \(k\in\{1,\ldots,K_{\mathrm{fold}}\}\); \texttt{index}
  \par\smallskip\noindent\emph{Definition.} \texttt{Generic cross-\allowbreak{}fitting fold.}
  \item \texttt{Ifold} --- \texttt{partition of sample indices}; \texttt{held-\allowbreak{}out folds}
  \par\smallskip\noindent\emph{Definition.} \((\mathcal I_1,\ldots,\mathcal I_{K_{\mathrm{fold}}})\) partitions \(\{1,\ldots,n\}\).
  \item \texttt{Jhn} --- \texttt{positive integer sequence}; \(J_n^h\in\mathbb N\); \texttt{outcome grid dimension}
  \par\smallskip\noindent\emph{Definition.} \texttt{Number of nonbase outcome-\allowbreak{}bridge representatives per arm.}
  \item \texttt{Jqn} --- \texttt{positive integer sequence}; \(J_n^q\in\mathbb N\); \texttt{treatment grid dimension}
  \par\smallskip\noindent\emph{Definition.} \texttt{Number of nonbase treatment-\allowbreak{}bridge representatives per arm.}
  \item \texttt{j} --- \texttt{outcome-\allowbreak{}grid index}; \(j\in\{0,\ldots,J_n^h\}\); \texttt{index}
  \par\smallskip\noindent\emph{Definition.} Outcome-bridge cell index, with \(j=0\) the base representative.
  \item \texttt{l} --- \texttt{treatment-\allowbreak{}grid index}; \(\ell\in\{0,\ldots,J_n^q\}\); \texttt{index}
  \par\smallskip\noindent\emph{Definition.} Treatment-bridge cell index, with \(\ell=0\) the base representative.
  \item \texttt{Jn} --- \texttt{positive integer sequence}; \(J_n\in\mathbb N\); \texttt{sieve dimension}
  \par\smallskip\noindent\emph{Definition.} \(J_n=2(J_n^h+J_n^q+J_n^hJ_n^q)\) is the number of arm-specific pure and four-corner contrast columns.
  \item \texttt{rhon} --- \texttt{positive regularization sequence}; \(\rho_n>0\); \texttt{projection regularization}
  \par\smallskip\noindent\emph{Definition.} \texttt{Ridge parameter in the empirical normal equations.}
  \item \texttt{Rn} --- \texttt{positive integer sequence}; \(R_n=(J_n^h+1)(J_n^q+1)\); \texttt{ensemble size}
  \par\smallskip\noindent\emph{Definition.} \texttt{Number of outcome-\allowbreak{}-\allowbreak{}treatment bridge cells per arm in the rectangular affine ensemble.}
  \item \texttt{Cout} --- \texttt{closed linear subspace}; \texttt{outcome-\allowbreak{}only fixed-\allowbreak{}rank normal space}
  \par\smallskip\noindent\emph{Definition.} \(\mathcal C_{\mathrm{out},P}=\overline{\operatorname{span}}\{F_a(h,m),R_a(m,u)\}\) is the outcome-solvability plus fixed-rank normal space.
  \item \texttt{r} --- \texttt{ensemble-\allowbreak{}component index}; \(r\in\{1,\ldots,R_n\}\); \texttt{index}
  \par\smallskip\noindent\emph{Definition.} \texttt{Generic affine-\allowbreak{}ensemble component.}
  \item \texttt{MCX} --- \texttt{local statistical model}; \texttt{outcome-\allowbreak{}only fixed-\allowbreak{}rank comparator}
  \par\smallskip\noindent\emph{Definition.} \(\mathcal M_{\mathrm{out}}(P^\star)\) is the positive finite outcome-only fixed-rank comparator germ with the explicit gauge-selected parameter in \(\mathrm{def:chen-xie-model}\).
  \item \texttt{TCX} --- \texttt{closed tangent space}; \texttt{outcome-\allowbreak{}only comparator tangent space}
  \par\smallskip\noindent\emph{Definition.} \(\mathcal T_{\mathrm{out},P^\star}\) is the closed span of scores of actual paths in the outcome-only fixed-rank comparator germ.
  \item \texttt{zalpha} --- \texttt{normal critical value}; \texttt{Wald critical value}
  \par\smallskip\noindent\emph{Definition.} \(z_{1-\alpha/2}=\Phi^{-1}(1-\alpha/2)\).
  \item \texttt{p} --- \texttt{density}; \texttt{base density}
  \par\smallskip\noindent\emph{Definition.} \(p=dP/d\nu\).
  \item \texttt{s} --- \texttt{mean-\allowbreak{}zero square-\allowbreak{}integrable function}; \texttt{path score}
  \par\smallskip\noindent\emph{Definition.} \(s\in L_0^2(P)\) is a DQM path score.
  \item \texttt{U} --- \texttt{latent random element}; \texttt{causal-\allowbreak{}identification variable}
  \par\smallskip\noindent\emph{Definition.} \(U\) is an unobserved confounder in a latent extension of the observed law.
  \item \texttt{Ya} --- \texttt{potential outcome}; \texttt{causal outcome}
  \par\smallskip\noindent\emph{Definition.} \(Y(a)\) is the outcome under treatment \(a\).
  \item \texttt{epsilon} --- \texttt{likelihood-\allowbreak{}ratio radius}; \(\varepsilon\in(0,1)\); \texttt{local chart radius}
  \par\smallskip\noindent\emph{Definition.} \(\varepsilon\) is the fixed radius of the canonical saturated likelihood-ratio ball.
  \item \texttt{Pprime} --- \texttt{probability law}; \texttt{nearby law}
  \par\smallskip\noindent\emph{Definition.} \(P'\ll P\) is a law in the canonical saturated domain \(\mathcal V_P\).
  \item \texttt{pprime} --- \texttt{density}; \texttt{nearby density}
  \par\smallskip\noindent\emph{Definition.} \(p'=dP'/d\nu\).
  \item \texttt{Vchart} --- \texttt{canonical saturated likelihood-\allowbreak{}ratio ball}; \texttt{saturated local domain}
  \par\smallskip\noindent\emph{Definition.} \(\mathcal V_P=\{P'\ll P:\|p'/p-1\|_{L^\infty(P)}<\varepsilon,\ E_P[p'/p]=1,\ E_{P'}[Y^2]<\infty\}\).
  \item \texttt{Uchart} --- \texttt{set of probability laws}; \texttt{saturated local operator stratum}
  \par\smallskip\noindent\emph{Definition.} \(\mathcal U_P\) is the full fixed-defect, spectrally separated transported-operator stratum inside \(\mathcal V_P\).
  \item \texttt{Lout} --- \texttt{set of probability laws}; \texttt{outcome-\allowbreak{}solvability locus}
  \par\smallskip\noindent\emph{Definition.} \(\mathcal L_{\mathrm{out}}(P)\) is the outcome-bridge solvability locus in \(\mathcal U_P\).
  \item \texttt{Ltrt} --- \texttt{set of probability laws}; \texttt{treatment-\allowbreak{}solvability locus}
  \par\smallskip\noindent\emph{Definition.} \(\mathcal L_{\mathrm{trt}}(P)\) is the treatment-bridge solvability locus in \(\mathcal U_P\).
  \item \texttt{Mjoint} --- \texttt{statistical model}; \(\mathcal M_{\mathrm{joint}}(P)\); \texttt{local observed statistical model}
  \par\smallskip\noindent\emph{Definition.} \texttt{The intersection of the full fixed-\allowbreak{}defect operator stratum with the outcome-\allowbreak{} and treatment-\allowbreak{}solvability loci,\allowbreak{} admitting every DQM curve that remains in that intersection.}
  \item \texttt{Pt} --- \texttt{regular dominated path}; \(t\mapsto P_t\); \texttt{pathwise perturbation}
  \par\smallskip\noindent\emph{Definition.} Any DQM path through \(P\) whose nearby laws remain in the repaired saturated joint model; no bridge lift is included in the path datum.
  \item \texttt{Yaz} --- \texttt{joint-\allowbreak{}intervention potential outcome}; \texttt{Chen-\allowbreak{}-\allowbreak{}Xie comparison outcome}
  \par\smallskip\noindent\emph{Definition.} \(Y(a,z)\) is the outcome under interventions setting treatment to \(a\) and the negative-control exposure to \(z\).
  \item \texttt{Waz} --- \texttt{joint-\allowbreak{}intervention negative-\allowbreak{}control outcome}; \texttt{Chen-\allowbreak{}-\allowbreak{}Xie comparison proxy outcome}
  \par\smallskip\noindent\emph{Definition.} \(W(a,z)\) is the proxy outcome under interventions setting treatment to \(a\) and the negative-control exposure to \(z\).
  \item \texttt{hdagger} --- \texttt{latent-\allowbreak{}compatible outcome bridge}; \texttt{causal-\allowbreak{}identification bridge}
  \par\smallskip\noindent\emph{Definition.} \(h_a^\dagger\in\mathcal H\) satisfies the latent conditional-mean compatibility restriction in \(\mathrm{ass:latent-bridge-compatibility}\).
  \item \texttt{Ba} --- \texttt{finite arm probability matrix}; \texttt{finite constraint coordinate}
  \par\smallskip\noindent\emph{Definition.} \(B_a[z,w]=P(A=a,Z=z,W=w)\) on a finite chart.
  \item \texttt{ba} --- \texttt{finite arm moment vector}; \texttt{finite outcome moment}
  \par\smallskip\noindent\emph{Definition.} \(b_a[z]=E[I_a\mathbf 1\{Z=z\}Y]\).
  \item \texttt{dW} --- \texttt{finite marginal vector}; \texttt{shared proxy marginal}
  \par\smallskip\noindent\emph{Definition.} \(d[w]=P(W=w)\).
  \item \texttt{Xi} --- \texttt{finite constraint map}; \texttt{finite-\allowbreak{}chart defining map}
  \par\smallskip\noindent\emph{Definition.} \(\Xi\) is the joint Schur-complement constraint map defined in \(\mathrm{def:finite-constraint-map}\).
  \item \texttt{hbar} --- \texttt{population outcome-\allowbreak{}bridge dictionary}; \texttt{population outcome-\allowbreak{}bridge representatives}
  \par\smallskip\noindent\emph{Definition.} \((h^0_{aj}):a=0,1,\ j=0,\ldots,J_n^h\) is a population dictionary with \(h^0_{aj}\in\mathcal H_a(P)\).
  \item \texttt{qbar} --- \texttt{population treatment-\allowbreak{}bridge dictionary}; \texttt{population treatment-\allowbreak{}bridge representatives}
  \par\smallskip\noindent\emph{Definition.} \((q^0_{a\ell}):a=0,1,\ \ell=0,\ldots,J_n^q\) is a population dictionary with \(q^0_{a\ell}\in\mathcal Q_a(P)\).
  \item \texttt{Dpop} --- \texttt{population random-\allowbreak{}vector dictionary}; \texttt{population normal sieve}
  \par\smallskip\noindent\emph{Definition.} \(D^0_{J_n}(O)\in\mathbb R^{J_n}\) concatenates, for each arm, the population pure-outcome, pure-treatment, and four-corner cell-signal contrasts displayed in \(\mathrm{def:population-affine-sieve}\).
  \item \texttt{kappaJ} --- \texttt{positive conditioning sequence}; \texttt{population Gram lower bound}
  \par\smallskip\noindent\emph{Definition.} \(\kappa_{J_n}=\lambda_{\min}^{+}\{E[D^0_{J_n}D^{0\top}_{J_n}]\}\).
  \item \texttt{eh2} --- \texttt{nonnegative stochastic rate}; \texttt{outcome bridge rate}
  \par\smallskip\noindent\emph{Definition.} \(e_{h,2,n}\) bounds foldwise \(L^2(P)\) outcome-bridge error.
  \item \texttt{eh4} --- \texttt{nonnegative stochastic rate}; \texttt{outcome bridge fourth-\allowbreak{}moment rate}
  \par\smallskip\noindent\emph{Definition.} \(e_{h,4,n}\) bounds foldwise \(L^4(P)\) outcome-bridge error.
  \item \texttt{eq2} --- \texttt{nonnegative stochastic rate}; \texttt{treatment bridge rate}
  \par\smallskip\noindent\emph{Definition.} \(e_{q,2,n}\) bounds foldwise \(L^2(P)\) treatment-bridge error.
  \item \texttt{eq4} --- \texttt{nonnegative stochastic rate}; \texttt{treatment bridge fourth-\allowbreak{}moment rate}
  \par\smallskip\noindent\emph{Definition.} \(e_{q,4,n}\) bounds foldwise \(L^4(P)\) treatment-bridge error.
  \item \texttt{rout} --- \texttt{nonnegative stochastic rate}; \texttt{outcome residual rate}
  \par\smallskip\noindent\emph{Definition.} \(r_{h,n}\) bounds foldwise outcome conditional-moment residuals.
  \item \texttt{radj} --- \texttt{nonnegative stochastic rate}; \texttt{adjoint residual rate}
  \par\smallskip\noindent\emph{Definition.} \(r_{q,n}\) bounds foldwise adjoint conditional-moment residuals.
  \item \texttt{Pcirc} --- \texttt{strictly positive finite observed law}; \texttt{singular branch witness}
  \par\smallskip\noindent\emph{Definition.} \(P^\circ\) is the singular equal-arm law in \(\mathrm{def:singular-branch-witness}\).
  \item \texttt{Psquare} --- \texttt{strictly positive finite observed law}; \texttt{different-\allowbreak{}row-\allowbreak{}space witness}
  \par\smallskip\noindent\emph{Definition.} \(P^\square\) is the positive \(4\)-by-\(4\), rank-two different-row-space law in \(\mathrm{def:four-by-four-witness}\).
  \item \texttt{ScoreP} --- \texttt{set of path scores}; \(\dot{\mathcal P}_P\subseteq L_0^2(P)\); \texttt{actual score set}
  \par\smallskip\noindent\emph{Definition.} Scores at zero of all regular dominated paths \(P_t\) in \(\mathcal M_{\mathrm{joint}}\).
  \item \texttt{TP} --- \texttt{closed tangent space}; \(\mathcal T_P\subseteq L_0^2(P)\); \texttt{observed-\allowbreak{}data tangent space}
  \par\smallskip\noindent\emph{Definition.} \(\mathcal T_P=\overline{\operatorname{span}}\,\dot{\mathcal P}_P\).
  \item \texttt{LiftHandle} --- \texttt{proof-\allowbreak{}construction handle}; \texttt{open-\allowbreak{}question handle}
  \par\smallskip\noindent\emph{Definition.} The coupled score-lifting construction defined in \(\mathrm{def:score-lifting-handle}\).
  \item \texttt{theta} --- \texttt{real functional}; \(\theta:\mathcal M_{\mathrm{joint}}\to\mathbb R\); \texttt{observed joint-\allowbreak{}bridge functional}
  \par\smallskip\noindent\emph{Definition.} \(\theta(P)=E[h_1(W,X)-h_0(W,X)]\) for arbitrary \(h_a\in\mathcal H_a(P)\); treatment-bridge existence makes the value invariant across outcome-bridge representatives.
  \item \texttt{phi} --- \texttt{centered bridge score}; \((h_0,h_1,q_0,q_1)\mapsto\phi_{h,q}\); \texttt{bridge-\allowbreak{}dependent gradient candidate}
  \par\smallskip\noindent\emph{Definition.} \(\phi_{h,q}=\sum_{a=0}^1 t_aS_a(h_a,q_a)-\theta(P)\).
  \item \texttt{PhiP} --- \texttt{family of square-\allowbreak{}integrable scores}; \(\Phi_P\subseteq L_0^2(P)\); \texttt{valid bridge-\allowbreak{}score family}
  \par\smallskip\noindent\emph{Definition.} \(\Phi_P=\{\phi_{h,q}:h_a\in\mathcal H_a(P),\ q_a\in\mathcal Q_a(P),\ \phi_{h,q}\in L_0^2(P)\}\).
  \item \texttt{phi0} --- \texttt{reference bridge score}; \(\phi_0\in\Phi_P\); \texttt{projection origin}
  \par\smallskip\noindent\emph{Definition.} \texttt{An arbitrary fixed square-\allowbreak{}integrable valid bridge score.}
  \item \texttt{AP} --- closed affine subset of \(L_0^2(P)\); \(\mathcal A_P=\overline{\operatorname{aff}}(\Phi_P)\); \texttt{quotient-\allowbreak{}gradient affine space}
  \par\smallskip\noindent\emph{Definition.} \texttt{Closed affine hull of all square-\allowbreak{}integrable valid bridge scores.}
  \item \texttt{phieff} --- \texttt{square-\allowbreak{}integrable mean-\allowbreak{}zero function}; \(\phi_{\mathrm{eff}}\in L_0^2(P)\); \texttt{affine-\allowbreak{}hull minimum-\allowbreak{}norm gradient candidate}
  \par\smallskip\noindent\emph{Definition.} When score-lifting completeness holds and \(\Phi_P\neq\varnothing\), \(\phi_{\mathrm{eff}}\) is the minimum-norm element of \(\mathcal A_P\) and the canonical gradient in the actual joint-model tangent space.
  \item \texttt{Veff} --- \texttt{nonnegative real functional}; \texttt{affine quotient variance candidate}
  \par\smallskip\noindent\emph{Definition.} When \(\Phi_P\neq\varnothing\), \(V_{\mathrm{eff}}(P)=E[\phi_{\mathrm{eff}}^2]\).
  \item \texttt{phisel} --- \texttt{square-\allowbreak{}integrable mean-\allowbreak{}zero selected-\allowbreak{}pair score}; \(\phi_{\mathrm{sel}}\in L_0^2(P^\star)\); \texttt{Zhang-\allowbreak{}-\allowbreak{}Bennett selected-\allowbreak{}pair comparator score}
  \par\smallskip\noindent\emph{Definition.} \(\phi_{\mathrm{sel}}=\sum_{a=0}^1t_aS_a(h_a^{\mathrm{sel}},q_a^{\mathrm{sel}})-\theta(P^\star)\).
  \item \texttt{Vsel} --- \texttt{positive real number}; \texttt{selected-\allowbreak{}pair asymptotic variance}
  \par\smallskip\noindent\emph{Definition.} \(V_{\mathrm{sel}}=E_{P^\star}[\phi_{\mathrm{sel}}^2]\).
  \item \texttt{alphajn} --- \texttt{nonnegative approximation sequence}; \texttt{normal-\allowbreak{}sieve approximation error}
  \par\smallskip\noindent\emph{Definition.} \(\alpha_{J_n}=\inf_{b\in\mathbb R^{J_n}}\|\Pi_{\mathcal C_P}\phi_0-b^\top D^0_{J_n}\|_{L^2(P)}\), where \(D^0_{J_n}\) is the population normal dictionary.
  \item \texttt{bJ} --- \texttt{population projection coefficient}; \texttt{population sieve projector}
  \par\smallskip\noindent\emph{Definition.} \(b_{J_n}=\arg\min_b E[(\phi_0-b^\top D^0_{J_n})^2]\) on the nonzero Gram eigenspace.
  \item \texttt{Afold} --- \texttt{fold-\allowbreak{}specific index set}; \texttt{nuisance construction sample}
  \par\smallskip\noindent\emph{Definition.} \(\mathcal A_k\subseteq\{1,\ldots,n\}\setminus\mathcal I_k\) is the construction subsample.
  \item \texttt{brho} --- \texttt{nonnegative deterministic rate}; \texttt{ridge bias rate}
  \par\smallskip\noindent\emph{Definition.} \(b_{\rho,n}=\rho_n\|b_{J_n}\|_2/\kappa_{J_n}\) is the explicit ridge-bias bound.
  \item \texttt{lambda0} --- \texttt{population affine weights}; \texttt{population ensemble weights}
  \par\smallskip\noindent\emph{Definition.} \(\lambda^0_{a,j\ell}\) are the scalar exact affine-conversion weights obtained from the population coefficient \(b_{J_n}\) by the displayed deterministic map.
  \item \texttt{Traink} --- \texttt{construction sigma-\allowbreak{}field}; \texttt{bridge construction information}
  \par\smallskip\noindent\emph{Definition.} \(\mathcal D_{A,k}=\sigma(O_i:i\in\mathcal A_k)\).
  \item \texttt{BJ} --- \texttt{positive deterministic sequence}; \texttt{sieve moment envelope}
  \par\smallskip\noindent\emph{Definition.} \(B_{J_n}\) is a fourth-moment envelope for dictionary coordinates and the reference score.
  \item \texttt{hhat} --- \texttt{cross-\allowbreak{}fitted outcome-\allowbreak{}bridge learner}; \texttt{construction-\allowbreak{}sample outcome bridge}
  \par\smallskip\noindent\emph{Definition.} \(\widehat h_{aj}^{(-k,A)}\) is trained only on \(\mathcal D_{A,k}\) by a penalized-minimax or Galerkin bridge criterion.
  \item \texttt{phiout} --- \texttt{square-\allowbreak{}integrable mean-\allowbreak{}zero function}; \texttt{outcome-\allowbreak{}only selected-\allowbreak{}parameter canonical gradient}
  \par\smallskip\noindent\emph{Definition.} \(\phi_{\mathrm{out}}=\phi_{\mathrm{sel}}-\Pi_{\mathcal C_{\mathrm{out},P^\star}}\phi_{\mathrm{sel}}\) is the canonical gradient of the explicit outcome-only comparator at \(P^\star\).
  \item \texttt{qhat} --- \texttt{cross-\allowbreak{}fitted treatment-\allowbreak{}bridge learner}; \texttt{construction-\allowbreak{}sample treatment bridge}
  \par\smallskip\noindent\emph{Definition.} \(\widehat q_{a\ell}^{(-k,A)}\) is trained only on \(\mathcal D_{A,k}\) by a penalized-minimax or Galerkin bridge criterion.
  \item \texttt{Cfold} --- \texttt{fold-\allowbreak{}specific index set}; \texttt{projection calibration sample}
  \par\smallskip\noindent\emph{Definition.} \(\mathcal C_k\subseteq\{1,\ldots,n\}\setminus\mathcal I_k\) is the calibration subsample, disjoint from \(\mathcal A_k\).
  \item \texttt{Vout} --- \texttt{positive real number}; \texttt{outcome-\allowbreak{}only selected-\allowbreak{}parameter efficiency bound}
  \par\smallskip\noindent\emph{Definition.} \(V_{\mathrm{out}}(P^\star)=E_{P^\star}[\phi_{\mathrm{out}}^2]\).
  \item \texttt{Dn} --- \texttt{random-\allowbreak{}vector dictionary}; \(\widehat D_{J_n}^{(-k,A)}(O)\in\mathbb R^{J_n}\); \texttt{nested-\allowbreak{}split normal-\allowbreak{}space sieve}
  \par\smallskip\noindent\emph{Definition.} The fold-specific dictionary is learned on \(\mathcal A_k\), evaluated on \(\mathcal C_k\) for calibration, and then evaluated on \(\mathcal I_k\) for estimation.
  \item \texttt{JWPprime} --- \texttt{canonical isometry}; \texttt{marginal Hilbert-\allowbreak{}space transport}
  \par\smallskip\noindent\emph{Definition.} \(J_W(P')f=f\sqrt{dP'_{W,X}/dP_{W,X}}\) transports \(L^2(P'_{W,X})\) to the fixed space \(\mathcal H\).
  \item \texttt{JGaPprime} --- \texttt{arm-\allowbreak{}specific canonical isometry}; \texttt{arm Hilbert-\allowbreak{}space transport}
  \par\smallskip\noindent\emph{Definition.} \(J_{G,a}(P')u=u\sqrt{d(I_aP'_{A,Z,X})/d(I_aP_{A,Z,X})}\) transports the arm space at \(P'\) to \(\mathcal G_a\).
  \item \texttt{Calibk} --- \texttt{calibration sigma-\allowbreak{}field}; \texttt{calibration information}
  \par\smallskip\noindent\emph{Definition.} \(\mathcal D_{C,k}=\sigma(O_i:i\in\mathcal C_k)\).
  \item \texttt{Ttildea} --- \texttt{transported bounded operator}; \texttt{fixed-\allowbreak{}coordinate bridge operator}
  \par\smallskip\noindent\emph{Definition.} \(\widetilde T_a(P')=J_{G,a}(P')T_a(P')J_W(P')^{-1}:\mathcal H\to\mathcal G_a\).
  \item \texttt{muReda} --- \texttt{nonnegative operator modulus}; \texttt{reduced minimum modulus}
  \par\smallskip\noindent\emph{Definition.} \(\mu_a(P')=\inf\{\|\widetilde T_a(P')u\|:u\perp\ker\widetilde T_a(P'),\ \|u\|=1\}\).
  \item \texttt{nC} --- \texttt{positive integer}; \texttt{calibration sample size}
  \par\smallskip\noindent\emph{Definition.} \(n_{C,k}=|\mathcal C_k|\).
  \item \texttt{Ehat} --- \texttt{calibration empirical average}; \(f\mapsto\widehat E_{C,k}[f]=n_{C,k}^{-1}\sum_{i\in\mathcal C_k}f(O_i)\); \texttt{calibration criterion}
  \par\smallskip\noindent\emph{Definition.} \texttt{Empirical average over the calibration observations,\allowbreak{} conditional on the construction sigma-\allowbreak{}field.}
  \item \texttt{PiNtilde} --- \texttt{orthogonal projection}; \texttt{transported kernel projection}
  \par\smallskip\noindent\emph{Definition.} \(\widetilde\Pi_{N,a}(P')\) projects \(\mathcal H\) onto \(\ker\widetilde T_a(P')\).
  \item \texttt{PiKtilde} --- \texttt{orthogonal projection}; \texttt{transported cokernel projection}
  \par\smallskip\noindent\emph{Definition.} \(\widetilde\Pi_{K,a}(P')\) projects \(\mathcal G_a\) onto \(\ker\widetilde T_a(P')^*\).
  \item \texttt{hout} --- \texttt{gauge-\allowbreak{}selected outcome bridge}; \texttt{single-\allowbreak{}valued comparator selection}
  \par\smallskip\noindent\emph{Definition.} \(h_a^{\mathrm{out}}(P')\) is the unique nearby outcome bridge with \(h_a^{\mathrm{out}}(P')(0)=a/5\) on the finite comparator germ.
  \item \texttt{taun} --- \texttt{positive deterministic sequence}; \texttt{spectral regularization}
  \par\smallskip\noindent\emph{Definition.} \(\tau_n>0\) is the calibration Gram eigenvalue cutoff.
  \item \texttt{Pihat} --- \texttt{random orthogonal projector}; \texttt{calibration spectral projector}
  \par\smallskip\noindent\emph{Definition.} \(\widehat\Pi_k^C\) projects onto calibration Gram eigenvectors with eigenvalue greater than \(\tau_n\).
  \item \texttt{thetaout} --- \texttt{real functional}; \texttt{single-\allowbreak{}valued comparator parameter}
  \par\smallskip\noindent\emph{Definition.} \(\theta_{\mathrm{out}}(P')=\sum_{a=0}^1t_a d(P')^\top h_a^{\mathrm{out}}(P')\) on the outcome-only comparator germ.
  \item \texttt{bhat} --- \texttt{spectrally regularized projection coefficient}; \(\widehat b^{(-k,C\mid A)}\in\mathbb R^{J_n}\); \texttt{calibrated Galerkin coefficient}
  \par\smallskip\noindent\emph{Definition.} The coefficient is zero below the calibration cutoff \(\tau_n\) and equals \((\widehat\gamma+\rho_n)^{-1}\) times the cross-moment coordinate on each retained calibration eigenvector.
  \item \texttt{Mqn} --- \texttt{positive deterministic sequence}; \texttt{treatment multiplier envelope}
  \par\smallskip\noindent\emph{Definition.} \(M_{q,n}\) bounds the fourth norm of \(1-I_aq^0_{a\ell}\) uniformly over retained cells.
  \item \texttt{lambdahat} --- \texttt{nested-\allowbreak{}cross-\allowbreak{}fitted affine-\allowbreak{}weight learner}; \texttt{calibrated affine-\allowbreak{}ensemble weight}
  \par\smallskip\noindent\emph{Definition.} \(\widehat\lambda_{a,j\ell}^{(-k,C\mid A)}\) is the deterministic coefficient-to-weight map applied to \(\widehat b^{(-k,C\mid A)}\).
  \item \texttt{Myn} --- \texttt{positive deterministic sequence}; \texttt{outcome multiplier envelope}
  \par\smallskip\noindent\emph{Definition.} \(M_{y,n}\) bounds the fourth norm of \(I_a(Y-h^0_{aj})\) uniformly over retained cells.
  \item \texttt{BDn} --- \texttt{positive deterministic sequence}; \texttt{calibration dictionary envelope}
  \par\smallskip\noindent\emph{Definition.} \(B_{D,n}\) bounds both population \(L^2\) dictionary coordinates and their conditional calibration tails.
  \item \texttt{etaGram} --- \texttt{nonnegative stochastic rate}; \texttt{calibration Gram concentration}
  \par\smallskip\noindent\emph{Definition.} \(\eta_{G,n}\) is the conditional calibration Gram concentration rate defined in \(\mathrm{def:inference-gauges}\).
  \item \texttt{Ln} --- \texttt{positive deterministic sequence}; \texttt{coefficient and affine-\allowbreak{}weight envelope}
  \par\smallskip\noindent\emph{Definition.} \(L_n\) bounds population and learned coefficient norms and the total variation of population and learned affine weights.
  \item \texttt{psihat} --- \texttt{nested-\allowbreak{}cross-\allowbreak{}fitted affine-\allowbreak{}ensemble signal}; \texttt{estimated affine-\allowbreak{}quotient signal}
  \par\smallskip\noindent\emph{Definition.} \(\widehat\psi_i^{(-k,C\mid A)}\) uses construction-sample bridges, calibration-sample cutoff coefficients, and evaluation only on \(i\in\mathcal I_k\).
  \item \texttt{thetahat} --- \texttt{real estimator}; \texttt{affine-\allowbreak{}ensemble one-\allowbreak{}step estimator}
  \par\smallskip\noindent\emph{Definition.} \(\widehat\theta_n=n^{-1}\sum_{i=1}^n\widehat\psi_i\).
  \item \texttt{Vhat} --- \texttt{nonnegative estimator}; \texttt{cross-\allowbreak{}fitted variance estimator}
  \par\smallskip\noindent\emph{Definition.} \(\widehat V_n=n^{-1}\sum_{i=1}^n(\widehat\psi_i-\widehat\theta_n)^2\).
  \item \texttt{CI} --- \texttt{random interval}; \texttt{Wald confidence interval}
  \par\smallskip\noindent\emph{Definition.} \(\mathrm{CI}_{1-\alpha}=\left[\widehat\theta_n\pm z_{1-\alpha/2}\sqrt{\widehat V_n/n}\right]\).
  \item \texttt{B0n} --- \texttt{positive deterministic sequence}; \texttt{calibration base-\allowbreak{}signal envelope}
  \par\smallskip\noindent\emph{Definition.} \(B_{0,n}\) bounds both the population \(L^2\) base signal and its conditional calibration tail.
  \item \texttt{etaCross} --- \texttt{nonnegative stochastic rate}; \texttt{calibration cross-\allowbreak{}moment concentration}
  \par\smallskip\noindent\emph{Definition.} \(\eta_{c,n}\) is the conditional calibration score--dictionary concentration rate defined in \(\mathrm{def:inference-gauges}\).
  \item \texttt{dcoord} --- \texttt{nonnegative deterministic array}; \texttt{coordinate dictionary error}
  \par\smallskip\noindent\emph{Definition.} \(d_{c,n}\), \(c=1,\ldots,J_n\), bounds the \(L^2(P)\) error of learned dictionary coordinate \(c\).
  \item \texttt{deltaD} --- \texttt{nonnegative stochastic rate}; \texttt{dimension-\allowbreak{}aware dictionary learning error}
  \par\smallskip\noindent\emph{Definition.} \(\delta_{D,n}=(\sum_{c=1}^{J_n}d_{c,n}^2)^{1/2}\) bounds the learned dictionary error in \(L^2(P;\ell^2)\).
  \item \texttt{d0n} --- \texttt{nonnegative deterministic sequence}; \texttt{base signal error}
  \par\smallskip\noindent\emph{Definition.} \(d_{0,n}\) bounds the \(L^2(P)\) error of the learned base contrast.
  \item \texttt{LambdaJ} --- \texttt{nonnegative deterministic sequence}; \texttt{population Gram upper eigenvalue}
  \par\smallskip\noindent\emph{Definition.} \(\Lambda_{J_n}=\lambda_{\max}E[D^0_{J_n}D_{J_n}^{0\top}]\).
  \item \texttt{zetaG} --- \texttt{nonnegative stochastic rate}; \texttt{total Gram perturbation}
  \par\smallskip\noindent\emph{Definition.} \(\zeta_{G,n}\) is the total learned calibration Gram perturbation defined in \(\mathrm{def:inference-gauges}\).
  \item \texttt{zetaC} --- \texttt{nonnegative stochastic rate}; \texttt{total cross-\allowbreak{}moment perturbation}
  \par\smallskip\noindent\emph{Definition.} \(\zeta_{c,n}\) is the total learned calibration cross-moment perturbation defined in \(\mathrm{def:inference-gauges}\).
  \item \texttt{etaProj} --- \texttt{nonnegative stochastic rate}; \texttt{spectral projector rate}
  \par\smallskip\noindent\emph{Definition.} \(\eta_{\Pi,n}=\zeta_{G,n}/\kappa_{J_n}\) bounds the calibration spectral-projector error.
  \item \texttt{Gamman} --- \texttt{nonnegative stochastic rate}; \texttt{learned signal convergence rate}
  \par\smallskip\noindent\emph{Definition.} \(\Gamma_n\) is the complete learned efficient-signal error gauge defined in \(\mathrm{def:inference-gauges}\).
\end{itemize}

\section{Assumptions}

\begin{assumption}[ass:dominated-law]\label{ass:dominated-law}
\(P\ll\nu\).
\par\smallskip\noindent\emph{Standard} (Dominated semiparametric model, \cite{vanDerVaart1998Asymptotic}).
\end{assumption}

\begin{assumption}[ass:iid-sampling]\label{ass:iid-sampling}
\(O_1,\ldots,O_n\stackrel{\mathrm{i.i.d.}}{\sim}P\).
\par\smallskip\noindent\emph{Standard} (Independent and identically distributed sampling, \cite{vanDerVaart1998Asymptotic}).
\end{assumption}

\begin{assumption}[ass:outcome-second-moment]\label{ass:outcome-second-moment}
\(E[Y^2]<\infty\).
\par\smallskip\noindent\emph{Standard} (Finite outcome second moment, \cite{BennettEtAl2026StronglyIdentified}).
\end{assumption}

\begin{assumption}[ass:local-uniform-equivalence]\label{ass:local-uniform-equivalence}
\(0<\varepsilon<1\).
\par\smallskip\noindent\emph{Novel.} Fixing a radius below one makes the full likelihood-ratio ball positive and uniformly equivalent to the base law; saturation itself is imposed by the chart definition.
\end{assumption}

\begin{assumption}[ass:transported-operator-continuity]\label{ass:transported-operator-continuity}
\(\|\widetilde T_a(P')-\widetilde T_a(P)\|_{\mathrm{op}}\to0\) as \(P'\to P\) in Hellinger distance for \(a=0,1\).
\par\smallskip\noindent\emph{Novel.} This fixed-coordinate continuity holds for positive finite-proxy charts and for conditional finite-matrix models with continuously varying cell laws; it permits moving marginal Hilbert spaces without selecting a bridge.
\end{assumption}

\begin{assumption}[ass:uniform-reduced-modulus]\label{ass:uniform-reduced-modulus}
For each \(a\in\{0,1\}\), \(\mu_a(P)>0\).
\par\smallskip\noindent\emph{Novel.} A positive base reduced minimum modulus is exactly the closed-range condition; the full stratum definition imposes the explicit nearby separation threshold.
\end{assumption}

\begin{assumption}[ass:outcome-bridge-existence]\label{ass:outcome-bridge-existence}
\(\mathcal H_a(P')\neq\varnothing\) for \(a=0,1\).
\par\smallskip\noindent\emph{Standard} (Square-integrable proximal outcome-bridge existence, \cite{ZhangLiMiaoTchetgen2023Nonuniqueness}).
\end{assumption}

\begin{assumption}[ass:treatment-bridge-existence]\label{ass:treatment-bridge-existence}
\(\mathcal Q_a(P')\neq\varnothing\) for \(a=0,1\).
\par\smallskip\noindent\emph{Standard} (Square-integrable proximal treatment-bridge existence, \cite{ZhangLiMiaoTchetgen2023Nonuniqueness}).
\end{assumption}

\begin{assumption}[ass:causal-consistency]\label{ass:causal-consistency}
\(Y=Y(A)\) almost surely.
\par\smallskip\noindent\emph{Standard} (Consistency, \cite{MiaoGengTchetgen2018Proxy}).
\end{assumption}

\begin{assumption}[ass:latent-exchangeability]\label{ass:latent-exchangeability}
\(Y(a)\perp(A,Z)\mid(U,X)\) for \(a=0,1\).
\par\smallskip\noindent\emph{Standard} (Latent exchangeability with outcome negative-control exposure, \cite{MiaoGengTchetgen2018Proxy}).
\end{assumption}

\begin{assumption}[ass:latent-positivity]\label{ass:latent-positivity}
\(0<P(A=a\mid U,X)<1\) almost surely for \(a=0,1\).
\par\smallskip\noindent\emph{Standard} (Latent treatment positivity, \cite{MiaoGengTchetgen2018Proxy}).
\end{assumption}

\begin{assumption}[ass:w-treatment-negative-control]\label{ass:w-treatment-negative-control}
\(W\perp(A,Z)\mid(U,X)\).
\par\smallskip\noindent\emph{Standard} (Treatment negative-control outcome condition, \cite{MiaoGengTchetgen2018Proxy}).
\end{assumption}

\begin{assumption}[ass:latent-bridge-compatibility]\label{ass:latent-bridge-compatibility}
\(E[h_a^\dagger(W,X)\mid U,X]=E[Y(a)\mid U,X]\) for \(a=0,1\).
\par\smallskip\noindent\emph{Standard} (Latent-compatible proximal outcome bridge, \cite{MiaoGengTchetgen2018Proxy}).
\end{assumption}

\begin{assumption}[ass:population-dictionary-approximation]\label{ass:population-dictionary-approximation}
\(\alpha_{J_n}\to0\) as \(J_n\to\infty\).
\par\smallskip\noindent\emph{Standard} (Dense Galerkin approximation, \cite{AiChen2012SequentialMoments}).
\end{assumption}

\begin{assumption}[ass:outcome-learning-rates]\label{ass:outcome-learning-rates}
Uniformly over folds, arms, and retained outcome-grid indices, \(\|\widehat h_{aj}^{(-k,A)}-h^0_{aj}\|_2=O_p(e_{h,2,n})\) and \(\|\widehat h_{aj}^{(-k,A)}-h^0_{aj}\|_4=O_p(e_{h,4,n})\).
\par\smallskip\noindent\emph{Standard} (Cross-fitted bridge estimation rates, \cite{BennettEtAl2023SourceConditionDR}).
\end{assumption}

\begin{assumption}[ass:treatment-learning-rates]\label{ass:treatment-learning-rates}
Uniformly over folds, arms, and retained treatment-grid indices, \(\|\widehat q_{a\ell}^{(-k,A)}-q^0_{a\ell}\|_2=O_p(e_{q,2,n})\) and \(\|\widehat q_{a\ell}^{(-k,A)}-q^0_{a\ell}\|_4=O_p(e_{q,4,n})\).
\par\smallskip\noindent\emph{Standard} (Cross-fitted dual bridge estimation rates, \cite{BennettEtAl2023SourceConditionDR}).
\end{assumption}

\begin{assumption}[ass:outcome-residual-rate]\label{ass:outcome-residual-rate}
Uniformly over folds, arms, and retained outcome-grid indices, \(\|T_a(\widehat h_{aj}^{(-k,A)}-h^0_{aj})\|_2=O_p(r_{h,n})\).
\par\smallskip\noindent\emph{Standard} (Conditional-moment residual rate, \cite{BennettEtAl2023SourceConditionDR}).
\end{assumption}

\begin{assumption}[ass:adjoint-residual-rate]\label{ass:adjoint-residual-rate}
Uniformly over folds, arms, and retained treatment-grid indices, \(\|T_a^*(\widehat q_{a\ell}^{(-k,A)}-q^0_{a\ell})\|_2=O_p(r_{q,n})\).
\par\smallskip\noindent\emph{Standard} (Adjoint conditional-moment residual rate, \cite{BennettEtAl2023SourceConditionDR}).
\end{assumption}

\begin{assumption}[ass:gram-conditioning]\label{ass:gram-conditioning}
The population Gram satisfies \(\kappa_{J_n}>0\); the cutoff and ridge obey \(0<\tau_n<\kappa_{J_n}/2\), \(\zeta_{G,n}=o_p(\tau_n)\), and \(\rho_n=o(\kappa_{J_n})\).
\par\smallskip\noindent\emph{Standard} (Restricted Galerkin eigenvalue and cutoff condition, \cite{AiChen2012SequentialMoments}).
\end{assumption}

\begin{assumption}[ass:empirical-moment-perturbation]\label{ass:empirical-moment-perturbation}
Conditional on \(\mathcal D_{A,k}\), the calibration observations are independent and, for every unit vector \(u\in\mathbb R^{J_n}\), the centered variables \(X_{i,u}=(u^\top\widehat D_{J_n}^{(-k,A)}(O_i))^2-E[(u^\top\widehat D_{J_n}^{(-k,A)})^2\mid\mathcal D_{A,k}]\) and \(Z_{i,u}=u^\top\widehat D_{J_n}^{(-k,A)}(O_i)\widehat B^{(-k,A)}(O_i)-E[u^\top\widehat D_{J_n}^{(-k,A)}\widehat B^{(-k,A)}\mid\mathcal D_{A,k}]\) satisfy \(E[\exp(\lambda X_{i,u})\mid\mathcal D_{A,k}]\le\exp(C\lambda^2B_{D,n}^4)\) for \(|\lambda|\le c/B_{D,n}^2\) and \(E[\exp(\lambda Z_{i,u})\mid\mathcal D_{A,k}]\le\exp(C\lambda^2B_{D,n}^2B_{0,n}^2)\) for \(|\lambda|\le c/(B_{D,n}B_{0,n})\), with fixed positive constants \(c,C\).
\par\smallskip\noindent\emph{Novel.} The independent calibration layer and its conditional exponential-moment convention are explicit estimator-side conditions needed for a learned singular Gram.
\end{assumption}

\begin{assumption}[ass:moment-envelope-growth]\label{ass:moment-envelope-growth}
\[ \sup_{a,j,\ell}\|1-I_aq^0_{a\ell}\|_4\le M_{q,n},\quad \sup_{a,j,\ell}\|I_a(Y-h^0_{aj})\|_4\le M_{y,n},\quad \max_{c\le J_n}\|D^0_{J_n,c}\|_2\le B_{D,n},\quad \|B^0\|_2\le B_{0,n}; \] conditional exponential-moment bounds are imposed by \(\mathrm{ass:empirical-moment-perturbation}\).
\par\smallskip\noindent\emph{Novel.} These explicit multiplier and population calibration envelopes are the primitive bounds used by the nested projection argument; they are not attributed to the mixed-bias references.
\end{assumption}

\begin{assumption}[ass:ridge-bias-rate]\label{ass:ridge-bias-rate}
\(b_{\rho,n}=\rho_n\|b_{J_n}\|_2/\kappa_{J_n}=o(1)\).
\par\smallskip\noindent\emph{Standard} (Vanishing Galerkin ridge bias, \cite{AiChen2012SequentialMoments}).
\end{assumption}

\begin{assumption}[ass:affine-weight-envelope]\label{ass:affine-weight-envelope}
Uniformly over folds, \(\|b_{J_n}\|_2\vee\|\widehat b^{(-k,C\mid A)}\|_2\le L_n\) with probability tending to one, each arm's population and learned affine weights sum to one, and the sum of the absolute population and learned weights over both arms and all retained cells is at most \(L_n\).
\par\smallskip\noindent\emph{Novel.} Coefficient and affine total-variation constraints are enforceable in the finite calibration program and expose every resolvent multiplier.
\end{assumption}

\begin{assumption}[ass:projection-rate]\label{ass:projection-rate}
The corrected projection gauge in \(\mathrm{def:inference-gauges}\) satisfies \(\Gamma_n=o_p(1)\).
\par\smallskip\noindent\emph{Novel.} This is the exact balance produced by the nested calibration and cutoff perturbation proof.
\end{assumption}

\begin{assumption}[ass:primitive-mixed-bias-rate]\label{ass:primitive-mixed-bias-rate}
\(\sqrt nL_n\min\{r_{h,n}e_{q,2,n},r_{q,n}e_{h,2,n}\}=o(1)\).
\par\smallskip\noindent\emph{Standard} (Cross-fitted primal-dual product-rate condition, \cite{BennettEtAl2023SourceConditionDR}).
\end{assumption}

\begin{assumption}[ass:training-measurability]\label{ass:training-measurability}
For every fold \(k\), all fitted bridges and learned dictionary functions are \(\mathcal D_{A,k}\)-measurable, the calibration Gram, cross-moment, spectral projector, coefficient, and affine weights are \(\mathcal D_{A,k}\vee\mathcal D_{C,k}\)-measurable, and neither sigma-field contains observations indexed by \(\mathcal I_k\).
\par\smallskip\noindent\emph{Standard} (Nested cross-fitting measurability, \cite{BennettEtAl2023SourceConditionDR}).
\end{assumption}

\begin{assumption}[ass:heldout-signal-integrability]\label{ass:heldout-signal-integrability}
Uniformly over folds and retained cells, the evaluation-fold cell signals and adjusted signal in \(\mathrm{def:affine-ensemble-estimator}\) are square-integrable conditionally on \(\mathcal D_{A,k}\vee\mathcal D_{C,k}\) with probability tending to one.
\par\smallskip\noindent\emph{Standard} (Conditional square-integrability for cross-fitted scores, \cite{BennettEtAl2023SourceConditionDR}).
\end{assumption}

\section{Definitions}

\begin{definition}[def:local-operator-chart]\label{def:local-operator-chart}
\par\smallskip\noindent\emph{Defined object.} \texttt{canonical saturated fixed-\allowbreak{}defect closed-\allowbreak{}range chart}
\par\smallskip\noindent\emph{Construction.} Fix \(\varepsilon\in(0,1)\) and define the saturated likelihood-ratio neighborhood
\[
\mathcal V_P=\left\{P'\ll P:\left\|\frac{p'}p-1\right\|_{L^\infty(P)}<\varepsilon,
\ E_P\!\left[\frac{p'}p\right]=1,\ E_{P'}[Y^2]<\infty\right\}.
\]
The full transported fixed-defect, separated stratum in this neighborhood is
\[
\mathcal U_P=\left\{P'\in\mathcal V_P:
\begin{array}{l}
\dim\ker\widetilde T_a(P')=\dim\ker\widetilde T_a(P),\\
\dim\ker\widetilde T_a(P')^*=\dim\ker\widetilde T_a(P)^*,\\
\mu_a(P')>\mu_a(P)/2
\end{array}
\text{ for }a=0,1\right\}.
\]
Thus \(\mathcal U_P\) contains every law in the full likelihood-ratio ball having the displayed defect indices and spectral separation; it is not an analyst-selected subcollection.  The map \(P'\mapsto\widetilde T_a(P')\) is required to be Hellinger-to-operator-norm continuous by \(\mathrm{ass:transported-operator-continuity}\).
\par\smallskip\noindent\emph{Member properties.} \texttt{ass:\allowbreak{}local-\allowbreak{}uniform-\allowbreak{}equivalence}; \texttt{ass:\allowbreak{}transported-\allowbreak{}operator-\allowbreak{}continuity}; \texttt{ass:\allowbreak{}uniform-\allowbreak{}reduced-\allowbreak{}modulus}.
\end{definition}

\begin{definition}[def:outcome-solvability-locus]\label{def:outcome-solvability-locus}
\par\smallskip\noindent\emph{Defined object.} \texttt{outcome-\allowbreak{}bridge solvability locus}
\par\smallskip\noindent\emph{Construction.} \(\mathcal L_{\mathrm{out}}(P)=\{P'\in\mathcal U_P:\mathcal H_a(P')\neq\varnothing\text{ for }a=0,1\}\).
\par\smallskip\noindent\emph{Member properties.} \texttt{ass:\allowbreak{}outcome-\allowbreak{}bridge-\allowbreak{}existence}.
\end{definition}

\begin{definition}[def:treatment-solvability-locus]\label{def:treatment-solvability-locus}
\par\smallskip\noindent\emph{Defined object.} \texttt{treatment-\allowbreak{}bridge solvability locus}
\par\smallskip\noindent\emph{Construction.} \(\mathcal L_{\mathrm{trt}}(P)=\{P'\in\mathcal U_P:\mathcal Q_a(P')\neq\varnothing\text{ for }a=0,1\}\).
\par\smallskip\noindent\emph{Member properties.} \texttt{ass:\allowbreak{}treatment-\allowbreak{}bridge-\allowbreak{}existence}.
\end{definition}

\begin{definition}[def:joint-model]\label{def:joint-model}
\par\smallskip\noindent\emph{Defined object.} \texttt{saturated local joint nonunique proximal bridge chart}
\par\smallskip\noindent\emph{Construction.} \[
\mathcal M_{\mathrm{joint}}(P)=\mathcal U_P\cap\mathcal L_{\mathrm{out}}(P)\cap\mathcal L_{\mathrm{trt}}(P).
\]
Every DQM curve \(t\mapsto P_t\) with \(P_0=P\) and \(P_t\in\mathcal M_{\mathrm{joint}}(P)\) for all sufficiently small \(|t|\) is admissible.  Its score \(s\in L_0^2(P)\) is determined by
\[
\sqrt{p_t}=\sqrt p\{1+ts/2\}+o(t)\quad\text{in }L^2(\nu).
\]
No bridge representative, bridge derivative, tangent equality, or score-lifting property is part of this definition.
\par\smallskip\noindent\emph{Member properties.} \texttt{ass:\allowbreak{}dominated-\allowbreak{}law}; \texttt{ass:\allowbreak{}outcome-\allowbreak{}second-\allowbreak{}moment}; \texttt{ass:\allowbreak{}outcome-\allowbreak{}bridge-\allowbreak{}existence}; \texttt{ass:\allowbreak{}treatment-\allowbreak{}bridge-\allowbreak{}existence}.
\end{definition}

\begin{definition}[def:score-family]\label{def:score-family}
\par\smallskip\noindent\emph{Defined object.} \texttt{valid bridge doubly robust score family}
\par\smallskip\noindent\emph{Construction.} \(\Phi_P=\{\phi_{h,q}=\sum_{a=0}^1t_aS_a(h_a,q_a)-\theta(P):h_a\in\mathcal H_a(P),\ q_a\in\mathcal Q_a(P),\ \phi_{h,q}\in L_0^2(P)\}\).
\par\smallskip\noindent\emph{Inputs.} \texttt{def:\allowbreak{}joint-\allowbreak{}model}.
\end{definition}

\begin{definition}[def:normal-space]\label{def:normal-space}
\par\smallskip\noindent\emph{Defined object.} \texttt{three-\allowbreak{}class observable normal space}
\par\smallskip\noindent\emph{Construction.} \(\mathcal C_P=\overline{\operatorname{span}}\{F_a(h,m),G_a(q,u),R_a(m,u):a\in\{0,1\},\ h\in\mathcal H_a(P),\ q\in\mathcal Q_a(P),\ m\in\mathcal K_a(P),\ u\in\mathcal N_a(P),\ F_a,G_a,R_a\in L_0^2(P)\}\).
\par\smallskip\noindent\emph{Inputs.} \texttt{def:\allowbreak{}joint-\allowbreak{}model}.
\end{definition}

\begin{definition}[def:affine-score-hull]\label{def:affine-score-hull}
\par\smallskip\noindent\emph{Defined object.} \texttt{proximal bridge-\allowbreak{}score affine hull}
\par\smallskip\noindent\emph{Construction.} \(\mathcal A_P=\overline{\operatorname{aff}}(\Phi_P)\subseteq L_0^2(P)\).
\par\smallskip\noindent\emph{Inputs.} \texttt{def:\allowbreak{}score-\allowbreak{}family}.
\end{definition}

\begin{definition}[def:quotient-gradient]\label{def:quotient-gradient}
\par\smallskip\noindent\emph{Defined object.} \texttt{proximal bridge-\allowbreak{}quotient gradient candidate}
\par\smallskip\noindent\emph{Construction.} When \(\Phi_P\neq\varnothing\), \(\phi_{\mathrm{eff}}=\operatorname*{argmin}_{\varphi\in\mathcal A_P}E[\varphi^2]\) and \(V_{\mathrm{eff}}(P)=E[\phi_{\mathrm{eff}}^2]\).
\par\smallskip\noindent\emph{Inputs.} \texttt{def:\allowbreak{}affine-\allowbreak{}score-\allowbreak{}hull}.
\end{definition}

\begin{definition}[def:score-lifting-handle]\label{def:score-lifting-handle}
\par\smallskip\noindent\emph{Defined object.} \texttt{coupled score-\allowbreak{}lifting handle}
\par\smallskip\noindent\emph{Construction.} Linearize the joint outcome-solvability, treatment-solvability, and moving-rank maps; quotient their derivative by \(\mathcal N_a(P)\) and \(\mathcal K_a(P)\); orthogonalize the two arms' derivative equations against the shared \((W,X)\) marginal; and use closed-range right inverses plus finite-support chart localization to synthesize path-score approximants for a target element of \(\mathcal C_P^\perp\).
\par\smallskip\noindent\emph{Inputs.} \texttt{def:\allowbreak{}joint-\allowbreak{}model}; \texttt{def:\allowbreak{}normal-\allowbreak{}space}.
\end{definition}

\begin{definition}[def:finite-witness]\label{def:finite-witness}
\par\smallskip\noindent\emph{Defined object.} \texttt{strictly positive nonattainment witness}
\par\smallskip\noindent\emph{Construction.} Let \(X\) be absent, \(U^\star\sim\operatorname{Bernoulli}(1/2)\), \(P^\star(A=1\mid U^\star=u)=(1+2u)/4\), and let \(Z,W\in\{0,1,2\}\) be conditionally independent given \(U^\star\), each with probability vector \((3/5,3/10,1/10)\) at \(U^\star=0\) and \((1/10,3/10,3/5)\) at \(U^\star=1\). Given \((A,U^\star)\), let \(Y\sim\operatorname{Bernoulli}(1/5+A/5+2U^\star/5)\), conditionally independent of \((Z,W)\). The observed marginal law is \(P^\star\).
\end{definition}

\begin{definition}[def:finite-constraint-map]\label{def:finite-constraint-map}
\par\smallskip\noindent\emph{Defined object.} \texttt{finite positive rank-\allowbreak{}two joint constraint map}
\par\smallskip\noindent\emph{Construction.} For \(Z,W\in\{0,1,2\}\), choose a nonsingular rank-two minor \(C_a\) and partition \(B_a=\left[\begin{smallmatrix}C_a&D_a\\E_a&F_a^{B}\end{smallmatrix}\right]\), \(b_a=(b_{a,\mathrm{top}},b_{a,\mathrm{bot}})\), and \(d=(d_{\mathrm{top}},d_{\mathrm{bot}})\). Define \(\Xi(P)=\{F_a^{B}-E_aC_a^{-1}D_a,\ b_{a,\mathrm{bot}}-E_aC_a^{-1}b_{a,\mathrm{top}},\ d_{\mathrm{bot}}-D_a^\top C_a^{-\top}d_{\mathrm{top}}:a=0,1\}\), with the same marginal vector \(d\) in both arms.
\par\smallskip\noindent\emph{Inputs.} \texttt{def:\allowbreak{}finite-\allowbreak{}witness}.
\end{definition}

\begin{definition}[def:chen-xie-model]\label{def:chen-xie-model}
\par\smallskip\noindent\emph{Defined object.} \texttt{outcome-\allowbreak{}only fixed-\allowbreak{}rank selected-\allowbreak{}parameter comparator germ}
\par\smallskip\noindent\emph{Construction.} On the positive \(3\)-by-\(3\), rank-two outcome-solvability component through \(P^\star\), let
\[
\mathcal M_{\mathrm{out}}(P^\star)=\mathcal U_{P^\star}\cap\mathcal L_{\mathrm{out}}(P^\star)
\]
with no treatment-solvability restriction.  For \(P'\) in a sufficiently small component, define \(h^{\mathrm{out}}_a(P')\) to be the unique outcome bridge satisfying the gauge
\[
h^{\mathrm{out}}_a(P')(0)=a/5.
\]
The base null vector has first coordinate \(3\), so the same coordinate is nonzero for the nearby transported kernel; solving the chosen nonsingular rank-two minor makes \(h^{\mathrm{out}}_a(P')\) a single-valued smooth rational function of the cell probabilities. Define
\[
\theta_{\mathrm{out}}(P')=\sum_{a=0}^1t_a\,d(P')^\top h^{\mathrm{out}}_a(P').
\]
At \(P^\star\), \(h^{\mathrm{out}}_a(P^\star)=h_a^{\mathrm{sel}}\).  Write \(\mathcal T_{\mathrm{out},P^\star}\) for this germ's tangent space,
\[
\phi_{\mathrm{out}}=\phi_{\mathrm{sel}}-\Pi_{\mathcal C_{\mathrm{out},P^\star}}\phi_{\mathrm{sel}},
\qquad
V_{\mathrm{out}}(P^\star)=E_{P^\star}[\phi_{\mathrm{out}}^2].
\]
\par\smallskip\noindent\emph{Inputs.} \texttt{def:\allowbreak{}local-\allowbreak{}operator-\allowbreak{}chart}; \texttt{def:\allowbreak{}outcome-\allowbreak{}solvability-\allowbreak{}locus}; \texttt{def:\allowbreak{}finite-\allowbreak{}witness}.
\end{definition}

\begin{definition}[def:population-affine-sieve]\label{def:population-affine-sieve}
\par\smallskip\noindent\emph{Defined object.} \texttt{population bridge and normal dictionaries}
\par\smallskip\noindent\emph{Construction.} Choose \(h^0_{aj}\in\mathcal H_a(P)\), \(j=0,\ldots,J_n^h\), and \(q^0_{a\ell}\in\mathcal Q_a(P)\), \(\ell=0,\ldots,J_n^q\), taking the base cell so that \(\phi_0=\sum_{a=0}^1t_aS_a(h^0_{a0},q^0_{a0})-\theta(P)\). Put \(S^0_{a,j\ell}=S_a(h^0_{aj},q^0_{a\ell})\) and concatenate \(D^{h,0}_{a,j}=t_a(S^0_{a,j0}-S^0_{a,00})\), \(D^{q,0}_{a,\ell}=t_a(S^0_{a,0\ell}-S^0_{a,00})\), and \(D^{hq,0}_{a,j\ell}=t_a(S^0_{a,j\ell}-S^0_{a,j0}-S^0_{a,0\ell}+S^0_{a,00})\) into \(D^0_{J_n}\). Define \(b_{J_n}=\arg\min_bE[(\phi_0-b^\top D^0_{J_n})^2]\) on the nonzero Gram eigenspace, \(\kappa_{J_n}=\lambda_{\min}^{+}E[D^0_{J_n}D_{J_n}^{0\top}]\), and \(\alpha_{J_n}=\inf_b\|\Pi_{\mathcal C_P}\phi_0-b^\top D^0_{J_n}\|_2\).
\par\smallskip\noindent\emph{Inputs.} \texttt{def:\allowbreak{}score-\allowbreak{}family}; \texttt{def:\allowbreak{}normal-\allowbreak{}space}.
\end{definition}

\begin{definition}[def:affine-ensemble-estimator]\label{def:affine-ensemble-estimator}
\par\smallskip\noindent\emph{Defined object.} \texttt{nested cross-\allowbreak{}fitted penalized-\allowbreak{}minimax Galerkin affine ensemble}
\par\smallskip\noindent\emph{Construction.} For each evaluation fold \(\mathcal I_k\), split its complement deterministically into disjoint index sets \(\mathcal A_k\) and \(\mathcal C_k\), with \(\min_k|\mathcal A_k|\wedge|\mathcal C_k|\ge c n\) for a fixed \(c>0\).  Fit every \(\widehat h_{aj}^{(-k,A)}\) and \(\widehat q_{a\ell}^{(-k,A)}\) using only \(\mathcal D_{A,k}=\sigma(O_i:i\in\mathcal A_k)\), and define \(\widehat S_{a,j\ell}^{(-k,A)}=S_a(\widehat h_{aj}^{(-k,A)},\widehat q_{a\ell}^{(-k,A)})\), \(\widehat B^{(-k,A)}=\sum_at_a\widehat S_{a,00}^{(-k,A)}\), and the pure-outcome, pure-treatment, and four-corner dictionary \(\widehat D_{J_n}^{(-k,A)}\) by the formulas in \(\mathrm{def:population-affine-sieve}\) with population cells replaced by these learned cells.  Conditional on \(\mathcal D_{A,k}\), evaluate these fixed functions only on \(\mathcal C_k\) and form
\[
 \widehat G_k^C=|\mathcal C_k|^{-1}\sum_{i\in\mathcal C_k}
 \widehat D_{J_n}^{(-k,A)}(O_i)\widehat D_{J_n}^{(-k,A)}(O_i)^\top,
 \qquad
 \widehat c_k^C=|\mathcal C_k|^{-1}\sum_{i\in\mathcal C_k}
 \widehat D_{J_n}^{(-k,A)}(O_i)\widehat B^{(-k,A)}(O_i).
\]
If \(\widehat G_k^C=\sum_m\widehat\gamma_{km}\widehat u_{km}\widehat u_{km}^\top\), define the calibration spectral projector and coefficient by
\[
 \widehat\Pi_k^C=\sum_{m:\widehat\gamma_{km}>\tau_n}\widehat u_{km}\widehat u_{km}^\top,
 \qquad
 \widehat b^{(-k,C\mid A)}=
 \sum_{m:\widehat\gamma_{km}>\tau_n}
 \frac{\widehat u_{km}\widehat u_{km}^\top\widehat c_k^C}
      {\widehat\gamma_{km}+\rho_n}.
\]
Thus the coefficient is zero on the discarded empirical eigenspace and is uniquely ridge-regularized on the retained eigenspace.  Partition it into pure and four-corner blocks and use the deterministic formulas
\[
 \widehat\lambda_{a,j\ell}=-\widehat b^{hq}_{a,j\ell},\quad
 \widehat\lambda_{a,j0}=-\widehat b^h_{a,j}+\sum_{\ell\ge1}\widehat b^{hq}_{a,j\ell},\quad
 \widehat\lambda_{a,0\ell}=-\widehat b^q_{a,\ell}+\sum_{j\ge1}\widehat b^{hq}_{a,j\ell},
\]
\[
 \widehat\lambda_{a,00}=1+\sum_{j\ge1}\widehat b^h_{a,j}
 +\sum_{\ell\ge1}\widehat b^q_{a,\ell}
 -\sum_{j,\ell\ge1}\widehat b^{hq}_{a,j\ell}
\]
to obtain \(\widehat\lambda_{a,j\ell}^{(-k,C\mid A)}\).  For \(i\in\mathcal I_k\), evaluate only
\[
 \widehat\psi_i^{(-k,C\mid A)}=
 \sum_{a=0}^1t_a\sum_{j=0}^{J_n^h}\sum_{\ell=0}^{J_n^q}
 \widehat\lambda_{a,j\ell}^{(-k,C\mid A)}
 \widehat S_{a,j\ell}^{(-k,A)}(O_i).
\]
Finally set \(\widehat\theta_n=n^{-1}\sum_i\widehat\psi_i^{(-k(i),C\mid A)}\) and \(\widehat V_n=n^{-1}\sum_i\{\widehat\psi_i^{(-k(i),C\mid A)}-\widehat\theta_n\}^2\).  Construction, calibration, and evaluation observations are therefore separated within each evaluation fold.
\par\smallskip\noindent\emph{Inputs.} \texttt{def:\allowbreak{}score-\allowbreak{}family}; \texttt{def:\allowbreak{}normal-\allowbreak{}space}; \texttt{def:\allowbreak{}affine-\allowbreak{}score-\allowbreak{}hull}; \texttt{ass:\allowbreak{}training-\allowbreak{}measurability}; \texttt{ass:\allowbreak{}heldout-\allowbreak{}signal-\allowbreak{}integrability}.
\end{definition}

\begin{definition}[def:inference-gauges]\label{def:inference-gauges}
\par\smallskip\noindent\emph{Defined object.} \texttt{dimension-\allowbreak{}aware learned-\allowbreak{}projection gauges}
\par\smallskip\noindent\emph{Construction.} Let \(M_{q,n}\) and \(M_{y,n}\) be uniform fourth-norm envelopes for \(1-I_aq^0_{a\ell}\) and \(I_a(Y-h^0_{aj})\), respectively.  For a learned cell define
\[
 d^S_{a,j\ell,n}=M_{q,n}e_{h,4,n}+M_{y,n}e_{q,4,n}+e_{h,4,n}e_{q,4,n}.
\]
For each coordinate \(c\) of the pure and four-corner dictionary, let \(d_{c,n}\) be the sum of \(d^S_{a,j\ell,n}\) over the at most four cells entering that coordinate, and define
\[
 \delta_{D,n}=\left(\sum_{c=1}^{J_n}d_{c,n}^2\right)^{1/2}
 \le \sqrt{J_n}\max_{c\le J_n}d_{c,n},
 \qquad
 d_{0,n}=\sum_{a=0}^1d^S_{a,00,n}.
\]
Let \(B_{D,n}\) and \(B_{0,n}\) be the conditional calibration envelopes in \(\mathrm{ass:empirical-moment-perturbation}\), put \(n_{C,k}=|\mathcal C_k|\) and \(n_{C,\min}=\min_k n_{C,k}\), and set
\[
 \eta_{G,n}=B_{D,n}^2\left\{\sqrt{\frac{J_n}{n_{C,\min}}}
 +\frac{J_n}{n_{C,\min}}\right\},
 \qquad
 \eta_{c,n}=B_{D,n}B_{0,n}\sqrt{\frac{J_n}{n_{C,\min}}}.
\]
For \(G_{J_n}=E[D^0_{J_n}D_{J_n}^{0\top}]\), let \(\Lambda_{J_n}=\lambda_{\max}(G_{J_n})\) and define the total learned-to-population perturbations
\[
 \zeta_{G,n}=\eta_{G,n}+2\sqrt{J_n}B_{D,n}\delta_{D,n}+\delta_{D,n}^2,
 \qquad
 \zeta_{c,n}=\eta_{c,n}+B_{0,n}\delta_{D,n}
 +(\sqrt{J_n}B_{D,n}+\delta_{D,n})d_{0,n},
\]
and the spectral-projector gauge \(\eta_{\Pi,n}=\zeta_{G,n}/\kappa_{J_n}\).  Finally define
\[
 \Gamma_n=\alpha_{J_n}+d_{0,n}+L_n\delta_{D,n}
 +\frac{\sqrt{\Lambda_{J_n}}}{\kappa_{J_n}}
 \left\{\zeta_{c,n}+(\zeta_{G,n}+\rho_n)L_n
 +\sqrt{\Lambda_{J_n}}B_{0,n}\eta_{\Pi,n}\right\}.
\]
\par\smallskip\noindent\emph{Inputs.} \texttt{def:\allowbreak{}population-\allowbreak{}affine-\allowbreak{}sieve}; \texttt{def:\allowbreak{}affine-\allowbreak{}ensemble-\allowbreak{}estimator}.
\end{definition}

\begin{definition}[def:singular-branch-witness]\label{def:singular-branch-witness}
\par\smallskip\noindent\emph{Defined object.} \texttt{singular equal-\allowbreak{}arm branch witness}
\par\smallskip\noindent\emph{Construction.} Let \(P^\circ\) be obtained from \(P^\star\) by replacing \(P^\star(A=1\mid U^\star=u)=(1+2u)/4\) with \(P^\circ(A=1\mid U^\star=u)=1/2\) for \(u=0,1\), leaving the conditional laws of \(Z,W,Y\) unchanged.
\par\smallskip\noindent\emph{Inputs.} \texttt{def:\allowbreak{}finite-\allowbreak{}witness}.
\end{definition}

\begin{definition}[def:four-by-four-witness]\label{def:four-by-four-witness}
\par\smallskip\noindent\emph{Defined object.} \texttt{positive four-\allowbreak{}level transverse-\allowbreak{}row-\allowbreak{}space witness}
\par\smallskip\noindent\emph{Construction.} Let \(Z,W\in\{0,1,2,3\}\), put
\[
d=(1,2,3,4)^\top/10,\quad z_\circ=(1,1,1,1)^\top/4,
\]
\[
\ell_0=r_0=(1,-1,0,0)^\top,\qquad
\ell_1=r_1=(0,1,-1,0)^\top,
\]
and define
\[
B_a=z_\circ(d/2)^\top+10^{-3}\ell_ar_a^\top,
\quad a=0,1.
\]
These two matrices define the joint probabilities of \((A,Z,W)\). Conditional on \((A=a,Z=z,W=w)\), let \(Y\) be Bernoulli with mean \(h_a^\square(w)\), where
\[
h_0^\square=(1/5,2/5,3/5,4/5)^\top,qquad
h_1^\square=(1/4,9/20,13/20,17/20)^\top.
\]
Denote the resulting positive \(64\)-cell law by \(P^\square\).
\end{definition}

\section{Main results}

\begin{lemma}[lem:normal-affine-identity]\label{lem:normal-affine-identity}
For every \(a\), \(h\in\mathcal H_a(P)\), \(q\in\mathcal Q_a(P)\), \(u\in\mathcal N_a(P)\), and \(m\in\mathcal K_a(P)\) on the square-integrable product domain, \[S_a(h+u,q+m)-S_a(h,q)=F_a(h,m)-G_a(q,u)-R_a(m,u).\] Pure shifts and the four-corner contrast recover the three terms separately, so \[\mathcal C_P=\overline{\operatorname{span}}\{\phi-\phi':\phi,\phi'\in\Phi_P\}\quad\text{and}\quad\mathcal A_P=\phi_0+\mathcal C_P\] for every \(\phi_0\in\Phi_P\).
\end{lemma}
\begin{proof}
Fix an arm and suppress its subscript.  By the definitions of \(S,F,G,R\),
\[
\begin{aligned}
S(h+u,q+m)-S(h,q)
&=u+I(q+m)(Y-h-u)-Iq(Y-h)\\
&=Im(Y-h)+u-Iqu-Imu\\
&=F(h,m)-G(q,u)-R(m,u).
\end{aligned}
\]
Here \(h+u\in\mathcal H_a(P)\) because \(T_au=0\), and \(q+m\in\mathcal Q_a(P)\) because \(T_a^*m=0\).  All equalities are in \(L^2(P)\) on the stipulated product domain.  Setting \(u=0\) gives the pure treatment-bridge shift \(F(h,m)\), and setting \(m=0\) gives the pure outcome-bridge shift \(-G(q,u)\).  Moreover
\[
S(h+u,q+m)-S(h+u,q)-S(h,q+m)+S(h,q)=-R(m,u),
\]
so four bridge-score differences generate the interaction term.

For two arbitrary bridge pairs, write their differences as \(u_a\in\mathcal N_a(P)\) and \(m_a\in\mathcal K_a(P)\).  Applying the first identity arm by arm and multiplying by \(t_a\) shows that every difference \(\phi-\phi'\) belongs to \(\mathcal C_P\).  Conversely, the pure-shift and four-corner identities show that every generator \(F,G,R\), with either sign and arm sign absorbed into its scalar coefficient, lies in the linear span of such score differences.  Taking closed linear spans gives
\[
\mathcal C_P=\overline{\operatorname{span}}\{\phi-\phi':\phi,\phi'\in\Phi_P\}.
\]
Finally, the affine hull of a nonempty set with reference point \(\phi_0\) is \(\phi_0\) plus the linear span of all its pairwise differences.  Closure yields \(\mathcal A_P=\phi_0+\mathcal C_P\).
\end{proof}
\par\smallskip\noindent \textit{Justification.} The identity isolates the interaction normal missed by optimizing over one bridge coordinate at a time. \textit{Gap.} BennettEtAl2026StronglyIdentified proves the mixed-bias identity for a fixed primal-dual pair; it does not identify the closed affine geometry generated by all nonunique proximal pairs. \textit{Consumer.} The identity converts a normal-equation projection into an implementable affine ensemble of ordinary proximal doubly robust signals.

\begin{lemma}[lem:operator-kernel-perturbation]\label{lem:operator-kernel-perturbation}
Let \(A,A':\mathcal H\to\mathcal G_a\) have closed range, reduced minimum moduli at least \(\mu>0\), and \(\varepsilon=\|A'-A\|_{\mathrm{op}}<\mu\). Then
\[
\operatorname{gap}(\ker A',\ker A)\le\varepsilon/\mu,\qquad
\operatorname{gap}(\ker A'^*,\ker A^*)\le\varepsilon/\mu.
\]
With the symmetric gap convention for closed subspaces, the same bounds control the norms of the associated orthogonal-projection differences. Applied to \(A=\widetilde T_a(P)\) and \(A'=\widetilde T_a(P')\), operator-norm continuity and base-law closed range therefore imply
\(\|\widetilde\Pi_{N,a}(P')-\widetilde\Pi_{N,a}(P)\|_{\mathrm{op}}\to0\) and
\(\|\widetilde\Pi_{K,a}(P')-\widetilde\Pi_{K,a}(P)\|_{\mathrm{op}}\to0\) along the repaired stratum. When these projections have finite rank, a bound below one also forces equality of their ranks. No differentiability of projections or bridge selections, solvability-submanifold claim, or path-score realization follows from these bounds.
\end{lemma}
\begin{proof}
Let \(x\in\ker A'\) with \(\|x\|=1\), and decompose \(x=x_0+x_1\) with \(x_0\in\ker A\) and \(x_1\perp\ker A\).  The reduced-modulus bound for \(A\) gives
\[
\mu\|x_1\|\le\|Ax_1\|=\|Ax\|=\|(A-A')x\|\le\varepsilon.
\]
Thus \(\operatorname{dist}(x,\ker A)\le\varepsilon/\mu\).  Interchanging \(A,A'\) proves the reverse directed-gap bound.  Applying the same argument to \(A^*,A'^*\) is legitimate because a closed-range operator and its adjoint have the same positive reduced minimum modulus: on \((\ker A)^{\perp}\), polar decomposition identifies \(A^*A\) and \(AA^*\) on their nonzero spectral subspaces.  This proves both displayed gap bounds.

For orthogonal projections \(P_M,P_N\), decompose a unit vector into its \(M\) and \(M^\perp\) components and apply the two directed-gap bounds; equivalently, the two off-diagonal restrictions of \(P_M-P_N\) have norms equal to the directed gaps.  Hence
\[
\|P_M-P_N\|_{\mathrm{op}}
=\max\{\operatorname{gap}(M,N),\operatorname{gap}(N,M)\}.
\]
Take \(A=\widetilde T_a(P)\), \(A'=\widetilde T_a(P')\).  The repaired stratum gives reduced moduli bounded below by \(\mu_a(P)/2\), while transported operator continuity makes \(\varepsilon\to0\); the two projection differences therefore converge to zero.  Finally, if \(\|P_M-P_N\|<1\), then \(P_N|_M\) and \(P_M|_N\) are injective.  When the ranks are finite, the two injections imply \(\dim M=\dim N\).  None of these arguments constructs a statistical path or differentiates a projection.
\end{proof}
\par\smallskip\noindent \textit{Justification.} The perturbation bound derives moving kernel and cokernel continuity from the repaired spectral stratum. \textit{Gap.} Subspace continuity does not establish path-score realization. \textit{Consumer.} The result supplies stable transported null coordinates without treating their differentiability as assumed.

\begin{openendedquestion}[oeq:score-lifting-completeness]\label{oeq:score-lifting-completeness}
For the saturated likelihood-ratio domain \(\mathcal V_P\), the full fixed-defect separated stratum \(\mathcal U_P\), and the actual path class in \(\mathcal M_{\mathrm{joint}}(P)\), does every \(s\in\mathcal C_P^\perp\) lie in the \(L^2(P)\)-closure of finite linear combinations of scores of admissible DQM paths, including at intersections with moving kernels and a shared \((W,X)\) marginal? Equivalently, does
\[
\mathcal T_P=\mathcal C_P^\perp
\]
hold on every such joint solvability germ?  Neither equality nor bridge-solvability path realization is included in the chart definition.
\end{openendedquestion}
\par\smallskip\noindent \textit{Justification.} This remains the reverse-inclusion problem for the full saturated fixed-defect intersection, not for an analyst-selected subchart. \textit{Gap.} The finite constant-rank and branch calculations do not provide a general coupled Hilbert-space incidence chart, and operator continuity alone supplies no differentiable path lift. \textit{Consumer.} Its resolution determines whether the affine projection is the information bound throughout the unrestricted joint model.

\begin{proposition}[prop:causal-identification]\label{prop:causal-identification}
Under the latent-extension assumptions, \(h_a^\dagger\) satisfies \(E[Y\mid A=a,Z,X]=E[h_a^\dagger(W,X)\mid A=a,Z,X]\). Hence \(h_a^\dagger\in\mathcal H_a(P)\). If treatment bridges exist, then every \(h_a\in\mathcal H_a(P)\) has \(E[h_a(W,X)]=E[h_a^\dagger(W,X)]=E[Y(a)]\), and therefore the observed joint-bridge functional obeys \(\theta(P)=E[Y(1)-Y(0)]\).
\end{proposition}
\begin{proof}
Fix \(a\in\{0,1\}\).  Assumption \texttt{ass:latent-positivity} gives
\[
P(A=a)=E\{P(A=a\mid U,X)\}>0,
\]
so all arm-specific conditional expectations below are well defined up to
\(P_{Z,X\mid A=a}\)-null sets.  The integrability needed for the conditional
expectations follows from \(h_a^\dagger\in\mathcal H=L^2(P_{W,X})\), from the
conditional mean appearing in \texttt{ass:latent-bridge-compatibility}, and
from the second-moment clause of \texttt{def:joint-model}.

On \(\{A=a\}\), assumption \texttt{ass:causal-consistency} gives
\(Y=Y(A)=Y(a)\).  Successive applications of the tower property, followed in
order by \texttt{ass:latent-exchangeability},
\texttt{ass:latent-bridge-compatibility}, and
\texttt{ass:w-treatment-negative-control}, yield
\[
\begin{aligned}
E[Y\mid A=a,Z,X]
&=E[Y(a)\mid A=a,Z,X]\\
&=E\!\left[E\{Y(a)\mid U,X,A,Z\}\mid A=a,Z,X\right]\\
&=E\!\left[E\{Y(a)\mid U,X\}\mid A=a,Z,X\right]\\
&=E\!\left[E\{h_a^\dagger(W,X)\mid U,X\}\mid A=a,Z,X\right]\\
&=E\!\left[E\{h_a^\dagger(W,X)\mid U,X,A,Z\}\mid A=a,Z,X\right]\\
&=E[h_a^\dagger(W,X)\mid A=a,Z,X].
\end{aligned}
\]
Indeed, the fourth equality uses that conditional independence of \(W\) and
\((A,Z)\) given \((U,X)\) implies the same conditional-mean equality for the
square-integrable measurable function \(h_a^\dagger(W,X)\).  By the
definitions of \(g_a\) and \(T_a\), the displayed identity is
\(T_a h_a^\dagger=g_a\).  Since \(h_a^\dagger\in\mathcal H\), the definition
of \(\mathcal H_a(P)\) therefore gives
\(h_a^\dagger\in\mathcal H_a(P)\).

Now let \(h_a\in\mathcal H_a(P)\) be arbitrary and put
\(d_a=h_a-h_a^\dagger\in\mathcal H\).  The preceding conclusion and the
definition of \(\mathcal H_a(P)\) imply
\[
T_a d_a=T_a h_a-T_a h_a^\dagger=g_a-g_a=0
\quad\text{in }\mathcal G_a.
\]
At the base law, assumption \texttt{ass:treatment-bridge-existence}, as
incorporated in \texttt{def:joint-model}, permits a choice
\(q_a\in\mathcal Q_a(P)\).  By the definition of \(\mathcal Q_a(P)\),
\(T_a^*q_a=1\) in \(\mathcal H\).  Because \(d_a\in\mathcal H\),
\(q_a\in\mathcal G_a\), and \(T_a\) is bounded, all the following inner
products are finite.  The adjoint identity then gives
\[
\begin{aligned}
E[h_a(W,X)-h_a^\dagger(W,X)]
&=\langle 1,d_a\rangle_{\mathcal H}\\
&=\langle T_a^*q_a,d_a\rangle_{\mathcal H}\\
&=\langle q_a,T_a d_a\rangle_{\mathcal G_a}\\
&=0.
\end{aligned}
\]
Thus every outcome-bridge representative has the same marginal mean as
\(h_a^\dagger\); no uniqueness of either bridge was used.

Finally, the tower property and
\texttt{ass:latent-bridge-compatibility} imply
\[
E[h_a^\dagger(W,X)]
=E\!\left[E\{h_a^\dagger(W,X)\mid U,X\}\right]
=E\!\left[E\{Y(a)\mid U,X\}\right]
=E[Y(a)].
\]
Consequently, for arbitrary \(h_a\in\mathcal H_a(P)\),
\[
E[h_a(W,X)]=E[h_a^\dagger(W,X)]=E[Y(a)].
\]
Applying this identity for \(a=1\) and \(a=0\), and then using the declared
definition of the observed joint-bridge functional \(\theta\), gives
\[
\theta(P)
=E[h_1(W,X)-h_0(W,X)]
=E[Y(1)]-E[Y(0)]
=E[Y(1)-Y(0)],
\]
where the last equality follows from integrability and linearity of
expectation.  This proves both bridge membership and the claimed causal
interpretation.
\end{proof}
\par\smallskip\noindent \textit{Justification.} This proposition separates an observed-law efficiency problem from the assumptions needed to interpret its target as an ATE. \textit{Gap.} MiaoGengTchetgen2018Proxy gives proximal causal identification; the proposition records exactly which latent extension makes the observed joint-bridge functional causal without adding bridge uniqueness or completeness. \textit{Consumer.} The SUPPORT analysis may invoke the quotient efficiency results for an ATE only after checking this proximal causal interpretation.

\begin{theorem}[thm:quotient-efficiency]\label{thm:quotient-efficiency}
Suppose that \(\mathcal T_P=\mathcal C_P^\perp\).  Suppose also that there are bridges \(h_a\in\mathcal H_a(P)\) and \(q_a\in\mathcal Q_a(P)\) for which \(\phi_{h,q}\in L_0^2(P)\), and that the following pathwise product condition holds.  For every admissible DQM path \(t\mapsto P_t\) with score \(s\), there are \(h_{a,t}\in\mathcal H_a(P_t)\) and \(q_{a,t}\in\mathcal Q_a(P_t)\) such that, with \(T_{a,t}\) and \(T_{a,t}^*\) denoting the operators under \(P_t\),
\[
 \min\!\left\{
 \|T_{a,t}(h_{a,t}-h_a)\|_{\mathcal G_a(P_t)}
 \|q_{a,t}-q_a\|_{\mathcal G_a(P_t)},
 \|h_{a,t}-h_a\|_{\mathcal H(P_t)}
 \|T_{a,t}^*(q_{a,t}-q_a)\|_{\mathcal H(P_t)}
 \right\}=o(|t|)
\]
for \(a=0,1\).  Then \(\theta\) is pathwise differentiable at \(P\), and \(\phi_{h,q}\) represents its derivative along every actual path.  Moreover, every gradient is in \(\mathcal A_P\), and, for any \(\phi_0\in\Phi_P\),
\[
 \phi_{\mathrm{eff}}
 =\phi_0-\Pi_{\mathcal C_P}\phi_0
 =\operatorname*{argmin}_{\varphi\in\mathcal A_P}E[\varphi^2],
 \qquad
 V_{\mathrm{eff}}(P)=\inf_{\varphi\in\mathcal A_P}E[\varphi^2].
\]
No converse from pathwise differentiability to \(\Phi_P\neq\varnothing\) is asserted.
\end{theorem}
\begin{proof}
Fix an admissible DQM path \(t\mapsto P_t\) through \(P\), and denote its score by \(s\).  Put \(r_t=dP_t/dP\) and \(z_t=\sqrt{r_t}-1\).  The DQM expansion in \(\mathrm{def:joint-model}\) gives
\[
 z_t=\frac{t}{2}s+o_{L^2(P)}(t),
 \qquad \|z_t\|_{L^2(P)}=O(|t|).
\]
Because every sufficiently local \(P_t\) belongs to \(\mathcal V_P\), \(\mathrm{def:local-operator-chart}\) and \(\mathrm{ass:local-uniform-equivalence}\) give \(1-\varepsilon<r_t<1+\varepsilon\), \(P\)-almost surely.  Consequently \(\sup_{|t|<t_0}\|z_t\|_{L^\infty(P)}<\infty\) for some \(t_0>0\).  We first verify the expectation derivative needed below.  For fixed \(f\in L^2(P)\),
\[
 E_{P_t}f-E_Pf
 =E_P\{f(r_t-1)\}
 =2E_P(fz_t)+E_P(fz_t^2).
\]
The first term is \(tE_P(fs)+o(t)\) by the DQM expansion and Cauchy--Schwarz.  For every \(M>0\),
\[
 \frac{E_P(|f|z_t^2)}{|t|}
 \le
 M\frac{\|z_t\|_2^2}{|t|}
 +\|f\mathbf 1\{|f|>M\}\|_2
   \frac{\|z_t\|_\infty\|z_t\|_2}{|t|}.
\]
For fixed \(M\), the first term tends to zero.  The limsup of the second term is bounded by a constant times \(\|f\mathbf 1\{|f|>M\}\|_2\), which tends to zero as \(M\to\infty\).  Hence
\[
 E_{P_t}f-E_Pf=tE_P(fs)+o(t).                                      \tag{1}
\]
This proves (1) for every fixed square-integrable \(f\), without a boundedness condition on \(f\).

Choose the bridge representatives in the theorem and set
\[
 B_{h,q}=\sum_{a=0}^1t_aS_a(h_a,q_a)
          =\theta(P)+\phi_{h,q}.
\]
The equality follows from \(\mathrm{def:score-family}\), and the base bridge equations give \(E_P B_{h,q}=\theta(P)\): indeed,
\[
 E_P[I_aq_a\{Y-h_a(W,X)\}]
 =E_P[I_aq_a\,E_P\{Y-h_a(W,X)\mid A=a,Z,X\}]=0.
\]
The hypothesis \(\phi_{h,q}\in L_0^2(P)\) implies \(B_{h,q}\in L^2(P)\).  Applying (1), and using \(E_Ps=0\), yields
\[
 E_{P_t}B_{h,q}-E_PB_{h,q}
 =tE_P(\phi_{h,q}s)+o(t).                                           \tag{2}
\]

It remains to compare the fixed-bridge expectation in (2) with the moving target.  For each arm, take the representatives \(h_{a,t}\in\mathcal H_a(P_t)\) and \(q_{a,t}\in\mathcal Q_a(P_t)\) supplied by the pathwise product condition, and write
\(\theta_a(P_t)=E_{P_t}h_{a,t}(W,X)\).  This arm mean is independent of the selected outcome bridge.  In fact, if \(h'_{a,t}\) is another solution and \(u_t=h'_{a,t}-h_{a,t}\), then \(T_{a,t}u_t=0\), while \(T_{a,t}^*q_{a,t}=1\); therefore
\[
 E_{P_t}u_t
 =E_{P_t}(I_aq_{a,t}u_t)
 =E_{P_t}\{I_aq_{a,t}T_{a,t}u_t\}=0.                               \tag{3}
\]
All terms in (3) are integrable by Cauchy--Schwarz.  The same justification applies below: membership of the displayed bridges in their \(L^2(P_t)\) spaces supplies the moving terms, and the likelihood-ratio bound transfers the fixed bridges from \(L^2(P)\) to \(L^2(P_t)\).

Let \(\Delta h_{a,t}=h_{a,t}-h_a\).  Outcome-bridge solvability and treatment-bridge solvability under \(P_t\), respectively, give
\[
 E_{P_t}[I_aq_a\{Y-h_{a,t}(W,X)\}]=0,
 \qquad
 E_{P_t}[(1-I_aq_{a,t})\Delta h_{a,t}(W,X)]=0.                       \tag{4}
\]
Using \(Y-h_a=(Y-h_{a,t})+\Delta h_{a,t}\) in the definition of \(S_a\), the two identities in (4) imply the exact mixed-bias identity
\[
 \begin{aligned}
 \theta_a(P_t)-E_{P_t}S_a(h_a,q_a)
 &=E_{P_t}[(1-I_aq_a)\Delta h_{a,t}]\\
 &=E_{P_t}[I_a(q_{a,t}-q_a)\Delta h_{a,t}].                         \tag{5}
 \end{aligned}
\]
Conditioning the last expectation first on \((A,Z,X)\) and then, separately, on \((W,X)\), and applying Cauchy--Schwarz, gives
\[
 \begin{aligned}
 |\theta_a(P_t)-E_{P_t}S_a(h_a,q_a)|
 &\le
 \|T_{a,t}\Delta h_{a,t}\|_{\mathcal G_a(P_t)}
 \|q_{a,t}-q_a\|_{\mathcal G_a(P_t)},\\
 |\theta_a(P_t)-E_{P_t}S_a(h_a,q_a)|
 &\le
 \|\Delta h_{a,t}\|_{\mathcal H(P_t)}
 \|T_{a,t}^*(q_{a,t}-q_a)\|_{\mathcal H(P_t)}.                    \tag{6}
 \end{aligned}
\]
The stated product condition says that the smaller of the two right-hand sides in (6) is \(o(|t|)\), for each of the two arms.  Taking the signed arm contrast in (5), and combining it with (2), proves
\[
 \theta(P_t)-\theta(P)=tE_P(\phi_{h,q}s)+o(t).                      \tag{7}
\]
Thus \(\theta\) is pathwise differentiable on the actual score set, with gradient \(\phi_{h,q}\).  The continuous linear functional \(v\mapsto E_P(\phi_{h,q}v)\) extends (7) from finite linear combinations of actual scores to their \(L^2(P)\)-closure \(\mathcal T_P\).

Fix any \(\phi_0\in\Phi_P\).  By \(\mathrm{lem:normal-affine-identity}\),
\[
 \mathcal A_P=\phi_0+\mathcal C_P.                                 \tag{8}
\]
Define
\[
 \phi_* =\phi_0-\Pi_{\mathcal C_P}\phi_0.
\]
Because \(-\Pi_{\mathcal C_P}\phi_0\in\mathcal C_P\), equation (8) gives \(\phi_*\in\mathcal A_P\); because it is the orthogonal-projection residual, \(\phi_*\perp\mathcal C_P\).  For every \(c\in\mathcal C_P\), Pythagoras gives
\[
 \|\phi_0+c\|_2^2
 =\|\phi_*\|_2^2+\|\Pi_{\mathcal C_P}\phi_0+c\|_2^2.           \tag{9}
\]
Hence \(\phi_*\) is the unique minimum-norm element of \(\mathcal A_P\).

Now use the theorem's tangent hypothesis \(\mathcal T_P=\mathcal C_P^\perp\).  It gives \(\phi_*\in\mathcal T_P\).  Moreover, every member of \(\Phi_P\), including \(\phi_0\), is a gradient: its difference from the already verified \(\phi_{h,q}\) belongs to \(\mathcal C_P=\mathcal T_P^\perp\) by (8).  Therefore, for every \(v\in\mathcal T_P\),
\[
 E_P(\phi_*v)
 =E_P(\phi_0v)-E_P\{(\Pi_{\mathcal C_P}\phi_0)v\}
 =\dot\theta_P(v).                                                  \tag{10}
\]
Thus \(\phi_*\) is the gradient lying in the tangent space, hence the canonical gradient.  If \(g\in L_0^2(P)\) is any gradient, then (10) and the gradient property of \(\phi_0\) imply \(g-\phi_0\perp\mathcal T_P\).  Since
\[
 \mathcal T_P^\perp=(\mathcal C_P^\perp)^\perp=\mathcal C_P,
\]
closedness of \(\mathcal C_P\) gives \(g\in\phi_0+\mathcal C_P=\mathcal A_P\).  Conversely, adding an element of \(\mathcal C_P=\mathcal T_P^\perp\) to \(\phi_0\) leaves all path derivatives unchanged, so every element of \(\mathcal A_P\) is a gradient.  Equations (9)--(10) and \(\mathrm{def:quotient-gradient}\) now yield
\[
 \phi_{\mathrm{eff}}
 =\phi_0-\Pi_{\mathcal C_P}\phi_0
 =\operatorname*{argmin}_{\varphi\in\mathcal A_P}E_P(\varphi^2),
 \qquad
 V_{\mathrm{eff}}(P)=\inf_{\varphi\in\mathcal A_P}E_P(\varphi^2).
\]
The argument starts from an explicitly assumed square-integrable bridge score and therefore makes no converse assertion that pathwise differentiability forces \(\Phi_P\neq\varnothing\).
\end{proof}
\par\smallskip\noindent \textit{Justification.} The Hilbert projection and complete gradient-affine-space calculation are proved once tangent equality and bridge-score pathwise validity are available. \textit{Gap.} Neither reverse score lifting nor unbounded-product bridge-score validity, including the claimed necessity of \(\Phi_P\ne\varnothing\), follows from the frozen assumptions. \textit{Consumer.} Analysts must use the unconditional positive finite results or state both missing premises explicitly.

\begin{theorem}[thm:single-pair-attainment]\label{thm:single-pair-attainment}
Conditional on \(\mathrm{oeq:score-lifting-completeness}\), fix \(\phi_0=\phi_{h,q}\in\Phi_P\). A shifted pair \((h_a+u_a,q_a+m_a)_{a=0}^1\) attains \(V_{\mathrm{eff}}(P)\) if and only if there are \(u_a\in\mathcal N_a(P)\) and \(m_a\in\mathcal K_a(P)\) such that \[\Pi_{\mathcal C_P}\phi_0=\sum_{a=0}^1t_a\{-F_a(h_a,m_a)+G_a(q_a,u_a)+R_a(m_a,u_a)\}.\] Otherwise every single-pair score has variance strictly above the affine-hull bound, while four-corner affine ensembles can approximate it.
\end{theorem}
\begin{proof}
Fix \(\phi_0=\phi_{h,q}\in\Phi_P\).  Null shifts preserve the bridge equations: if \(u_a\in\mathcal N_a(P)=\ker T_a\) and \(m_a\in\mathcal K_a(P)=\ker T_a^*\), then
\[
 h_a+u_a\in\mathcal H_a(P),
 \qquad
 q_a+m_a\in\mathcal Q_a(P).
\]
Direct expansion of \(S_a(h_a+u_a,q_a+m_a)-S_a(h_a,q_a)\), the algebraic identity underlying \(\mathrm{lem:normal-affine-identity}\), gives pointwise
\[
 \phi_{h+u,q+m}
 =\phi_0+
 \sum_{a=0}^1t_a\{F_a(h_a,m_a)-G_a(q_a,u_a)-R_a(m_a,u_a)\}.          \tag{11}
\]
By \(\mathrm{lem:normal-affine-identity}\), \(\mathcal A_P=\phi_0+\mathcal C_P\).  Since \(\mathcal C_P\) is closed, orthogonal projection and Pythagoras show directly that the unique minimum-norm element of this affine set, as named in \(\mathrm{def:quotient-gradient}\), is
\[
 \phi_{\mathrm{eff}}=\phi_0-\Pi_{\mathcal C_P}\phi_0.            \tag{12}
\]
The shifted pair attains \(V_{\mathrm{eff}}(P)\) exactly when its score equals the unique minimizer (12).  Subtracting \(\phi_0\) from (11)--(12) shows that this happens if and only if
\[
 \Pi_{\mathcal C_P}\phi_0
 =\sum_{a=0}^1t_a\{-F_a(h_a,m_a)+G_a(q_a,u_a)+R_a(m_a,u_a)\},        \tag{13}
\]
as an equality in \(L^2(P)\).  Necessity is legitimate even when the three displayed products are not separately square-integrable, because their sum in (11) is the difference of two square-integrable scores.  Conversely, if null shifts satisfy (13) as an \(L^2(P)\) equality, then the pointwise identity (11) makes their combined score equal to \(\phi_{\mathrm{eff}}\), hence square-integrable and attaining the bound.

If (13) has no solution, every valid single-pair score \(\varphi\in\Phi_P\) differs from \(\phi_{\mathrm{eff}}\).  Since \(\varphi-\phi_{\mathrm{eff}}\in\mathcal C_P\) and \(\phi_{\mathrm{eff}}\perp\mathcal C_P\),
\[
 E_P(\varphi^2)
 =V_{\mathrm{eff}}(P)+\|\varphi-\phi_{\mathrm{eff}}\|_2^2
 >V_{\mathrm{eff}}(P).                                             \tag{14}
\]
Finally, \(\phi_{\mathrm{eff}}\in\overline{\operatorname{aff}}(\Phi_P)\) by \(\mathrm{def:affine-score-hull}\), so a sequence of finite affine combinations of valid bridge scores converges to it in \(L^2(P)\).  The pure-shift and four-corner identities in \(\mathrm{lem:normal-affine-identity}\) express the generating \(F_a\), \(G_a\), and \(R_a\) directions by such finite bridge-score contrasts.  Thus the approximating combinations may be realized as the asserted four-corner affine ensembles.
\end{proof}
\par\smallskip\noindent \textit{Justification.} The factorization equation determines when choosing one representative can equal global across-pair efficiency. \textit{Gap.} ZhangLiMiaoTchetgen2023Nonuniqueness establishes locally efficient regular inference under its additional regularity, while CuiEtAl2024SemiparametricProximal obtains single-pair efficiency under uniqueness and surjectivity; neither gives this two-fiber attainment criterion. \textit{Consumer.} The SUPPORT workflow can test whether fitted bridge scores leave variance-reducing affine directions.

\begin{theorem}[thm:finite-witness-nonattainment]\label{thm:finite-witness-nonattainment}
The law \(P^\star\) lies in the joint chart, has positive mass on all \(36\) observed cells, and has rank-two arm operators with \(\mathcal N_a(P^\star)=\mathcal K_a(P^\star)=\operatorname{span}\{(3,-7,3)^\top\}\). For the displayed constraint map \(\Xi\), the Jacobian has constant full row rank on a relatively open positive neighborhood of \(P^\star\); only on this finite chart the joint solvability locus is thereby proved to be a submanifold, with \(\mathcal T_P=\mathcal C_P^\perp\). At \(P^\star\), \(\theta(P^\star)=1/5\), \(V_{\mathrm{eff}}(P^\star)=773723312/334026875\), and the exact coefficient obstruction \(410302072/87437547015>0\) in each arm implies that no single bridge pair attains the bound. The strict nonattainment persists on a smaller relatively open chart.
\end{theorem}
\begin{proof}
Summing out \(U^\star\) gives minimum observed cell probability \(9/1000\) and the two arm matrices
\[
B_0=\begin{pmatrix}
109/800&57/800&3/100\\
57/800&9/200&27/800\\
3/100&27/800&39/800
\end{pmatrix},
\quad
B_1=\begin{pmatrix}
39/800&27/800&3/100\\
27/800&9/200&57/800\\
3/100&57/800&109/800
\end{pmatrix}.
\]
Both have rank two, and direct multiplication gives
\[
B_av=0,\qquad B_a^\top v=0,
\quad v=(3,-7,3)^\top.
\]
Their common column-sum vector is \(d=(7/20,3/10,7/20)^\top\).  The outcome moment vectors are
\[
b_0=(21/400,9/200,21/400)^\top,
\qquad b_1=(3/50,21/200,37/200)^\top.
\]
Multiplication verifies \(B_ah_a^{\mathrm{sel}}=b_a\) and
\(B_a^\top q_a^{\mathrm{sel}}=d\).  Hence both bridge fibers are nonempty, and their directions are exactly \(\operatorname{span}\{v\}\).  Finite support gives all required moments and a positive reduced minimum modulus.  Therefore \(P^\star\) belongs to its repaired saturated joint chart.  Since \(h_1^{\mathrm{sel}}-h_0^{\mathrm{sel}}=1/5\), \(\theta(P^\star)=1/5\).

Use the ordered normal vector \(D=(F_0,G_0,R_0,F_1,G_1,R_1)^\top\) based on these selected bridges and \(v\).  Exact cell summation gives
\[
\Gamma=E[DD^\top]=
\begin{pmatrix}
1323/500&-6/5&63/20&0&0&0\\
-6/5&343/5&-21&0&-21&0\\
63/20&-21&441/2&0&0&0\\
0&0&0&1323/500&6/5&-63/20\\
0&-21&0&6/5&343/5&-21\\
0&0&0&-63/20&-21&441/2
\end{pmatrix},
\]
whose determinant is \(16253767998768573/12500000>0\), and
\[
c=E[D\phi_{\mathrm{sel}}]
=(24/125,8/25,6/5,-24/125,8/25,6/5)^\top.
\]
The differentials of the rank, outcome-solvability, and treatment-solvability Schur equations are nonsingular scalar changes of, respectively, the \(R,F,G\) covectors.  Positive definiteness of \(\Gamma\) therefore proves that \(D\Xi(P^\star)\) has full row rank six.  The finite constraint-tangent lemma supplies actual positive paths for every vector in the common kernel, so
\(\mathcal T_{P^\star}=\mathcal C_{P^\star}^\perp\) unconditionally.

Solving the normal equations gives
\[
\beta=\Gamma^{-1}c=\left(
\frac{16264}{229047},\frac{12512}{1145235},\frac{131512}{24049935},
-\frac{16264}{229047},\frac{12512}{1145235},\frac{131512}{24049935}
\right)^\top.
\]
Since \(E[\phi_{\mathrm{sel}}^2]=4432/1875\), orthogonal projection yields
\[
V_{\mathrm{eff}}(P^\star)
=\frac{4432}{1875}-c^\top\Gamma^{-1}c
=\frac{773723312}{334026875}.
\]

Every bridge pair is \(h_a^{\mathrm{sel}}+c_av\), \(q_a^{\mathrm{sel}}+d_av\).  Relative to \((F_a,G_a,R_a)\), its score change has coefficients
\[
\alpha_a=t_ad_a,\qquad \gamma_a=-t_ac_a,
\qquad \zeta_a=-t_ac_ad_a=t_a\alpha_a\gamma_a.
\]
To equal the projected score, the coefficient vector \(-\beta\) must satisfy this quadratic identity.  For each arm, exact substitution instead gives
\[
\beta_{R,a}+t_a\beta_{F,a}\beta_{G,a}
=\frac{410302072}{87437547015}>0.
\]
Because \(\Gamma\) is positive definite, coefficient equality is equivalent to function equality; hence no single pair attains the bound.  The Schur minor, \(D\Xi\), \(\Gamma\), and the displayed obstruction vary continuously in the positive cell probabilities.  Their nonzero determinants and strict inequality persist after shrinking the relatively open rank-two chart, proving local persistence.
\end{proof}
\par\smallskip\noindent \textit{Justification.} Exact matrices, Gram determinant, projection coefficients, variance, and factorization obstruction are reproduced. \textit{Gap.} The result is a positive finite fixed-rank theorem and does not settle the general lifting question. \textit{Consumer.} It is a reproducible certificate that no individual bridge pair need attain the affine bound.

\begin{theorem}[thm:selected-pair-strict-gap]\label{thm:selected-pair-strict-gap}
At \(P^\star\), the displayed functions satisfy \(T_ah_a^{\mathrm{sel}}=g_a\) and \(T_a^*q_a^{\mathrm{sel}}=1\). Their score has
\[
V_{\mathrm{sel}}=\frac{4432}{1875}
=V_{\mathrm{eff}}(P^\star)+\frac{9496288}{200416125}
>V_{\mathrm{eff}}(P^\star).
\]
Fixing one nonzero-null-coordinate gauge for each outcome and treatment bridge gives smooth rational continuations of these representatives on a sufficiently small positive rank-two neighborhood, and continuity of the finite normal Gram matrix preserves the strict variance gap there. This is an exact selected-score comparison; it does not assert that any cited estimator has this selected limit or import an efficiency claim from a different statistical model.
\end{theorem}
\begin{proof}
Write vectors in the coordinate order \(0,1,2\).  By \(\mathrm{def:finite-witness}\) and the exact marginalization in \(\mathrm{thm:finite-witness-nonattainment}\),
\[
B_0=\begin{pmatrix}
109/800&57/800&3/100\\
57/800&9/200&27/800\\
3/100&27/800&39/800
\end{pmatrix},\qquad
B_1=\begin{pmatrix}
39/800&27/800&3/100\\
27/800&9/200&57/800\\
3/100&57/800&109/800
\end{pmatrix},
\]
\[
b_0=\begin{pmatrix}21/400\\9/200\\21/400\end{pmatrix},\qquad
b_1=\begin{pmatrix}3/50\\21/200\\37/200\end{pmatrix},\qquad
d=\begin{pmatrix}7/20\\3/10\\7/20\end{pmatrix}.
\]
The displayed selected representatives are
\[
h_0^{\mathrm{sel}}=\begin{pmatrix}0\\2/5\\4/5\end{pmatrix},\quad
h_1^{\mathrm{sel}}=\begin{pmatrix}1/5\\3/5\\1\end{pmatrix},\quad
q_0^{\mathrm{sel}}=\begin{pmatrix}0\\8/3\\16/3\end{pmatrix},\quad
q_1^{\mathrm{sel}}=\begin{pmatrix}16/3\\8/3\\0\end{pmatrix}.
\]
Direct multiplication gives the four identities
\[
B_0h_0^{\mathrm{sel}}=b_0,\qquad
B_1h_1^{\mathrm{sel}}=b_1,
\qquad
B_0^\top q_0^{\mathrm{sel}}=d,\qquad
B_1^\top q_1^{\mathrm{sel}}=d.
\]
Indeed, the first two products have numerators, after use of a common denominator \(4000\), equal respectively to
\[
(210,180,210)^\top/4000=b_0,qquad
(240,420,740)^\top/4000=b_1,
\]
and the last two products both equal \((7/20,3/10,7/20)^\top\).

We now translate these unconditional moment equations into the asserted conditional bridge equations.  Put
\(r_a(z)=P^\star(A=a,Z=z)=\sum_wB_a[z,w]\).  Positivity of all \(36\) observed cells, proved in \(\mathrm{thm:finite-witness-nonattainment}\), implies \(r_a(z)>0\) and \(d(w)=P^\star(W=w)>0\).  By the definitions of \(T_a\), \(g_a\), \(B_a\), and \(b_a\),
\[
(T_ah)(z)=\frac{(B_ah)[z]}{r_a(z)},
\qquad
g_a(z)=\frac{b_a[z]}{r_a(z)}.
\]
Consequently \(B_ah_a^{\mathrm{sel}}=b_a\) implies
\(T_ah_a^{\mathrm{sel}}=g_a\).  Similarly, the definition of the adjoint gives
\[
(T_a^*q)(w)
=E_{P^\star}[I_aq(Z)\mid W=w]
=\frac{(B_a^\top q)[w]}{d(w)}.
\]
Thus \(B_a^\top q_a^{\mathrm{sel}}=d\), together with \(d(w)>0\), implies
\(T_a^*q_a^{\mathrm{sel}}=1\).  This proves all four bridge restrictions without invoking uniqueness.

For the variance calculation, \(h_1^{\mathrm{sel}}(w)-h_0^{\mathrm{sel}}(w)=1/5\) for every \(w\), and \(\mathrm{thm:finite-witness-nonattainment}\) gives \(\theta(P^\star)=1/5\).  Hence the centered selected score reduces to
\[
\phi_{\mathrm{sel}}
=I_1q_1^{\mathrm{sel}}(Z)\{Y-h_1^{\mathrm{sel}}(W)\}
-I_0q_0^{\mathrm{sel}}(Z)\{Y-h_0^{\mathrm{sel}}(W)\}
=t_Aq_A^{\mathrm{sel}}(Z)\{Y-h_A^{\mathrm{sel}}(W)\}.
\]
Under \(\mathrm{def:finite-witness}\), conditional on \((A=a,U^\star=u)\), the variables \(Z\), \(W\), and \(Y\) are mutually independent, the law of each of \(Z,W\) is
\[
p_0=(3/5,3/10,1/10),\qquad p_1=(1/10,3/10,3/5),
\]
and
\[
\mu_{a,u}:=P^\star(Y=1\mid A=a,U^\star=u)=\frac{1+a+2u}{5}.
\]
Therefore
\[
E[(q_a^{\mathrm{sel}}(Z))^2(Y-h_a^{\mathrm{sel}}(W))^2\mid A=a,U^\star=u]
=Q_{a,u}H_{a,u},
\]
where
\[
Q_{a,u}=\sum_{z=0}^2p_u(z)(q_a^{\mathrm{sel}}(z))^2,
\qquad
H_{a,u}=\sum_{w=0}^2p_u(w)
\left[\mu_{a,u}\{1-h_a^{\mathrm{sel}}(w)\}^2
+(1-\mu_{a,u})\{h_a^{\mathrm{sel}}(w)\}^2\right].
\]
Substitution of the displayed vectors gives the following exact table:
\[
\begin{array}{c|c|c|c|c|c}
u&a&P^\star(U^\star=u,A=a)&Q_{a,u}&H_{a,u}&P^\star(U^\star=u,A=a)Q_{a,u}H_{a,u}\\ \hline
0&0&3/8&224/45&29/125&812/1875\\
0&1&1/8&96/5&39/125&468/625\\
1&0&1/8&96/5&39/125&468/625\\
1&1&3/8&224/45&29/125&812/1875
\end{array}
\]
Since \(I_0I_1=0\), summing the last column proves
\[
V_{\mathrm{sel}}
=E_{P^\star}[\phi_{\mathrm{sel}}^2]
=2\frac{812}{1875}+2\frac{468}{625}
=\frac{4432}{1875}.
\]
The established witness theorem gives
\(V_{\mathrm{eff}}(P^\star)=773723312/334026875\).  Exact subtraction yields
\[
V_{\mathrm{sel}}-V_{\mathrm{eff}}(P^\star)
=\frac{4432}{1875}-\frac{773723312}{334026875}
=\frac{9496288}{200416125}>0,
\]
which proves the claimed strict comparison at \(P^\star\).

It remains to justify persistence, including the asserted gauge construction.  Work in the relatively open positive rank-two joint-solvability chart supplied by \(\mathrm{def:finite-constraint-map}\) and \(\mathrm{thm:finite-witness-nonattainment}\).  At \(P^\star\), both \(\ker(B_a)\) and \(\ker(B_a^\top)\) are spanned by
\(v=(3,-7,3)^\top\), so every coordinate used below is nonzero in the relevant null vector.

For the outcome bridge in arm \(a\), impose the gauge
\[
h_a(P)(0)=h_a^{\mathrm{sel}}(0)=a/5.
\]
Because \(v(0)\ne0\) and \(B_a(P^\star)\) has rank two, the two columns indexed by \(w=1,2\) are linearly independent.  Hence there is a pair of row indices \(\mathscr I_a\) for which
\(B_a(P^\star)[\mathscr I_a,\{1,2\}]\) is nonsingular.  On the neighborhood where this determinant remains nonzero, define
\[
\begin{pmatrix}h_a(P)(1)\\h_a(P)(2)\end{pmatrix}
=B_a(P)[\mathscr I_a,\{1,2\}]^{-1}
\left{b_a(P)[\mathscr I_a]
-B_a(P)[\mathscr I_a,\{0\}]\frac a5\right}.
\]
This is a rational, hence smooth, function of the finite cell probabilities.  Outcome solvability and rank two imply that the remaining row equation also holds.  The gauge intersects each affine solution fiber once because every nonzero vector in \(\ker(B_a(P))\) has a nonzero zeroth coordinate after the neighborhood is shrunk.

For treatment bridges let \(r_0=0\) and \(r_1=2\), and impose
\[
q_a(P)(r_a)=q_a^{\mathrm{sel}}(r_a)=0.
\]
Since \(v(r_a)\ne0\), the two rows of \(B_a(P^\star)\) with indices different from \(r_a\) are linearly independent.  Choose two columns on which the resulting two-by-two minor is nonsingular and solve the corresponding two equations of
\(B_a(P)^\top q_a(P)=d(P)\) for the two unrestricted coordinates.  The resulting formula again consists only of matrix inversion and multiplication by finite moment coordinates, so it is rational and smooth while that minor remains nonzero.  Treatment solvability and rank two give the remaining equation.  These four sections agree with the displayed selected representatives at \(P^\star\) by uniqueness under their gauges.

For completeness, choose rationally varying null vectors \(u_a(P)\in\ker(B_a(P))\) and \(m_a(P)\in\ker(B_a(P)^\top)\) by fixing one coordinate that equals \(3\) at \(P^\star\).  Form the six normal generators
\[
D(P)=\bigl(F_0(h_0(P),m_0(P)),G_0(q_0(P),u_0(P)),R_0(m_0(P),u_0(P)),
F_1(h_1(P),m_1(P)),G_1(q_1(P),u_1(P)),R_1(m_1(P),u_1(P))\bigr)^\top.
\]
Because both kernels are one-dimensional, these six functions span \(\mathcal C_P\).  Let
\[
\Gamma(P)=E_P[D(P)D(P)^\top],qquad
c(P)=E_P[D(P)\phi_{\mathrm{sel}}(P)].
\]
The exact base-law Gram calculation in \(\mathrm{thm:finite-witness-nonattainment}\) is
\[
\Gamma(P^\star)=
\begin{pmatrix}
1323/500&-6/5&63/20&0&0&0\\
-6/5&343/5&-21&0&-21&0\\
63/20&-21&441/2&0&0&0\\
0&0&0&1323/500&6/5&-63/20\\
0&-21&0&6/5&343/5&-21\\
0&0&0&-63/20&-21&441/2
\end{pmatrix},
\qquad
\det\Gamma(P^\star)=\frac{16253767998768573}{12500000}>0.
\]
As a Gram matrix it is nonnegative definite, and its positive determinant makes it positive definite.  All entries of \(D(P)\), \(\Gamma(P)\), \(c(P)\), and the selected-score variance are continuous rational functions of the positive cell probabilities on the chosen chart.  Thus, after shrinking the chart, \(\Gamma(P)\) remains invertible.

The affine-normal identity used by \(\mathrm{thm:finite-witness-nonattainment}\) gives
\(\mathcal A_P=\phi_{\mathrm{sel}}(P)+\mathcal C_P\).  Hence finite-dimensional orthogonal projection onto the span of \(D(P)\) gives
\[
V_{\mathrm{eff}}(P)
=E_P[\phi_{\mathrm{sel}}(P)^2]-c(P)^\top\Gamma(P)^{-1}c(P),
\]
and therefore the continued selected-pair gap is
\[
\Delta(P):=E_P[\phi_{\mathrm{sel}}(P)^2]-V_{\mathrm{eff}}(P)
=c(P)^\top\Gamma(P)^{-1}c(P).
\]
This is continuous, and at the base law
\(\Delta(P^\star)=9496288/200416125>0\).  A final shrinking of the relatively open positive rank-two chart therefore ensures \(\Delta(P)>0\) throughout it.  This proves the stated smooth rational continuation and persistence of the strict variance gap.  The argument compares only the displayed bridge-score section with the quotient efficiency bound and makes no assertion about the probability limit of any published estimator.
\end{proof}
\par\smallskip\noindent \textit{Justification.} The narrowed theorem proves only the displayed selected score's exact strict gap and its gauge-continuation persistence. \textit{Gap.} Linking a published estimator to this particular selected pair requires a separate verified limit argument. \textit{Consumer.} The numerical benchmark is usable once an applied procedure's selected bridge limit is identified.

\begin{lemma}[lem:learned-projection-convergence]\label{lem:learned-projection-convergence}
For the nested construction/calibration/evaluation estimator, under the construction-sample bridge rates, multiplier and calibration envelopes, Gram cutoff, affine coefficient envelope, and corrected projection balance, the exact cell expansion implies
\[
 \max_k\|\widehat D_{J_n}^{(-k,A)}-D^0_{J_n}\|_{L^2(P;\ell^2)}=O_p(\delta_{D,n}),
 \qquad
 \delta_{D,n}=\left(\sum_{c=1}^{J_n}d_{c,n}^2\right)^{1/2}
 \le\sqrt{J_n}\max_{c\le J_n}d_{c,n}.
\]
The independent calibration sample satisfies the rates in \(\mathrm{lem:conditional-calibration-concentration}\).  If \(\Pi_{R,J_n}\) projects onto the positive population Gram eigenspace, then
\[
 \max_k\|\widehat\Pi_k^C-\Pi_{R,J_n}\|_{\mathrm{op}}=O_p(\eta_{\Pi,n}),
\]
and
\[
 \max_k\|\widehat b^{(-k,C\mid A)\top}\widehat D_{J_n}^{(-k,A)}
 -b_{J_n}^\top D^0_{J_n}\|_2
 =O_p\!\left[L_n\delta_{D,n}
 +\frac{\sqrt{\Lambda_{J_n}}}{\kappa_{J_n}}
 \left\{\zeta_{c,n}+(\zeta_{G,n}+\rho_n)L_n
 +\sqrt{\Lambda_{J_n}}B_{0,n}\eta_{\Pi,n}\right\}\right].
\]
Consequently, whenever the hypotheses of \(\mathrm{thm:quotient-efficiency}\) hold, the explicit coefficient-to-affine-weight identity gives
\[
 \max_k\|\widehat\psi^{(-k,C\mid A)}-\theta(P)-\phi_{\mathrm{eff}}\|_2
 =O_p(\Gamma_n)=o_p(1).
\]
\end{lemma}
\begin{proof}
Fix a fold \(k\) and suppress its construction and calibration superscripts.  Define
\[
 \Delta h_{aj}=\widehat h_{aj}^{(-k,A)}-h^0_{aj},\qquad
 \Delta q_{a\ell}=\widehat q_{a\ell}^{(-k,A)}-q^0_{a\ell}.
\]
Direct expansion of the signal in \(\mathrm{def:score-family}\) gives the exact cell identity
\[
 S_a(\widehat h_{aj},\widehat q_{a\ell})-S_a(h^0_{aj},q^0_{a\ell})
 =(1-I_aq^0_{a\ell})\Delta h_{aj}
 +I_a(Y-h^0_{aj})\Delta q_{a\ell}
 -I_a\Delta q_{a\ell}\Delta h_{aj}.                                      \tag{1}
\]
All products in (1) are square-integrable by \(\mathrm{ass:heldout-signal-integrability}\) and the fourth-moment hypotheses below.  By H\"older's inequality, \(\mathrm{ass:outcome-learning-rates}\), \(\mathrm{ass:treatment-learning-rates}\), and \(\mathrm{ass:moment-envelope-growth}\), uniformly over retained cells and folds,
\[
\begin{aligned}
 \|S_a(\widehat h_{aj},\widehat q_{a\ell})-S_a(h^0_{aj},q^0_{a\ell})\|_2
 &\le M_{q,n}\|\Delta h_{aj}\|_4+M_{y,n}\|\Delta q_{a\ell}\|_4
       +\|\Delta h_{aj}\|_4\|\Delta q_{a\ell}\|_4\\
 &=O_p(d^S_{a,j\ell,n}).                                                    \tag{2}
\end{aligned}
\]
Each pure or four-corner coordinate is a signed sum of at most four cell signals.  Thus the triangle inequality, followed by the definition in \(\mathrm{def:inference-gauges}\), yields
\[
 \|\widehat D-D^0\|_{L^2(P;\ell^2)}^2
 =\sum_{c=1}^{J_n}\|\widehat D_c-D_c^0\|_2^2
 =O_p\!\left(\sum_{c=1}^{J_n}d_{c,n}^2\right).
\]
Because \(K_{\mathrm{fold}}\) is fixed, the same bound holds for the maximum over folds.  Consequently,
\[
 \max_k\|\widehat D_{J_n}^{(-k,A)}-D^0_{J_n}\|_{L^2(P;\ell^2)}
 =O_p(\delta_{D,n}),\qquad
 \delta_{D,n}\le \sqrt{J_n}\max_{c\le J_n}d_{c,n}.                       \tag{3}
\]
Applying (2) to the two base cells gives, for
\[
 B^0=\sum_{a=0}^1t_aS_a(h^0_{a0},q^0_{a0})=\theta(P)+\phi_0,
\]
the further bound
\[
 \max_k\|\widehat B^{(-k,A)}-B^0\|_2=O_p(d_{0,n}).                         \tag{4}
\]

Put
\[
 G=E[D^0D^{0\top}],\quad c=E[D^0B^0],\quad
 \widetilde G=E[\widehat D\widehat D^\top\mid\mathcal D_{A,k}],\quad
 \widetilde c=E[\widehat D\widehat B\mid\mathcal D_{A,k}].
\]
The population envelope in \(\mathrm{ass:moment-envelope-growth}\) implies
\(\|D^0\|_{L^2(P;\ell^2)}\le\sqrt{J_n}B_{D,n}\) and \(\|B^0\|_2\le B_{0,n}\).  The Hilbert-space Cauchy--Schwarz inequality, (3), and (4) therefore give
\[
 \|\widetilde G-G\|_{\mathrm{op}}
 \le \|\widehat D-D^0\|_{L^2(P;\ell^2)}
       \{\|\widehat D\|_{L^2(P;\ell^2)}+\|D^0\|_{L^2(P;\ell^2)}\}
 =O_p\{2\sqrt{J_n}B_{D,n}\delta_{D,n}+\delta_{D,n}^2\},                  \tag{5}
\]
and
\[
 \|\widetilde c-c\|_2
 =O_p\!\left\{B_{0,n}\delta_{D,n}
 +(\sqrt{J_n}B_{D,n}+\delta_{D,n})d_{0,n}\right\}.                       \tag{6}
\]
The established \(\mathrm{lem:conditional-calibration-concentration}\) supplies calibration deviations of orders \(\eta_{G,n}\) and \(\eta_{c,n}\).  Combining that lemma with (5)--(6) and \(\mathrm{def:inference-gauges}\) proves
\[
 \max_k\|\widehat G_k^C-G\|_{\mathrm{op}}=O_p(\zeta_{G,n}),\qquad
 \max_k\|\widehat c_k^C-c\|_2=O_p(\zeta_{c,n}).                          \tag{7}
\]
In particular, the sharpened conditional rates used here are
\[
 \eta_{G,n}=B_{D,n}^2\left\{\sqrt{J_n/n_{C,\min}}+J_n/n_{C,\min}\right\},
 \qquad
 \eta_{c,n}=B_{D,n}B_{0,n}\sqrt{J_n/n_{C,\min}}.
\]
Thus, when the displayed envelopes are bounded, the calibration part is consistent under \(J_n/n_{C,\min}\to0\); no logarithmic dimension factor is required.

Let \(\Pi_R=\Pi_{R,J_n}\) project onto \(\operatorname{range}(G)\), and write
\(\delta_k=\|\widehat G_k^C-G\|_{\mathrm{op}}\).  By \(\mathrm{ass:gram-conditioning}\), every positive eigenvalue of \(G\) is at least \(\kappa_{J_n}\), \(0<\tau_n<\kappa_{J_n}/2\), and \(\delta_k=o_p(\tau_n)\) uniformly over the fixed number of folds.  The variational inequality
\[
 |u^\top(\widehat G_k^C-G)u|\le\delta_k\|u\|_2^2                         \tag{8}
\]
shows that every empirical eigenvalue paired with the population zero spectrum is at most \(\delta_k<\tau_n\), whereas the remaining eigenvalues are at least \(\kappa_{J_n}-\delta_k>\tau_n\), with probability tending to one.  Hence the thresholded operator
\[
 \widehat G_{k,\tau}=\widehat\Pi_k^C\widehat G_k^C\widehat\Pi_k^C
\]
has the same rank as \(G\), reduced minimum modulus at least \(\kappa_{J_n}/2\), and
\(\|\widehat G_{k,\tau}-G\|_{\mathrm{op}}\le2\delta_k\).  Applying the established \(\mathrm{lem:operator-kernel-perturbation}\) to these finite-dimensional self-adjoint operators and taking complements of their kernel projections yields
\[
 \max_k\|\widehat\Pi_k^C-\Pi_R\|_{\mathrm{op}}
 =O_p(\zeta_{G,n}/\kappa_{J_n})=O_p(\eta_{\Pi,n}).                         \tag{9}
\]

The population least-squares coefficient satisfies
\[
 Gb_{J_n}=c,\qquad b_{J_n}\in\operatorname{range}(G).                     \tag{10}
\]
Indeed, if \(v\in\ker G\), then \(E[(v^\top D^0)^2]=v^\top Gv=0\), so \(v^\top D^0=0\) almost surely and \(v^\top c=E[v^\top D^0B^0]=0\).  Hence \(c\in\operatorname{range}(G)\), which proves (10) on the nonzero Gram eigenspace specified in \(\mathrm{def:population-affine-sieve}\).  From \(\mathrm{def:affine-ensemble-estimator}\), the cutoff-ridge coefficient obeys
\[
 (\widehat G_k^C+\rho_n I)\widehat b=\widehat\Pi_k^C\widehat c_k^C,
 \qquad \widehat b=\widehat\Pi_k^C\widehat b.                            \tag{11}
\]
Subtracting (10) from (11), using \(G(I-\Pi_R)=0\), and projecting onto \(\operatorname{range}(G)\) gives
\[
 G\Pi_R(\widehat b-b_{J_n})
 =\Pi_R\{(G-\widehat G_k^C)\widehat b+(\widehat c_k^C-c)
 -\rho_n\widehat b-(I-\widehat\Pi_k^C)\widehat c_k^C\}.                 \tag{12}
\]
Since \(c=\Pi_Rc\), (7) and (9) imply
\[
 \|(I-\widehat\Pi_k^C)\widehat c_k^C\|_2
 \le \|\widehat c_k^C-c\|_2+\|(\Pi_R-\widehat\Pi_k^C)c\|_2
 =O_p\{\zeta_{c,n}+\eta_{\Pi,n}\|c\|_2\}.                             \tag{13}
\]
For each unit vector \(u\), Cauchy--Schwarz gives
\(
 |u^\top c|\le\{E[(u^\top D^0)^2]\}^{1/2}\|B^0\|_2
\), and therefore
\[
 \|c\|_2\le\sqrt{\Lambda_{J_n}}B_{0,n}.                                \tag{14}
\]
Invert \(G\) only on its positive eigenspace.  Equations (7), (12)--(14), the lower eigenvalue bound in \(\mathrm{ass:gram-conditioning}\), and the coefficient envelope in \(\mathrm{ass:affine-weight-envelope}\) prove
\[
 \max_k\|\Pi_R(\widehat b-b_{J_n})\|_2
 =O_p\!\left[
 \frac{\zeta_{c,n}+(\zeta_{G,n}+\rho_n)L_n
 +\sqrt{\Lambda_{J_n}}B_{0,n}\eta_{\Pi,n}}
 {\kappa_{J_n}}
 \right].                                                                  \tag{15}
\]
This retains the learned coefficient multiplying the Gram perturbation and the ridge term.  Furthermore,
\[
 \|v^\top D^0\|_2=\|G^{1/2}v\|_2
 \le\sqrt{\Lambda_{J_n}}\|\Pi_Rv\|_2,
 \qquad
 \|\widehat b^\top(\widehat D-D^0)\|_2\le L_n\delta_{D,n}.
\]
Together with (3) and (15), these inequalities give
\[
\begin{aligned}
 &\max_k\|\widehat b^{(-k,C\mid A)\top}\widehat D_{J_n}^{(-k,A)}
       -b_{J_n}^\top D^0_{J_n}\|_2\\
 &\quad=O_p\!\left[L_n\delta_{D,n}
 +\frac{\sqrt{\Lambda_{J_n}}}{\kappa_{J_n}}
 \left\{\zeta_{c,n}+(\zeta_{G,n}+\rho_n)L_n
 +\sqrt{\Lambda_{J_n}}B_{0,n}\eta_{\Pi,n}\right\}\right].             \tag{16}
\end{aligned}
\]

It remains to identify the adjusted signal.  Collecting the coefficient of every bridge cell in the deterministic coefficient-to-weight map in \(\mathrm{def:affine-ensemble-estimator}\) gives, for each arm,
\[
 \sum_{j=0}^{J_n^h}\sum_{\ell=0}^{J_n^q}\widehat\lambda_{a,j\ell}=1,
 \qquad
 \widehat\psi=\widehat B-\widehat b^\top\widehat D.                     \tag{17}
\]
The identical algebra applies to the population coefficient.  By the established \(\mathrm{lem:normal-affine-identity}\), every coordinate of \(D^0\) belongs to \(\mathcal C_P\).  Orthogonal decomposition therefore gives, for every \(b\in\mathbb R^{J_n}\),
\[
 \|\phi_0-b^\top D^0\|_2^2
 =\|\phi_0-\Pi_{\mathcal C_P}\phi_0\|_2^2
  +\|\Pi_{\mathcal C_P}\phi_0-b^\top D^0\|_2^2.                         \tag{18}
\]
The normal equation (10) shows that \(b_{J_n}\) minimizes the \(b\)-dependent term in (18), whose norm is therefore \(\alpha_{J_n}\).  Under the hypotheses of the established \(\mathrm{thm:quotient-efficiency}\),
\(\phi_{\mathrm{eff}}=\phi_0-\Pi_{\mathcal C_P}\phi_0\).  Since \(B^0=\theta(P)+\phi_0\), it follows that
\[
 \|B^0-b_{J_n}^\top D^0-\theta(P)-\phi_{\mathrm{eff}}\|_2
 =\alpha_{J_n}.                                                             \tag{19}
\]
Combining (4), (16), (17), and (19) gives exactly
\[
 \max_k\|\widehat\psi^{(-k,C\mid A)}-\theta(P)-\phi_{\mathrm{eff}}\|_2
 =O_p(\Gamma_n).
\]
Finally, \(\mathrm{ass:projection-rate}\) states \(\Gamma_n=o_p(1)\); its constituent approximation, conditioning, ridge, and learning requirements are supplied by the other displayed hypotheses.  This proves every assertion of the lemma.
\end{proof}
\par\smallskip\noindent \textit{Justification.} The exact cell expansion gives the dimension-aware dictionary error, independent calibration gives Gram and cross-moment concentration, spectral cutoff controls the population null space, and the resolvent retains every coefficient and ridge multiplier. The deterministic four-corner map then converts the projected signal into a valid affine bridge ensemble. \textit{Gap.} No proof gap remains under the displayed sieve approximation, multiplier, calibration, spectral, coefficient-envelope, and projection-gauge conditions; these conditions may be restrictive when the positive Gram spectrum decays quickly. \textit{Consumer.} The lemma supplies, rather than assumes, foldwise \(L^2\) convergence to the efficient signal required by the inference theorem.

\begin{theorem}[thm:affine-ensemble-inference]\label{thm:affine-ensemble-inference}
Assume the hypotheses of \(\mathrm{thm:quotient-efficiency}\), \(0<V_{\mathrm{eff}}(P)<\infty\), fixed \(K_{\mathrm{fold}}\), the nested-split conditions of \(\mathrm{lem:learned-projection-convergence}\), and
\[
 \sqrt nL_n\min\{r_{h,n}e_{q,2,n},r_{q,n}e_{h,2,n}\}=o(1).
\]
Then, pointwise under \(P\),
\[
 \sqrt n\{\widehat\theta_n-\theta(P)\}
 =n^{-1/2}\sum_{i=1}^n\phi_{\mathrm{eff}}(O_i)+o_p(1)
 \ \Rightarrow\ N\{0,V_{\mathrm{eff}}(P)\},
\]
\(\widehat V_n/V_{\mathrm{eff}}(P)\to_p1\), and
\[
 P\{\theta(P)\in\mathrm{CI}_{1-\alpha}\}\to1-\alpha.
\]
For a locally uniform regularity extension, fix \(T<\infty\) and a class \(\mathcal S\) of bounded DQM path scores, and write \(P_{n,t,s}=P_{t/\sqrt n,s}\).  Suppose uniformly over \(s\in\mathcal S\) and \(|t|\le T\): the hypotheses of \(\mathrm{thm:quotient-efficiency}\), including tangent equality and pathwise product validity, hold at every \(P_{n,t,s}\); the construction-sample bridge and multiplier bounds defining \(d_{c,n}\) and \(d_{0,n}\), the sieve approximation \(\alpha_{J_n}\), the calibration exponential-moment and concentration bounds \(\eta_{G,n},\eta_{c,n}\), the positive-eigenvalue gap, cutoff stability, coefficient and weight envelopes, and \(\Gamma_n=o_p(1)\) all hold under \(P_{n,t,s}\); the root-\(n\) mixed-bias gauge is \(o_p(1)\); \(\phi_{\mathrm{eff}}(P_{n,t,s})\to\phi_{\mathrm{eff}}(P)\) after transport to the fixed base space \(L^2(P)\) by the canonical likelihood-ratio isometry, and its transported squared family is uniformly integrable; and
\[
 \sup_{s,t}\left|\sqrt n\{\theta(P_{n,t,s})-\theta(P)\}-tE_P\{\phi_{\mathrm{eff}}(P)s\}\right|\to0,
\]
\[
 \sup_{s,t}\left|\log\frac{dP_{n,t,s}^n}{dP^n}
 -\frac{t}{\sqrt n}\sum_{i=1}^ns(O_i)+\frac{t^2}{2}E_P[s^2]\right|=o_P(1),
\]
with a uniform Lindeberg condition for \(\phi_{\mathrm{eff}}(P_{n,t,s})\), and suppose the arrays \(P_{n,t,s}^n\) and \(P^n\) are mutually contiguous uniformly on the declared class.  Then the preceding asymptotic linear expansion, variance consistency, and centered Gaussian limit hold uniformly over \(s\in\mathcal S\) and \(|t|\le T\).  In particular, the estimator is regular and locally efficient on every such uniformly controlled contiguous path class.
\end{theorem}
\begin{proof}
We first prove the pointwise assertions.  Let \(P_{n,k}\) denote the empirical measure on \(\mathcal I_k\), and define the random function
\[
 f_k=\widehat\psi^{(-k,C\mid A)}-\theta(P)-\phi_{\mathrm{eff}}.
\]
The proved \(\mathrm{lem:learned-projection-convergence}\), with its nested-split hypotheses, gives
\[
 \max_{k\le K_{\mathrm{fold}}}\|f_k\|_{L^2(P)}=o_p(1).                   \tag{20}
\]
By \(\mathrm{ass:training-measurability}\), \(f_k\) is measurable with respect to \(\mathcal D_{A,k}\vee\mathcal D_{C,k}\), while the observations in \(\mathcal I_k\) are independent of that sigma-field under \(\mathrm{ass:iid-sampling}\).  Conditional square-integrability follows from \(\mathrm{ass:heldout-signal-integrability}\).  Since all fold sizes are proportional to \(n\) and \(K_{\mathrm{fold}}\) is fixed, conditional variance calculation yields
\[
 E\!\left[
 \left.\left\{\frac{|\mathcal I_k|}{\sqrt n}(P_{n,k}-P)f_k\right\}^2
 \right|\mathcal D_{A,k}\vee\mathcal D_{C,k}\right]
 \le C\|f_k\|_2^2=o_p(1).                                                  \tag{21}
\]
Conditional Chebyshev inequality and a finite union bound imply
\[
 \sum_{k=1}^{K_{\mathrm{fold}}}\frac{|\mathcal I_k|}{\sqrt n}
 (P_{n,k}-P)f_k=o_p(1).                                                      \tag{22}
\]

We next bound the conditional drift.  For any learned cell, (1), outcome-bridge solvability, and treatment-bridge solvability give
\[
 P\{S_a(\widehat h,\widehat q)-S_a(h^0,q^0)\}
 =-E[I_a(\widehat q-q^0)(\widehat h-h^0)].                                 \tag{23}
\]
Indeed, treating the learned functions as fixed after conditioning on their training sigma-field,
\[
 E[(\widehat h-h^0)(1-I_aq^0)]
 =E[(\widehat h-h^0)(1-T_a^*q^0)]=0
\]
because \(q^0\in\mathcal Q_a(P)\), and
\[
 E[I_a(\widehat q-q^0)(Y-h^0)]
 =E[I_a(\widehat q-q^0)(g_a-T_ah^0)]=0
\]
because \(h^0\in\mathcal H_a(P)\).  Conditioning the last expectation in (23), first on \((A,Z,X)\) and then on \((W,X)\), and applying Cauchy--Schwarz in the corresponding Hilbert spaces gives both bounds
\[
 |E[I_a\Delta q\Delta h]|
 \le\|T_a\Delta h\|_{\mathcal G_a}\|\Delta q\|_{\mathcal G_a},
 \qquad
 |E[I_a\Delta q\Delta h]|
 \le\|\Delta h\|_{\mathcal H}\|T_a^*\Delta q\|_{\mathcal H}.         \tag{24}
\]
The affine weights sum to one within each arm by (17), and their total variation is at most \(L_n\) by \(\mathrm{ass:affine-weight-envelope}\).  Hence (23)--(24), \(\mathrm{ass:outcome-residual-rate}\), \(\mathrm{ass:adjoint-residual-rate}\), and the \(L^2\) parts of \(\mathrm{ass:outcome-learning-rates}\) and \(\mathrm{ass:treatment-learning-rates}\) prove
\[
 \max_k|P\widehat\psi^{(-k,C\mid A)}-\theta(P)|
 =O_p\!\left(L_n\min\{r_{h,n}e_{q,2,n},r_{q,n}e_{h,2,n}\}\right)
 =o_p(n^{-1/2}),                                                             \tag{25}
\]
where the last equality is precisely \(\mathrm{ass:primitive-mixed-bias-rate}\).

The folds partition the sample, \(P\phi_{\mathrm{eff}}=0\) by \(\mathrm{thm:quotient-efficiency}\), and therefore (22)--(25) give
\[
\begin{aligned}
 \sqrt n\{\widehat\theta_n-\theta(P)\}
 &=\frac1{\sqrt n}\sum_{i=1}^n\phi_{\mathrm{eff}}(O_i)
 +\sum_k\frac{|\mathcal I_k|}{\sqrt n}(P_{n,k}-P)f_k
 +\sum_k\frac{|\mathcal I_k|}{\sqrt n}Pf_k\\
 &=\frac1{\sqrt n}\sum_{i=1}^n\phi_{\mathrm{eff}}(O_i)+o_p(1).          \tag{26}
\end{aligned}
\]
Also by \(\mathrm{thm:quotient-efficiency}\),
\(E[\phi_{\mathrm{eff}}^2]=V_{\mathrm{eff}}(P)\in(0,\infty)\).  For every \(\epsilon>0\), integrability and dominated convergence give
\[
 E[\phi_{\mathrm{eff}}^2\mathbf 1\{|\phi_{\mathrm{eff}}|>\epsilon\sqrt n\}]
 \longrightarrow0.
\]
Thus \(\mathrm{lem:lindeberg-feller-central-limit}\), applied to
\(X_{ni}=n^{-1/2}\phi_{\mathrm{eff}}(O_i)\), implies
\[
 n^{-1/2}\sum_{i=1}^n\phi_{\mathrm{eff}}(O_i)
 \Rightarrow N\{0,V_{\mathrm{eff}}(P)\}.                                  \tag{27}
\]
For completeness, (26) and (27) have the same limit without invoking any additional result: for every continuity point \(x\) of the normal distribution and every \(\delta>0\), the distribution function of the left side of (26) is bounded above and below by those of the leading sum at \(x+\delta\) and \(x-\delta\), plus the probability that the remainder exceeds \(\delta\); first let \(n\to\infty\), then \(\delta\downarrow0\).  This proves the asserted pointwise Gaussian limit.

For variance estimation, (20), conditional Markov inequality, and the same foldwise conditioning imply
\[
 \frac1n\sum_{i=1}^nf_{k(i)}(O_i)^2=o_p(1).                                \tag{28}
\]
The weak law for \(\phi_{\mathrm{eff}}^2\) follows under only its finite first moment: for fixed \(M\), Chebyshev's inequality applies to the bounded variables \(\phi_{\mathrm{eff}}^2\wedge M\), and Markov's inequality controls the empirical remainder by
\(E[\phi_{\mathrm{eff}}^2\mathbf 1\{\phi_{\mathrm{eff}}^2>M\}]\), which vanishes as \(M\to\infty\).  Hence
\[
 \frac1n\sum_{i=1}^n\phi_{\mathrm{eff}}(O_i)^2
 \longrightarrow_pV_{\mathrm{eff}}(P).                                   \tag{29}
\]
Empirical Cauchy--Schwarz, (28), and (29) make the cross term negligible.  Moreover, (26)--(27) imply \(\widehat\theta_n-\theta(P)=O_p(n^{-1/2})\).  Using the definition of \(\widehat V_n\), we obtain
\[
\begin{aligned}
 \widehat V_n
 &=\frac1n\sum_{i=1}^n\{\widehat\psi_i-\theta(P)\}^2
   -\{\widehat\theta_n-\theta(P)\}^2\\
 &\longrightarrow_pV_{\mathrm{eff}}(P).                                   \tag{30}
\end{aligned}
\]
Since the limit is positive, division in (30) proves
\(\widehat V_n/V_{\mathrm{eff}}(P)\to_p1\).  Combining (26)--(27) with this ratio convergence, using the same distribution-function sandwich as above, shows that the studentized statistic converges to \(N(0,1)\).  Continuity of its distribution at \(\pm z_{1-\alpha/2}\) proves
\[
 P\{\theta(P)\in\mathrm{CI}_{1-\alpha}\}\longrightarrow1-\alpha.
\]

We now prove the locally uniform extension, without resolving \(\mathrm{oeq:score-lifting-completeness}\).  Let
\[
 \Xi_n=\{(s,t):s\in\mathcal S,\ |t|\le T\},\qquad
 Q_{n,\xi}=P_{n,t,s}\quad\text{for }\xi=(s,t),
\]
and abbreviate
\[
 \theta_{n,\xi}=\theta(Q_{n,\xi}),\qquad
 \phi_{n,\xi}=\phi_{\mathrm{eff}}(Q_{n,\xi}),\qquad
 V_{n,\xi}=E_{Q_{n,\xi}}[\phi_{n,\xi}^2].
\]
The theorem's uniform hypotheses expressly impose the hypotheses of \(\mathrm{thm:quotient-efficiency}\), including tangent equality and pathwise product validity, at every \(Q_{n,\xi}\).  Thus that theorem identifies \(\phi_{n,\xi}\) as the canonical gradient at each such law.  This is a premise at the local laws, not a proof of the unresolved score-lifting equality.

Choose an arbitrary deterministic sequence \(\xi_n\in\Xi_n\) and put \(Q_n=Q_{n,\xi_n}\).  Equations (1)--(19) apply under \(Q_n\): every bridge-learning, multiplier, sieve-approximation, calibration-concentration, positive-eigenvalue, cutoff, coefficient, weight, and projection-gauge premise used there is assumed uniformly over \(\Xi_n\).  Therefore
\[
 \max_k\|\widehat\psi_{n,k}-\theta(Q_n)-\phi_{n,\xi_n}\|_{L^2(Q_n)}
 =o_{Q_n}(1).                                                               \tag{31}
\]
The uniformly root-\(n\)-small mixed-bias gauge likewise makes the analogue of (25) \(o_{Q_n}(n^{-1/2})\).  Conditional evaluation-fold Chebyshev under \(Q_n\), exactly as in (21)--(26), then proves
\[
 \sqrt n\{\widehat\theta_n-\theta(Q_n)\}
 =\frac1{\sqrt n}\sum_{i=1}^n\phi_{n,\xi_n}(O_i)+o_{Q_n}(1).              \tag{32}
\]
Because \(\xi_n\) was arbitrary, (32) holds uniformly over \(\Xi_n\).  Indeed, failure of uniformity would provide \(\epsilon,\delta>0\), a subsequence, and choices \(\xi_n\) on that subsequence for which the remainder probability is at least \(\delta\), contradicting the sequencewise argument.

To compare the varying gradients in one Hilbert space, define the canonical likelihood-ratio isometry
\[
 \mathsf J_{n,\xi}f=f\sqrt{dQ_{n,\xi}/dP}:L^2(Q_{n,\xi})\longrightarrow L^2(P).
\]
It is an isometry because
\(\|\mathsf J_{n,\xi}f\|_{L^2(P)}^2=E_{Q_{n,\xi}}[f^2]\).  The assumed uniform transported \(L^2(P)\) convergence gives
\[
 \sup_{\xi\in\Xi_n}|V_{n,\xi}-V_{\mathrm{eff}}(P)|\longrightarrow0.       \tag{33}
\]
Every \(Q_{n,\xi}\) lies in \(\mathcal V_P\), so its likelihood ratio is between \(1-\varepsilon\) and \(1+\varepsilon\).  Hence, for every fixed \(\epsilon>0\),
\[
\begin{aligned}
 &E_{Q_{n,\xi}}[\phi_{n,\xi}^2
       \mathbf 1\{|\phi_{n,\xi}|>\epsilon\sqrt n\}]\\
 &\quad\le E_P[(\mathsf J_{n,\xi}\phi_{n,\xi})^2
       \mathbf 1\{|\mathsf J_{n,\xi}\phi_{n,\xi}|
       >\epsilon\sqrt{1-\varepsilon}\sqrt n\}],
                                                                                \tag{34}
\end{aligned}
\]
which tends to zero uniformly by uniform integrability of the transported squared family; this also verifies the theorem's displayed uniform Lindeberg requirement.  For every sequence \(\xi_n\), \(\mathrm{lem:lindeberg-feller-central-limit}\), (33), and (34) give, under \(Q_{n,\xi_n}^n\),
\[
 n^{-1/2}\sum_{i=1}^n\phi_{n,\xi_n}(O_i)
 \Rightarrow N\{0,V_{\mathrm{eff}}(P)\}.                                  \tag{35}
\]
The preceding contradiction argument upgrades (35) to uniform convergence in bounded-Lipschitz distance over \(\Xi_n\).  Combining (32) and (35) proves the asserted uniform centered Gaussian limit.

The variance argument is also uniform.  Along every sequence \(\xi_n\), conditional Markov inequality applied to (31) gives the analogue of (28).  For the leading squares, truncate \(\phi_{n,\xi_n}^2\) at \(M\).  Chebyshev's inequality makes the centered average of the truncated variables vanish because its variance is at most
\(M E_{Q_n}[\phi_{n,\xi_n}^2]/n\), which is uniformly \(O(M/n)\) by (33).  Uniform integrability in (34) controls both the empirical and population tails as \(M\to\infty\).  The arbitrary-sequence argument thus yields, for every \(\epsilon>0\),
\[
 \sup_{\xi\in\Xi_n}Q_{n,\xi}^n
 \left(\left|\widehat V_n-V_{n,\xi}\right|>\epsilon\right)
 \longrightarrow0.                                                         \tag{36}
\]
Equations (33) and (36), together with \(V_{\mathrm{eff}}(P)>0\), prove uniform ratio consistency.  The distribution-function sandwich used pointwise, now applied along every sequence and combined with the same contradiction argument, proves uniform Wald coverage for \(\theta(Q_{n,\xi})\).

Finally, the assumed uniform likelihood expansion and mutual contiguity certify that \(Q_{n,\xi}^n\) are the declared contiguous local experiments and make base-law and local-law negligible events interchangeable.  The assumed uniform target expansion states
\[
 \sqrt n\{\theta(Q_{n,\xi})-\theta(P)\}
 =tE_P[\phi_{\mathrm{eff}}(P)s]+o(1)                                      \tag{37}
\]
uniformly over \(\Xi_n\).  Centering at \(\theta(P)\) therefore adds the shift in (37), whereas centering at the true local target removes it and leaves the common centered limit in (35).  This is regularity on every declared uniformly controlled bounded-score contiguous class.  At every local law \(\mathrm{thm:quotient-efficiency}\) makes \(\phi_{n,\xi}\) canonical, and (33) makes its variance converge to \(V_{\mathrm{eff}}(P)\).  Thus the estimator attains the local semiparametric efficiency bound, completing the proof.
\end{proof}
\par\smallskip\noindent \textit{Justification.} Nested cross-fitting makes the evaluation-fold empirical remainder conditionally negligible, the exact primal-dual identity and root-\(n\) product condition control drift, and learned-score convergence yields asymptotic linearity. Truncation proves variance consistency under second moments, while the fixed-space transport and subsequence principle prove the locally uniform triangular-array extension. \textit{Gap.} The theorem remains conditional on \(\mathrm{thm:quotient-efficiency}\), hence on tangent equality and pathwise product validity at each relevant law. It does not settle \(\mathrm{oeq:score-lifting-completeness}\), and its uniform conclusion is limited to the explicitly controlled bounded-score contiguous classes. \textit{Consumer.} When those quotient and learning conditions are verified, the affine ensemble provides pointwise and locally uniform efficient standard errors and Wald coverage for the proximal bridge functional.

\begin{theorem}[thm:unique-bijective-reduction]\label{thm:unique-bijective-reduction}
Suppose \(T_a:\mathcal H\to\mathcal G_a\) and \(T_a^*:\mathcal G_a\to\mathcal H\) are bijective for \(a=0,1\). Then \(\mathcal N_a(P)=\mathcal K_a(P)=\{0\}\), each bridge fiber is a singleton, \(\mathcal C_P=\{0\}\), and \(\Phi_P\) is either empty or the singleton containing the unique bridge score; it is the latter whenever that score is square-integrable. Whenever the hypotheses and conclusion of \(\mathrm{thm:quotient-efficiency}\) apply, this unique score is \(\phi_{\mathrm{eff}}\). Unconditionally, at a strictly positive finite-support law for which the square arm matrices are nonsingular, the full-rank joint-solvability germ is relatively open in the probability simplex, \(\mathcal T_P=L_0^2(P)\), and the unique square-integrable bridge score is the canonical gradient.
\end{theorem}
\begin{proof}
For each arm, bijectivity of \(T_a\) implies injectivity, and bijectivity of \(T_a^*\) does the same for the adjoint.  Hence
\[
 \mathcal N_a(P)=\ker T_a=\{0\},
 \qquad
 \mathcal K_a(P)=\ker T_a^*=\{0\}.                                \tag{15}
\]
Surjectivity also supplies existence: \(g_a\in\mathcal G_a\) has a preimage under \(T_a\), while the constant function \(1\in\mathcal H\) has a preimage under \(T_a^*\).  If \(h_a,h'_a\in\mathcal H_a(P)\), then \(T_a(h_a-h'_a)=0\), so (15) gives \(h_a=h'_a\).  Likewise, if \(q_a,q'_a\in\mathcal Q_a(P)\), then \(T_a^*(q_a-q'_a)=0\), whence \(q_a=q'_a\).  Thus each bridge fiber is a singleton.  Every generator in \(\mathrm{def:normal-space}\) contains either \(m\in\mathcal K_a(P)\) or \(u\in\mathcal N_a(P)\), so (15) makes every generator zero and
\[
 \mathcal C_P=\{0\}.                                                \tag{16}
\]
The unique bridge representatives yield at most one candidate in \(\Phi_P\); they yield exactly one when their centered bridge score is square-integrable.  Whenever the hypotheses of \(\mathrm{thm:quotient-efficiency}\) hold, (16) gives
\[
 \phi_{\mathrm{eff}}
 =\phi_0-\Pi_{\{0\}}\phi_0=\phi_0,
\]
which proves the conditional general-model assertion.

It remains to prove the unconditional finite-support assertion.  Let the finitely many \((Z,X)\)-cells index rows and the finitely many \((W,X)\)-cells index columns.  For arm \(a\), define the joint-moment matrix, outcome vector, and marginal vector by
\[
 (B_a)_{ij}=P\{A=a,(Z,X)=i,(W,X)=j\},\quad
 (b_a)_i=E_P[I_a\mathbf 1\{(Z,X)=i\}Y],\quad
 d_j=P\{(W,X)=j\}.                                                  \tag{17}
\]
In these coordinates the two bridge equations and the arm target are
\[
 B_ah_a=b_a,
 \qquad B_a^\top q_a=d,
 \qquad \theta_a(P)=d^\top h_a.                                   \tag{18}
\]
The equivalence follows by multiplying the conditional equations defining \(T_a h_a=g_a\) and \(T_a^*q_a=1\) by their positive cell probabilities.  By hypothesis, each square matrix \(B_a\) is nonsingular.  Since determinant and cell probabilities are continuous functions of the probability vector, all cells remain positive and both matrices remain nonsingular on a sufficiently small relative neighborhood in the probability simplex.  On that neighborhood,
\[
 h_a(P')=B_a(P')^{-1}b_a(P'),
 \qquad q_a(P')=B_a(P')^{-\top}d(P')                               \tag{19}
\]
exist uniquely and vary continuously.  In the fixed Hilbert coordinates of \(\mathrm{def:local-operator-chart}\), the matrix of \(\widetilde T_a(P')\) is obtained from \(B_a(P')\) by left and right multiplication by positive diagonal matrices formed from the arm and marginal cell probabilities.  These diagonal matrices vary continuously while all cells are positive.  Thus the least singular value of \(\widetilde T_a(P')\) is continuous and positive at \(P\); after shrinking the neighborhood it exceeds one half its base value.  The kernel and adjoint-kernel dimensions remain zero.  Therefore the fixed-defect and reduced-modulus requirements of \(\mathrm{def:local-operator-chart}\), as well as both solvability requirements in \(\mathrm{def:joint-model}\), hold throughout this neighborhood.  The second-moment condition is automatic on a finite sample space.  This proves that the full-rank joint-solvability germ is relatively open.

Let \(s\in L_0^2(P)\).  Finite support makes \(s\) bounded.  Therefore, for all sufficiently small \(|t|\),
\[
 p_t(o)=p(o)\{1+ts(o)\}                                            \tag{20}
\]
is strictly positive, sums to one, lies in the relative open germ just obtained and in the likelihood-ratio ball, and is DQM with score \(s\), since
\[
 \sqrt{p_t(o)}=\sqrt{p(o)}\{1+ts(o)/2+O(t^2)\}
\]
uniformly over the finite support.  Hence every element of \(L_0^2(P)\) is an actual path score and
\[
 \mathcal T_P=L_0^2(P).                                             \tag{21}
\]

For completeness, differentiate the finite equations and thereby verify the gradient directly, without invoking score lifting.  Along any DQM path in this germ, finite-dimensional DQM gives
\(\dot p(o)=p(o)s(o)\), and hence derivatives \(\dot B_a\), \(\dot b_a\), and \(\dot d\) in (17).  Nonsingularity and (19) give differentiable bridge representatives.  Differentiating (18), and using \(B_a^\top q_a=d\), yields
\[
 \begin{aligned}
 \dot\theta_a
 &=\dot d^\top h_a+d^\top\dot h_a
 =\dot d^\top h_a+q_a^\top B_a\dot h_a\\
 &=\dot d^\top h_a+q_a^\top(\dot b_a-\dot B_ah_a)
 =E_P[S_a(h_a,q_a)s].                                               \tag{22}
 \end{aligned}
\]
The last equality follows term by term from (17): its three terms are respectively \(E_P(h_as)\), \(E_P(I_aq_aYs)\), and \(-E_P(I_aq_ah_as)\).  Taking the signed arm contrast in (22), and using \(E_Ps=0\), proves
\[
 \dot\theta_P(s)=E_P(\phi_{h,q}s).
\]
All variables are bounded on the finite sample space, so this unique bridge score is square-integrable.  By (21), a gradient orthogonal to the full tangent must vanish; consequently the displayed bridge score is the unique gradient in \(L_0^2(P)\), and hence the canonical gradient.
\end{proof}
\par\smallskip\noindent \textit{Justification.} The repaired theorem separates algebraic uniqueness, conditional general efficiency, and unconditional finite full-rank efficiency. \textit{Gap.} Bijectivity does not itself create paths or guarantee an unbounded product score is square-integrable. \textit{Consumer.} Finite full-rank unique-bridge software retains the familiar canonical score endpoint.

\begin{theorem}[thm:score-lifting-fails-without-chart-richness]\label{thm:score-lifting-fails-without-chart-richness}
If the full saturated chart is replaced by an arbitrary analyst-selected subchart \(\mathcal S\) containing \(P\), write \(\mathcal T_{P,\mathcal S}\) for the closed span of scores of paths restricted to \(\mathcal S\). Pointwise bridge existence, defect indices, and reduced-modulus bounds do not imply reverse score lifting. In particular, for \(\mathcal S=\{P\}\) the actual tangent space is \(\{0\}\); whenever \(\mathcal C_P^\perp\ne\{0\}\), one has \(\mathcal T_{P,\mathcal S}\ne\mathcal C_P^\perp\).
\end{theorem}
\begin{proof}
Every path with range in the singleton \(\mathcal S=\{P\}\) is constant.  Its square-root density derivative is therefore zero, so its score is zero and the closed span of all such scores is \(\mathcal T_{P,\mathcal S}=\{0\}\).  The operator, its two kernels, its reduced minimum modulus, and both bridge fibers at the only law in \(\mathcal S\) are exactly their base-law objects; hence all pointwise operator and solvability restrictions can hold.  If \(\mathcal C_P^\perp\ne\{0\}\), choose \(s\ne0\) in that space.  Then \(s\notin\mathcal T_{P,\mathcal S}\).  This explicit witness proves that reverse lifting is a property of the admitted path class, not a consequence of pointwise operator restrictions.
\end{proof}
\par\smallskip\noindent \textit{Justification.} This explicit counterexample isolates why a canonical saturated domain is necessary before reverse lifting can even be posed. \textit{Gap.} Pointwise rank and bridge restrictions do not determine which statistical paths an analyst-selected subchart admits. \textit{Consumer.} The diagnostic prevents efficiency conclusions from being inferred from path-poor charts.

\begin{lemma}[lem:finite-full-rank-constraint-tangent]\label{lem:finite-full-rank-constraint-tangent}
Let \(P\) be a strictly positive law on a finite sample space, and let a neighborhood of \(P\) in the probability simplex be cut out by a continuously differentiable map \(\Xi\) whose derivative has full row rank at \(P\). Then every vector in \(\ker D\Xi(P)\) is the derivative of a continuously differentiable positive-probability path in the level set \(\Xi^{-1}\{\Xi(P)\}\). Consequently the closed score tangent is the kernel of the corresponding \(L^2(P)\) normal covectors.
\end{lemma}
\begin{proof}
Use the first \(N-1\) cell probabilities as coordinates on the open simplex, where \(N\) is the number of cells.  Full row rank permits a permutation \(x=(u,v)\) for which \(L=D_v\Xi(P)\) is square and nonsingular.  On a sufficiently small closed coordinate ball, continuity of \(D_v\Xi\) gives
\[
\sup_{(u,v)}\|I-L^{-1}D_v\Xi(u,v)\|_{\mathrm{op}}\le\frac12.
\]
For fixed nearby \(u\), the iteration
\[
v_{m+1}=v_m-L^{-1}\{\Xi(u,v_m)-\Xi(P)\}
\]
contracts distances by at most \(1/2\) and, after shrinking the ball, maps the ball into itself.  Its successive differences are bounded by a geometric series, hence it converges to a unique \(v(u)\) satisfying \(\Xi(u,v(u))=\Xi(P)\).  Subtracting the fixed-point equations, dividing by an increment, and using continuity of the derivative gives
\[
Dv(P)[\dot u]=-L^{-1}D_u\Xi(P)[\dot u].
\]
Thus, for any \((\dot u,\dot v)\in\ker D\Xi(P)\), the curve \(u_t=u_0+t\dot u\), \(v_t=v(u_t)\) has derivative \((\dot u,\dot v)\).  Strict positivity persists for small \(|t|\).  If \(\dot p\) is its cell-probability derivative, its score is \(s=\dot p/p\), and
\[
D\Xi(P)[\dot p]_j=\sum_o\dot p(o)N_j(o)=E_P[s(O)N_j(O)]
\]
for the coordinate gradient \(N_j\).  Hence the score tangent is exactly the common kernel of these inner products, which is the orthogonal complement of their span.
\end{proof}
\par\smallskip\noindent \textit{Justification.} The lemma turns a checked finite constraint Jacobian into actual positive-probability paths rather than only a formal linearized kernel. \textit{Gap.} A finite Gram computation alone does not establish realization by paths in the model. \textit{Consumer.} It supplies the unconditional tangent calculations for the finite witness and comparator germs.

\begin{theorem}[thm:outcome-only-fixed-rank-comparison]\label{thm:outcome-only-fixed-rank-comparison}
For the single-valued outcome-only fixed-rank comparator in \(\mathrm{def:chen-xie-model}\),
\[
\mathcal T_{\mathrm{out},P^\star}=\mathcal C_{\mathrm{out},P^\star}^{\perp},
\qquad
\mathcal C_{\mathrm{out},P^\star}
=\operatorname{span}\{F_0,R_0,F_1,R_1\}.
\]
Its canonical gradient is \(\phi_{\mathrm{out}}\), and
\[
V_{\mathrm{out}}(P^\star)=\frac{30544}{13125}.
\]
Moreover, after shrinking the joint and outcome-only germs to one common neighborhood of \(P^\star\), \(\mathcal M_{\mathrm{joint}}(P^\star)\subset\mathcal M_{\mathrm{out}}(P^\star)\) and \(\theta_{\mathrm{out}}=\theta\) on the common joint germ, so
\[
\phi_{\mathrm{eff}}=\Pi_{\mathcal T_{P^\star}}\phi_{\mathrm{out}},
\qquad
V_{\mathrm{out}}(P^\star)-V_{\mathrm{eff}}(P^\star)
=\frac{2166784}{200416125}>0.
\]
The extra joint-model normals are precisely the nonredundant components of \(\operatorname{span}(G_0,G_1)\) modulo \(\mathcal C_{\mathrm{out},P^\star}\). This theorem concerns the explicit selected fixed-rank parameter, not an unidentified neighboring-law bridge functional and not a claim about the canonical gradient in a different published model.
\end{theorem}
\begin{proof}
On the positive finite outcome-only germ, the rank-two Schur equation contributes the \(R_a\) differential and the equation \(b_a\in\operatorname{range}(B_a)\) contributes the \(F_a\) differential.  With the ordering \(D_{\mathrm{out}}=(F_0,R_0,F_1,R_1)^\top\), exact summation under \(P^\star\) gives
\[
\Gamma_{\mathrm{out}}=E[D_{\mathrm{out}}D_{\mathrm{out}}^\top]
=\begin{pmatrix}
1323/500&63/20&0&0\\
63/20&441/2&0&0\\
0&0&1323/500&-63/20\\
0&0&-63/20&441/2
\end{pmatrix},
\]
\[
c_{\mathrm{out}}=E[D_{\mathrm{out}}\phi_{\mathrm{sel}}]
=(24/125,6/5,-24/125,6/5)^\top,
\quad
\Gamma_{\mathrm{out}}^{-1}c_{\mathrm{out}}
=(8/119,8/1785,-8/119,8/1785)^\top.
\]
The matrix is positive definite.  The finite constraint-tangent lemma therefore gives
\(\mathcal T_{\mathrm{out},P^\star}=\mathcal C_{\mathrm{out},P^\star}^\perp\).

To verify that \(\phi_{\mathrm{sel}}\) is a gradient of the selected parameter, differentiate \(B_{a,t}h^{\mathrm{out}}_{a,t}=b_{a,t}\) and
\(\theta_{\mathrm{out},a}(P_t)=d_t^\top h^{\mathrm{out}}_{a,t}\).  At \(t=0\), \(B_a^\top q_a^{\mathrm{sel}}=d\), so
\[
\dot\theta_{\mathrm{out},a}
=\dot d^{\top}h_a^{\mathrm{sel}}+d^\top\dot h_a
=\dot d^{\top}h_a^{\mathrm{sel}}
+q_a^{\mathrm{sel}\top}(\dot b_a-\dot B_a h_a^{\mathrm{sel}})
=E[S_a(h_a^{\mathrm{sel}},q_a^{\mathrm{sel}})s].
\]
After taking the signed arm contrast and centering, this is \(E[\phi_{\mathrm{sel}}s]\).  Orthogonal projection onto the outcome tangent consequently gives the displayed \(\phi_{\mathrm{out}}\).  Since \(E[\phi_{\mathrm{sel}}^2]=4432/1875\), direct substitution yields
\[
V_{\mathrm{out}}=\frac{4432}{1875}-c_{\mathrm{out}}^\top
\Gamma_{\mathrm{out}}^{-1}c_{\mathrm{out}}
=\frac{30544}{13125}.
\]

After shrinking both germs to one common neighborhood of \(P^\star\), every joint-model law is an outcome-only law.  On this common joint germ, treatment-bridge existence gives \(d_t^\top u_t=0\) for every outcome-bridge null shift \(u_t\), so the selected outcome bridge and every other outcome bridge have the same mean; hence \(\theta_{\mathrm{out}}(P_t)=\theta(P_t)\).  Thus \(\phi_{\mathrm{out}}\) is also a gradient on the smaller joint tangent.  Projection nesting gives \(\phi_{\mathrm{eff}}=\Pi_{\mathcal T_{P^\star}}\phi_{\mathrm{out}}\), and Pythagoras gives the variance difference.  Using the exact joint variance from \(\mathrm{thm:finite-witness-nonattainment}\),
\[
\frac{30544}{13125}-\frac{773723312}{334026875}
=\frac{2166784}{200416125}>0.
\]
Finally, the joint normal span is generated by \(F,G,R\), whereas the outcome-only rank-stratum span is generated by \(F,R\); quotienting the former by the latter leaves exactly the stated nonredundant \(G\) components.
\end{proof}
\par\smallskip\noindent \textit{Justification.} The explicit gauge selection makes the comparator parameter single-valued on every neighboring outcome-only law. \textit{Gap.} This is not represented as Chen--Xie's exact published canonical gradient without importing that model. \textit{Consumer.} It isolates the variance gain from treatment-solvability normals with a common target on the nested germ.

\begin{theorem}[thm:singular-branch-tangent-witness]\label{thm:singular-branch-tangent-witness}
At \(P^\circ\), both arm matrices are equal and rank two, both bridge fibers are non-singleton, and \(q_0=q_1=2\) is a treatment bridge. The six displayed generators \((F_0,G_0,R_0,F_1,G_1,R_1)\) have rank five because \(G_0+G_1=0\). Nevertheless the actual local joint-model tangent closure satisfies
\[
\mathcal T_{P^\circ}=\mathcal C_{P^\circ}^{\perp}.
\]
Thus a singular six-equation Jacobian does not by itself create an additional normal direction.
\end{theorem}
\begin{proof}
The arm matrices are
\[
B_0=B_1=\begin{pmatrix}
37/400&21/400&3/100\\
21/400&9/200&21/400\\
3/100&21/400&37/400
\end{pmatrix},
\]
which have rank two and left and right null vector \(v=(3,-7,3)^\top\). Their column sums are \(d/2\), so \(B_a^\top(2\mathbf 1)=d\).  Since \(I_0+I_1=1\),
\[
G_0(2,v)+G_1(2,v)
=v(W)\{2I_0-1+2I_1-1\}=0.
\]
Exact row reduction of the six generator columns gives rank five.

Let \(\mathcal R_a\) be the row plane of \(B_a\), and let \(d_a=B_a^\top\mathbf1\).  Treatment solvability is equivalent to \(d_0+d_1\in\mathcal R_0\cap\mathcal R_1\).  If \(\mathcal R_0\ne\mathcal R_1\), their intersection is a line, so this condition is equivalent to proportionality of \(d_0,d_1\); conversely proportionality places their sum in both row planes.  If \(\mathcal R_0=\mathcal R_1\), the condition holds automatically.  Hence the germ is the union of the common-row-space branch and the branch on which \(A\) and \(W\) are independent, together with the two outcome-solvability equations.

For the common-row-space branch, a normal basis consists of \(F_0,F_1\) and the four functions \(v(W)\ell(A,Z)\), where \(\ell\) ranges over the four-dimensional left kernel of the stacked \(6\)-by-\(3\) arm matrix.  These six columns have rank six.  For the independence branch, a normal basis consists of \(F_0,F_1,R_0,R_1\) and two centered \(A\)-by-\(W\) independence contrasts; these six columns also have rank six.  Substituting the displayed rational matrix and reducing the combined twelve columns gives rank seven.  Their normal-space intersection therefore has dimension \(6+6-7=5\).  Each of the five independent \(F,G,R\) generators belongs to both branch normal spaces, so that intersection equals \(\mathcal C_{P^\circ}\).  Positive finite paths realize both smooth branch tangents.  The tangent closure of their union is the span of the two branch tangents, whose orthocomplement is the intersection of their normal spaces, proving the claim.
\end{proof}
\par\smallskip\noindent \textit{Justification.} This reproduces the singular branch test where a single full-rank constraint-map argument is unavailable. \textit{Gap.} Constant arm ranks do not imply that a chosen joint list of solvability equations has independent differentials. \textit{Consumer.} It validates the proposed normal span at the first branch-crossing stress test.

\begin{proposition}[prop:four-by-four-rank-two-tangent-check]\label{prop:four-by-four-rank-two-tangent-check}
The law \(P^\square\) is strictly positive, both arm matrices have rank two, their row spaces are distinct with one-dimensional intersection, and both outcome and treatment bridges are nonunique. On the full positive joint fixed-rank germ through \(P^\square\), the actual tangent space obeys
\[
\mathcal T_{P^\square}=\mathcal C_{P^\square}^{\perp},
\qquad \dim\mathcal C_{P^\square}=15.
\]
In particular, the positive \(4\)-by-\(4\), rank-two, different-row-space early-kill test produces no normal direction outside the three declared classes.
\end{proposition}
\begin{proof}
The two matrices are
\[
B_0=\begin{pmatrix}
27/2000&3/125&3/80&1/20\\
23/2000&13/500&3/80&1/20\\
1/80&1/40&3/80&1/20\\
1/80&1/40&3/80&1/20
\end{pmatrix},
\]
\[
B_1=\begin{pmatrix}
1/80&1/40&3/80&1/20\\
1/80&13/500&73/2000&1/20\\
1/80&3/125&77/2000&1/20\\
1/80&1/40&3/80&1/20
\end{pmatrix}.
\]
Every entry is at least \(23/2000\), and the smallest observed \((Y,A,Z,W)\) cell has probability \(23/10000\).  The factorization in the definition shows
\[
\operatorname{row}(B_a)=\operatorname{span}\{d,r_a\},
\quad \operatorname{rank}(B_a)=2,
\quad B_a^\top(2\mathbf1)=d.
\]
Moreover \(d,r_0,r_1\) are linearly independent, so the two row spaces meet only in \(\operatorname{span}\{d\}\).  The displayed Bernoulli regressions give \(b_a=B_ah_a^\square\), establishing both bridge restrictions.  Each left and right kernel has dimension two.

Choose the following exact kernel bases:
\[
\ker B_0^\top=\operatorname{span}\{(-1/2,-1/2,1,0)^\top,(-1/2,-1/2,0,1)^\top\},
\]
\[
\ker B_1^\top=\operatorname{span}\{(-2,1,1,0)^\top,(-1,0,0,1)^\top\},
\]
\[
\ker B_0=\operatorname{span}\{(-1,-1,1,0)^\top,(-4/3,-4/3,0,1)^\top\},
\]
\[
\ker B_1=\operatorname{span}\{(-5,1,1,0)^\top,(-4,0,0,1)^\top\}.
\]
They generate four \(F\), four \(G\), and eight \(R\) columns.  Their only linear relation is the \(G\)-relation induced by the one-dimensional common right kernel; deleting either copy of that common-kernel column leaves a \(15\)-by-\(15\) Gram minor with determinant
\[
\frac{27}{524288000000}>0.
\]
Thus \(\dim\mathcal C_{P^\square}=15\).

It remains to compare this linear calculation with actual paths.  Near the base, every rank-two row space satisfying treatment solvability has a basis \([d_t,r_{a,t}]\), because the intersection of the two row planes remains the single line spanned by the positive marginal \(d_t\).  Write
\(B_{a,t}=L_{a,t}[d_t,r_{a,t}]^\top\).  The identity that \(d_t\) is the column sum of \(B_{0,t}+B_{1,t}\), together with linear independence of \(d_t,r_{0,t},r_{1,t}\), is exactly
\[
\mathbf1^\top L_{0,t}=(\pi_{0,t},0),\qquad
\mathbf1^\top L_{1,t}=(\pi_{1,t},0),qquad
\pi_{0,t}+\pi_{1,t}=1.
\]
After fixing the usual two basis gauges, these are smooth free coordinates for the full treatment-solvability incidence germ.  The rank equations contribute the eight \(R\) normals, the two outcome equations per arm contribute the four \(F\) normals, and the three independent displayed treatment equations contribute the three independent \(G\) normals.  Hence the smooth germ has codimension \(15\), and its normal space contains the already independent \(15\)-dimensional space \(\mathcal C_{P^\square}\); equality follows.  Every coordinate curve remains positive for small \(|t|\), hence belongs to the repaired saturated likelihood-ratio chart.  Actual path scores therefore span \(\mathcal C_{P^\square}^\perp\).
\end{proof}
\par\smallskip\noindent \textit{Justification.} This is the prescribed positive rank-two different-row-space early-kill test. \textit{Gap.} The shared marginal creates one cross-arm dependence among treatment normals, which must agree with the incidence-germ codimension. \textit{Consumer.} The exact rank certificate detects any missing cross-arm normal before a general lifting proof is attempted.

\begin{lemma}[lem:conditional-calibration-concentration]\label{lem:conditional-calibration-concentration}
Under \(\mathrm{ass:empirical-moment-perturbation}\), conditionally on every construction sigma-field \(\mathcal D_{A,k}\),
\[
 \|\widehat G_k^C-E[\widehat D_{J_n}^{(-k,A)}\widehat D_{J_n}^{(-k,A)\top}\mid\mathcal D_{A,k}]\|_{\mathrm{op}}=O_p(\eta_{G,n}),
\]
and
\[
 \|\widehat c_k^C-E[\widehat D_{J_n}^{(-k,A)}\widehat B^{(-k,A)}\mid\mathcal D_{A,k}]\|_2=O_p(\eta_{c,n}),
\]
uniformly over the fixed number of folds.
\end{lemma}
\begin{proof}
Fix a fold \(k\) and condition on \(\mathcal D_{A,k}\).  By \(\mathrm{ass:empirical-moment-perturbation}\), the learned dictionary and base signal are then fixed functions, the observations indexed by \(\mathcal C_k\) are conditionally independent, and the two stated moment-generating-function bounds hold in every deterministic Euclidean direction.  Write \(J=J_n\), \(N=n_{C,k}\), and
\[
 M_k=\widehat G_k^C-E[\widehat D_{J_n}^{(-k,A)}
 \widehat D_{J_n}^{(-k,A)\top}\mid\mathcal D_{A,k}].
\]
For a deterministic unit vector \(u\in\mathbb R^J\), the definition of \(\widehat G_k^C\) in \(\mathrm{def:affine-ensemble-estimator}\) gives
\[
 u^\top M_ku=N^{-1}\sum_{i\in\mathcal C_k}X_{i,u}.
\]
For \(0<\lambda\le c/B_{D,n}^2\), conditional independence and the first moment-generating-function inequality in \(\mathrm{ass:empirical-moment-perturbation}\) imply
\[
 E\!\left[\left.\exp\!\left\{\lambda\sum_{i\in\mathcal C_k}X_{i,u}\right\}\right|\mathcal D_{A,k}\right]
 \le \exp(CN\lambda^2B_{D,n}^4).
\]
Conditional exponential Markov inequality, optimized over the permitted \(\lambda\), and the same argument applied to \(-X_{i,u}\), therefore yield constants \(c_1>0\), independent of \(n,k,u\), such that for every \(x>0\),
\[
 \Pr\!\left(\left.|u^\top M_ku|>x\ \right|\mathcal D_{A,k}\right)
 \le 2\exp\!\left[-c_1N\min\!\left\{
 \frac{x^2}{B_{D,n}^4},\frac{x}{B_{D,n}^2}\right\}\right]. \tag{1}
\]

Let \(\mathcal N_J\) be a maximal \(1/4\)-separated subset of the Euclidean unit sphere.  The balls of radius \(1/8\) centered at its elements are disjoint and contained in the ball of radius \(9/8\), so volume comparison gives \(|\mathcal N_J|\le 9^J\).  Maximality makes \(\mathcal N_J\) a \(1/4\)-net.  If \(M\) is symmetric, choose a unit vector \(v\) with \(|v^\top Mv|=\|M\|_{\mathrm{op}}\), and then choose \(u\in\mathcal N_J\) with \(\|u-v\|_2\le1/4\).  Since
\[
 |v^\top Mv-u^\top Mu|
 \le |(v-u)^\top Mv|+|u^\top M(v-u)|
 \le \tfrac12\|M\|_{\mathrm{op}},
\]
we obtain
\[
 \|M\|_{\mathrm{op}}\le2\max_{u\in\mathcal N_J}|u^\top Mu|. \tag{2}
\]
Set \(q_J=J\).  For a constant \(C_1>0\), put
\[
 x_J=C_1B_{D,n}^2\left\{\sqrt{\frac{J}{N}}+\frac{J}{N}\right\}.
\]
For \(r=J/N\), both \((\sqrt r+r)^2\) and \(\sqrt r+r\) are at least \(r\).  Consequently,
\[
 N\min\!\left\{\frac{x_J^2}{B_{D,n}^4},
 \frac{x_J}{B_{D,n}^2}\right\}
 \ge J\min\{C_1^2,C_1\}.
\]
Equation \((1)\), the preceding lower bound, and a union bound over \(\mathcal N_J\), together with \((2)\), give
\[
 \Pr\!\left(\left.\|M_k\|_{\mathrm{op}}>2x_J\ \right|\mathcal D_{A,k}\right)
 \le 2\exp\!\left[J\log 9-c_1J\min\{C_1^2,C_1\}\right].
\]
Choosing \(C_1\) sufficiently large, and increasing it further for any prescribed tail probability, proves
\[
 \|M_k\|_{\mathrm{op}}
 =O_p\!\left(B_{D,n}^2\left\{\sqrt{\frac{J_n}{n_{C,k}}}
 +\frac{J_n}{n_{C,k}}\right\}\right)
 =O_p(\eta_{G,n}), \tag{3}
\]
where the final equality uses \(n_{C,k}\ge n_{C,\min}\) and the corrected definition in \(\mathrm{def:inference-gauges}\).

For the cross-moment, define the conditionally centered vector
\[
 V_i=\widehat D_{J_n}^{(-k,A)}(O_i)\widehat B^{(-k,A)}(O_i)
 -E[\widehat D_{J_n}^{(-k,A)}\widehat B^{(-k,A)}\mid\mathcal D_{A,k}],
\]
so \(u^\top V_i=Z_{i,u}\).  Apply the second moment-generating-function bound in \(\mathrm{ass:empirical-moment-perturbation}\) to \(Z_{i,u}\) and \(-Z_{i,u}\).  Since \(e^z+e^{-z}-2\ge z^2\), division by \(\lambda^2\) and passage to the limit \(\lambda\to0\) give a constant \(C_2<\infty\) such that
\[
 E[(u^\top V_i)^2\mid\mathcal D_{A,k}]
 \le C_2B_{D,n}^2B_{0,n}^2. \tag{4}
\]
Taking in \((4)\) each coordinate vector, and using conditional centering and conditional independence to eliminate cross-observation terms, yields
\[
 \begin{aligned}
 E\!\left[\left.\left\|N^{-1}\sum_{i\in\mathcal C_k}V_i\right\|_2^2
 \right|\mathcal D_{A,k}\right]
 &=N^{-2}\sum_{i\in\mathcal C_k}\sum_{c=1}^J
 E[V_{i,c}^2\mid\mathcal D_{A,k}]\\
 &\le \frac{C_2JB_{D,n}^2B_{0,n}^2}{N}.
 \end{aligned} \tag{5}
\]
Conditional Markov inequality applied to \((5)\) proves
\[
 \|\widehat c_k^C-E[\widehat D_{J_n}^{(-k,A)}\widehat B^{(-k,A)}
 \mid\mathcal D_{A,k}]\|_2
 =O_p\!\left(B_{D,n}B_{0,n}\sqrt{\frac{J_n}{n_{C,k}}}\right)
 =O_p(\eta_{c,n}).
\]
Because \(K_{\mathrm{fold}}\) is fixed, a finite union bound makes both asserted bounds uniform over \(k\); independence across folds is neither asserted nor needed.  Finally, if \(B_{D,n}\) and \(B_{0,n}\) are bounded, the displayed calibration gauges vanish under \(J_n/n_{C,\min}\to0\).  Thus the consistency condition supplied by these bounds is \(J_n/n_{C,\min}\to0\), with no extraneous \(\log(1+J_n)\) factor.  This proves the lemma.
\end{proof}
\par\smallskip\noindent \textit{Justification.} Conditional on the construction sample, a sphere-net Bernstein argument with \(|\mathcal N_{J_n}|\le9^{J_n}\) controls the learned Gram operator at the log-free dimension rate, while a direct conditional second-moment calculation controls the score--dictionary vector; a finite fold union bound supplies the asserted uniformity. \textit{Gap.} No proof gap remains under the declared conditional moment-generating-function envelope and fixed-fold condition. With bounded envelopes, the resulting calibration gauges vanish when \(J_n/n_{C,\min}\to0\), without a \(\log(1+J_n)\) loss. \textit{Consumer.} The sharpened bounds propagate through \(\zeta_{G,n}\), \(\zeta_{c,n}\), \(\eta_{\Pi,n}\), and \(\Gamma_n\), and feed the singular-Gram projector and resolvent perturbation in \(\mathrm{lem:learned-projection-convergence}\).

\begin{lemma}[lem:lindeberg-feller-central-limit]\label{lem:lindeberg-feller-central-limit}
Let \(X_{n1},\ldots,X_{nk_n}\) be independent mean-zero real random variables.  If \(\sum_iE[X_{ni}^2]\to V\in(0,\infty)\) and, for every \(\epsilon>0\), \(\sum_iE[X_{ni}^2\mathbf 1\{|X_{ni}|>\epsilon\}]\to0\), then \(\sum_iX_{ni}\Rightarrow N(0,V)\).
\end{lemma}

\section{Interpretation}

With redundant proxies, a chosen bridge pair can define a valid score without minimizing variance over the affine score hull. The exact finite witness proves this phenomenon and supplies a computable variance-minimizing replacement on that germ. For the SUPPORT application, the nested affine ensemble supplies certified standard errors and Wald intervals when the quotient hypotheses and the displayed learning and local-uniform conditions are verified. It does not assert that the published proxy specification has nonunique bridges, and global efficiency on the unrestricted joint model remains conditional on the unresolved reverse score-lifting equality.

\section*{Technical internal limitation (diagnostic only)}

The central unresolved mathematical step is reverse score lifting on the unrestricted saturated infinite-dimensional intersection. Operator-norm continuity gives continuous kernel projections but not a local parameterization of the coupled bridge-solvability loci or actual path lifts for every element of \(\mathcal C_P^\perp\). The proved estimator theorem additionally requires a dense population dictionary, fourth-moment bridge and multiplier control, conditional exponential calibration tails, the stated lower bound \(\kappa_{J_n}\) on positive population Gram eigenvalues, a cutoff satisfying \(0<\tau_n<\kappa_{J_n}/2\) and \(\zeta_{G,n}=o_p(\tau_n)\), and the stated ridge sequence, bounded coefficient and affine-weight envelopes, a vanishing learned-projection gauge, and a root-\(n\) mixed-bias product. Its regularity extension is only over path classes on which these conditions, quotient validity at every local law, likelihood-ratio transport continuity, uniform integrability, LAN, and contiguity hold uniformly. Closed range excludes generic compact infinite-rank inverse problems; rank-changing unions, general direct-integral measurable charts, and a SUPPORT reanalysis remain outside the proved scope.

\section{Honest scope}

Unconditional results are the normal--affine algebra, operator perturbation, the positive finite regular witness, the singular branch witness, the transverse four-level witness, strict single-pair nonattainment, and the explicit outcome-only selected-parameter comparison. Conditional on the hypotheses of \(\mathrm{thm:quotient-efficiency}\), the nested affine ensemble now has proved learned-score convergence, pointwise asymptotic linearity, ratio-consistent variance estimation, and Wald coverage; it is regular and locally efficient on the explicitly controlled contiguous path classes where the same quotient and primitive conditions hold uniformly. The reverse score-lifting OEQ remains open, so no unconditional global tangent-equality claim is made. The observed bridge functional receives a causal ATE interpretation only from the separate latent identification proposition, and no efficiency claim is made for path-poor subcharts, rank-changing unions, generic compact ill-posed operators, or an external outcome-only model whose neighboring-law parameter has not been imported.

\begin{thebibliography}{99}

  \bibitem{BennettEtAl2026StronglyIdentified} Bennett, A., Kallus, N., Mao, X., Newey, W., Syrgkanis, V., and Uehara, M. (2026). Inference on strongly identified functionals of weakly identified functions. Journal of the Royal Statistical Society: Series B, 88(3), 998--1032. doi:10.1093/jrsssb/qkaf075.

  \bibitem{ZhangLiMiaoTchetgen2023Nonuniqueness} Zhang, J., Li, W., Miao, W., and Tchetgen Tchetgen, E. (2023). Proximal causal inference without uniqueness assumptions. Statistics \& Probability Letters, 198, 109836. doi:10.1016/j.spl.2023.109836.

  \bibitem{ChenXie2026LocalOveridentification} Chen, X. and Xie, H. (2026). On local overidentification and efficiency gains in modern causal inference and data combination. Cowles Foundation Discussion Paper 2497, arXiv:2510.16683.

  \bibitem{CuiEtAl2024SemiparametricProximal} Cui, Y., Pu, H., Shi, X., Miao, W., and Tchetgen Tchetgen, E. (2024). Semiparametric proximal causal inference. Journal of the American Statistical Association. doi:10.1080/01621459.2023.2191817.

  \bibitem{ImbensEtAl2025LongTerm} Imbens, G., Kallus, N., Mao, X., and Wang, Y. (2025). Long-term causal inference under persistent confounding via data combination. Journal of the Royal Statistical Society: Series B, 87(2), 362--388. doi:10.1093/jrsssb/qkae095.

  \bibitem{AiShan2025EfficientProximalATE} Ai, C. and Shan, J. (2025). Efficient estimation of average treatment effects with unmeasured confounding and proxies. Statistica Sinica, arXiv:2501.02214.

  \bibitem{SeveriniTripathi2012EfficiencyBounds} Severini, T. A. and Tripathi, G. (2012). Efficiency bounds for estimating linear functionals of nonparametric regression models with endogenous regressors. Journal of Econometrics, 170(2), 491--498. doi:10.1016/j.jeconom.2012.05.018.

  \bibitem{ChenSantos2018Overidentification} Chen, X. and Santos, A. (2018). Overidentification in regular models. Econometrica, 86(5), 1771--1817. doi:10.3982/ECTA13559.

  \bibitem{AiChen2012SequentialMoments} Ai, C. and Chen, X. (2012). The semiparametric efficiency bound for models of sequential moment restrictions containing unknown functions. Journal of Econometrics, 170(2), 442--457. doi:10.1016/j.jeconom.2012.05.015.

  \bibitem{MiaoGengTchetgen2018Proxy} Miao, W., Geng, Z., and Tchetgen Tchetgen, E. (2018). Identifying causal effects with proxy variables of an unmeasured confounder. Biometrika, 105(4), 987--993. doi:10.1093/biomet/asy038.

  \bibitem{BennettEtAl2023SourceConditionDR} Bennett, A., Kallus, N., Mao, X., Newey, W., Syrgkanis, V., and Uehara, M. (2023). Source condition double robust inference on functionals of inverse problems. arXiv:2307.13793.

  \bibitem{vanDerVaart1998Asymptotic} van der Vaart, A. W. (1998). Asymptotic Statistics. Cambridge University Press.

\end{thebibliography}

\end{document}